% This LaTeX file was created by Metamath on 21-Dec-2023 3:10 PM.

\begin{metadefinition}\blTitle{df-bl.bool} Boolean type is true or false
\begin{align}
\vdash ( x \in \,\mathrm{Bool} \leftrightarrow x \in \{ \mathsf{true} ,
    \mathsf{false} \} )\label{eq:df-bl.bool}\tag{df-bl.bool}
\end{align}
\end{metadefinition}

\begin{lemma}\blTitle{bl.anfal} Conjunction with False on the Left
\begin{align}
\vdash ( ( \bot \wedge \mw{\varphi} ) \leftrightarrow \bot
    )\label{eq:bl.anfal}\tag{bl.anfal}
\end{align}
\end{lemma}

\begin{proof}
\begin{align}
  1 &&  & \vdash \lnot \bot&(\mbox{\ref{eq:fal}})\notag
  \\ 2 &&  & \vdash \lnot ( \bot \wedge \mw{\varphi} )&(\mbox{\ref{eq:intnanr}},1)\notag
  \\ 3 &&  & \vdash ( ( \bot \wedge \mw{\varphi} ) \leftrightarrow \bot )&(\mbox{\ref{eq:bifal}},2)\notag
\end{align}
\end{proof}

\begin{lemma}\blTitle{bl.anfar} Conjunction with False on the Right
\begin{align}
\vdash ( ( \mw{\varphi} \wedge \bot ) \leftrightarrow \bot
    )\label{eq:bl.anfar}\tag{bl.anfar}
\end{align}
\end{lemma}

\begin{proof}
\begin{align}
  1 &&  & \vdash \lnot \bot&(\mbox{\ref{eq:fal}})\notag
  \\ 2 &&  & \vdash \lnot ( \mw{\varphi} \wedge \bot )&(\mbox{\ref{eq:intnan}},1)\notag
  \\ 3 &&  & \vdash ( ( \mw{\varphi} \wedge \bot ) \leftrightarrow \bot )&(\mbox{\ref{eq:bifal}},2)\notag
\end{align}
\end{proof}

\begin{lemma}\blTitle{bl.ortrl} Disjunction with True on the Left
\begin{align}
\vdash ( ( \top \vee \mw{\varphi} ) \leftrightarrow \top
    )\label{eq:bl.ortrl}\tag{bl.ortrl}
\end{align}
\end{lemma}

\begin{proof}
\begin{align}
  1 &&  & \vdash \top&(\mbox{\ref{eq:tru}})\notag
  \\ 2 &&  & \vdash ( \top \vee \mw{\varphi} )&(\mbox{\ref{eq:orci}},1)\notag
  \\ 3 &&  & \vdash ( ( \top \vee \mw{\varphi} ) \leftrightarrow \top )&(\mbox{\ref{eq:bitru}},2)\notag
\end{align}
\end{proof}

\begin{lemma}\blTitle{bl.ortrr} Disjunction with True on the Right
\begin{align}
\vdash ( ( \mw{\varphi} \vee \top ) \leftrightarrow \top
    )\label{eq:bl.ortrr}\tag{bl.ortrr}
\end{align}
\end{lemma}

\begin{proof}
\begin{align}
  1 &&  & \vdash \top&(\mbox{\ref{eq:tru}})\notag
  \\ 2 &&  & \vdash ( \mw{\varphi} \vee \top )&(\mbox{\ref{eq:olci}},1)\notag
  \\ 3 &&  & \vdash ( ( \mw{\varphi} \vee \top ) \leftrightarrow \top )&(\mbox{\ref{eq:bitru}},2)\notag
\end{align}
\end{proof}

\begin{lemma}\blTitle{bl.orfar} Disjunction with False on the Right
\begin{align}
\vdash ( ( \mw{\varphi} \vee \bot ) \leftrightarrow \mw{\varphi}
    )\label{eq:bl.orfar}\tag{bl.orfar}
\end{align}
\end{lemma}

\begin{proof}
\begin{align}
  1 &&  & \vdash ( ( \mw{\varphi} \vee \bot ) \leftrightarrow \mw{\varphi} )&(\mbox{\ref{eq:orfa}})\notag
\end{align}
\end{proof}

\begin{lemma}\blTitle{bl.seq} Superfluity of Equivalence
\begin{align}
\vdash ( ( \mw{\varphi} \leftrightarrow \top ) \rightarrow \mw{\varphi}
    )\label{eq:bl.seq}\tag{bl.seq}
\end{align}
\end{lemma}

\begin{proof}
\begin{align}
  1 &&  & \vdash ( ( \mw{\varphi} \leftrightarrow \top ) \rightarrow ( \top \rightarrow \mw{\varphi} ) )&(\mbox{\ref{eq:biimpr}})\notag
  \\ 2 &&  & \vdash ( \mw{\varphi} \leftrightarrow ( \top \rightarrow \mw{\varphi} ) )&(\mbox{\ref{eq:trut}})\notag
  \\ 3 &&  & \vdash ( ( \top \rightarrow \mw{\varphi} ) \rightarrow \mw{\varphi} )&(\mbox{\ref{biimpri}},2)\notag
  \\ 4 &&  & \vdash ( ( \mw{\varphi} \leftrightarrow \top ) \rightarrow \mw{\varphi} )&(\mbox{\ref{syl}},1,3)\notag
\end{align}
\end{proof}

\begin{lemma}\blTitle{bl.an2impor2} P and Q implies P or R.
\begin{align}
\vdash ( ( \mw{\varphi} \wedge \mw{\psi} ) \rightarrow ( \mw{\varphi} \vee
    \mw{\chi} ) )\label{eq:bl.an2impor2}\tag{bl.an2impor2}
\end{align}
\end{lemma}

\begin{proof}
\begin{align}
  1 &&  & \vdash ( ( \mw{\varphi} \wedge \mw{\psi} ) \rightarrow \mw{\varphi} )&(\mbox{\ref{simpl}})\notag
  \\ 2 &&  & \vdash ( \mw{\varphi} \rightarrow ( \mw{\varphi} \vee \mw{\chi} ) )&(\mbox{\ref{orc}})\notag
  \\ 3 &&  & \vdash ( ( \mw{\varphi} \wedge \mw{\psi} ) \rightarrow ( \mw{\varphi} \vee \mw{\chi} ) )&(\mbox{\ref{syl}},1,2)\notag
\end{align}
\end{proof}

\begin{lemma}\blTitle{bl.PandQorPisP} Law of And-Simplification: P and (Q
or P) is P
\begin{align}
\vdash ( ( \mw{\varphi} \wedge ( \mw{\psi} \vee \mw{\varphi} ) )
    \leftrightarrow \mw{\varphi} )\label{eq:bl.PandQorPisP}\tag{bl.PandQorPisP}
\end{align}
\end{lemma}

\begin{proof}
\begin{align}
  1 &&  & \vdash ( \mw{\varphi} \leftrightarrow ( \mw{\varphi} \wedge ( \mw{\varphi} \vee \mw{\psi} ) ) )&(\mbox{\ref{pm4.45}})\notag
  \\ 2 &&  & \vdash ( ( \mw{\varphi} \leftrightarrow ( \mw{\varphi} \wedge ( \mw{\varphi} \vee \mw{\psi} ) ) ) \leftrightarrow ( ( \mw{\varphi} \wedge ( \mw{\varphi} \vee \mw{\psi} ) ) \leftrightarrow \mw{\varphi} ) )&(\mbox{\ref{eq:bicom}})\notag
  \\ 3 &&  & \vdash ( ( \mw{\varphi} \vee \mw{\psi} ) \leftrightarrow ( \mw{\psi} \vee \mw{\varphi} ) )&(\mbox{\ref{orcom}})\notag
  \\ 4 &&  & \vdash ( ( \mw{\varphi} \wedge ( \mw{\varphi} \vee \mw{\psi} ) ) \leftrightarrow ( \mw{\varphi} \wedge ( \mw{\psi} \vee \mw{\varphi} ) ) )&(\mbox{\ref{anbi2i}},3)\notag
  \\ 5 &&  & \vdash ( ( ( \mw{\varphi} \wedge ( \mw{\varphi} \vee \mw{\psi} ) ) \leftrightarrow \mw{\varphi} ) \leftrightarrow ( ( \mw{\varphi} \wedge ( \mw{\psi} \vee \mw{\varphi} ) ) \leftrightarrow \mw{\varphi} ) )&(\mbox{\ref{eq:bibi1i}},4)\notag
  \\ 6 &&  & \vdash ( ( \mw{\varphi} \leftrightarrow ( \mw{\varphi} \wedge ( \mw{\varphi} \vee \mw{\psi} ) ) ) \leftrightarrow ( ( \mw{\varphi} \wedge ( \mw{\psi} \vee \mw{\varphi} ) ) \leftrightarrow \mw{\varphi} ) )&(\mbox{\ref{bitri}},2,5)\notag
  \\ 7 &&  & \vdash ( ( \mw{\varphi} \wedge ( \mw{\psi} \vee \mw{\varphi} ) ) \leftrightarrow \mw{\varphi} )&(\mbox{\ref{eq:mpbi}},1,6)\notag
\end{align}
\end{proof}

\begin{lemma}\blTitle{bl.PandPorQisP} Law of And-Simplification: P and (P
or Q) is P
\begin{align}
\vdash ( ( \mw{\varphi} \wedge ( \mw{\varphi} \vee \mw{\psi} ) )
    \leftrightarrow \mw{\varphi} )\label{eq:bl.PandPorQisP}\tag{bl.PandPorQisP}
\end{align}
\end{lemma}

\begin{proof}
\begin{align}
  1 &&  & \vdash ( ( \mw{\varphi} \wedge ( \mw{\psi} \vee \mw{\varphi} ) ) \leftrightarrow \mw{\varphi} )&(\mbox{\ref{eq:bl.PandQorPisP}})\notag
  \\ 2 &&  & \vdash ( ( \mw{\psi} \vee \mw{\varphi} ) \leftrightarrow ( \mw{\varphi} \vee \mw{\psi} ) )&(\mbox{\ref{orcom}})\notag
  \\ 3 &&  & \vdash ( ( \mw{\varphi} \wedge ( \mw{\psi} \vee \mw{\varphi} ) ) \leftrightarrow ( \mw{\varphi} \wedge ( \mw{\varphi} \vee \mw{\psi} ) ) )&(\mbox{\ref{anbi2i}},2)\notag
  \\ 4 &&  & \vdash ( ( ( \mw{\varphi} \wedge ( \mw{\psi} \vee \mw{\varphi} ) ) \leftrightarrow \mw{\varphi} ) \leftrightarrow ( ( \mw{\varphi} \wedge ( \mw{\varphi} \vee \mw{\psi} ) ) \leftrightarrow \mw{\varphi} ) )&(\mbox{\ref{eq:bibi1i}},3)\notag
  \\ 5 &&  & \vdash ( ( \mw{\varphi} \wedge ( \mw{\varphi} \vee \mw{\psi} ) ) \leftrightarrow \mw{\varphi} )&(\mbox{\ref{eq:mpbi}},1,4)\notag
\end{align}
\end{proof}

\begin{lemma}\blTitle{bl.PisPandQorP} Law of And-Simplification: P is P
and (Q or P)
\begin{align}
\vdash ( \mw{\varphi} \leftrightarrow ( \mw{\varphi} \wedge ( \mw{\psi} \vee
    \mw{\varphi} ) ) )\label{eq:bl.PisPandQorP}\tag{bl.PisPandQorP}
\end{align}
\end{lemma}

\begin{proof}
\begin{align}
  1 &&  & \vdash ( \mw{\varphi} \leftrightarrow ( \mw{\varphi} \wedge ( \mw{\varphi} \vee \mw{\psi} ) ) )&(\mbox{\ref{pm4.45}})\notag
  \\ 2 &&  & \vdash ( ( \mw{\varphi} \vee \mw{\psi} ) \leftrightarrow ( \mw{\psi} \vee \mw{\varphi} ) )&(\mbox{\ref{orcom}})\notag
  \\ 3 &&  & \vdash ( ( \mw{\varphi} \wedge ( \mw{\varphi} \vee \mw{\psi} ) ) \leftrightarrow ( \mw{\varphi} \wedge ( \mw{\psi} \vee \mw{\varphi} ) ) )&(\mbox{\ref{anbi2i}},2)\notag
  \\ 4 &&  & \vdash ( ( \mw{\varphi} \leftrightarrow ( \mw{\varphi} \wedge ( \mw{\varphi} \vee \mw{\psi} ) ) ) \leftrightarrow ( \mw{\varphi} \leftrightarrow ( \mw{\varphi} \wedge ( \mw{\psi} \vee \mw{\varphi} ) ) ) )&(\mbox{\ref{eq:bibi2i}},3)\notag
  \\ 5 &&  & \vdash ( \mw{\varphi} \leftrightarrow ( \mw{\varphi} \wedge ( \mw{\psi} \vee \mw{\varphi} ) ) )&(\mbox{\ref{eq:mpbi}},1,4)\notag
\end{align}
\end{proof}

\begin{lemma}\blTitle{bl.PorQandPisP} Law of Or-Simplification: P or (Q
and P) is P
\begin{align}
\vdash ( ( \mw{\varphi} \vee ( \mw{\psi} \wedge \mw{\varphi} ) )
    \leftrightarrow \mw{\varphi} )\label{eq:bl.PorQandPisP}\tag{bl.PorQandPisP}
\end{align}
\end{lemma}

\begin{proof}
\begin{align}
  1 &&  & \vdash ( \mw{\varphi} \leftrightarrow ( \mw{\varphi} \vee ( \mw{\varphi} \wedge \mw{\psi} ) ) )&(\mbox{\mmref{pm4.44}})\notag
  \\ 2 &&  & \vdash ( ( \mw{\varphi} \leftrightarrow ( \mw{\varphi} \vee ( \mw{\varphi} \wedge \mw{\psi} ) ) ) \leftrightarrow ( ( \mw{\varphi} \vee ( \mw{\varphi} \wedge \mw{\psi} ) ) \leftrightarrow \mw{\varphi} ) )&(\mbox{\ref{eq:bicom}})\notag
  \\ 3 &&  & \vdash ( ( \mw{\varphi} \wedge \mw{\psi} ) \leftrightarrow ( \mw{\psi} \wedge \mw{\varphi} ) )&(\mbox{\ref{eq:ancom}})\notag
  \\ 4 &&  & \vdash ( ( \mw{\varphi} \vee ( \mw{\varphi} \wedge \mw{\psi} ) ) \leftrightarrow ( \mw{\varphi} \vee ( \mw{\psi} \wedge \mw{\varphi} ) ) )&(\mbox{\ref{eq:orbi2i}},3)\notag
  \\ 5 &&  & \vdash ( ( ( \mw{\varphi} \vee ( \mw{\varphi} \wedge \mw{\psi} ) ) \leftrightarrow \mw{\varphi} ) \leftrightarrow ( ( \mw{\varphi} \vee ( \mw{\psi} \wedge \mw{\varphi} ) ) \leftrightarrow \mw{\varphi} ) )&(\mbox{\ref{eq:bibi1i}},4)\notag
  \\ 6 &&  & \vdash ( ( \mw{\varphi} \leftrightarrow ( \mw{\varphi} \vee ( \mw{\varphi} \wedge \mw{\psi} ) ) ) \leftrightarrow ( ( \mw{\varphi} \vee ( \mw{\psi} \wedge \mw{\varphi} ) ) \leftrightarrow \mw{\varphi} ) )&(\mbox{\ref{eq:bitri}},2,5)\notag
  \\ 7 &&  & \vdash ( ( \mw{\varphi} \vee ( \mw{\psi} \wedge \mw{\varphi} ) ) \leftrightarrow \mw{\varphi} )&(\mbox{\ref{eq:mpbi}},1,6)\notag
\end{align}
\end{proof}

\begin{lemma}\blTitle{bl.PorPandQisP} Law of Or-Simplification: P or (P
and Q) is P
\begin{align}
\vdash ( ( \mw{\varphi} \vee ( \mw{\varphi} \wedge \mw{\psi} ) )
    \leftrightarrow \mw{\varphi} )\label{eq:bl.PorPandQisP}\tag{bl.PorPandQisP}
\end{align}
\end{lemma}

\begin{proof}
\begin{align}
  1 &&  & \vdash ( ( \mw{\varphi} \vee ( \mw{\psi} \wedge \mw{\varphi} ) ) \leftrightarrow \mw{\varphi} )&(\mbox{\ref{eq:bl.PorQandPisP}})\notag
  \\ 2 &&  & \vdash ( ( \mw{\psi} \wedge \mw{\varphi} ) \leftrightarrow ( \mw{\varphi} \wedge \mw{\psi} ) )&(\mbox{\ref{eq:ancom}})\notag
  \\ 3 &&  & \vdash ( ( \mw{\varphi} \vee ( \mw{\psi} \wedge \mw{\varphi} ) ) \leftrightarrow ( \mw{\varphi} \vee ( \mw{\varphi} \wedge \mw{\psi} ) ) )&(\mbox{\ref{eq:orbi2i}},2)\notag
  \\ 4 &&  & \vdash ( ( ( \mw{\varphi} \vee ( \mw{\psi} \wedge \mw{\varphi} ) ) \leftrightarrow \mw{\varphi} ) \leftrightarrow ( ( \mw{\varphi} \vee ( \mw{\varphi} \wedge \mw{\psi} ) ) \leftrightarrow \mw{\varphi} ) )&(\mbox{\ref{eq:bibi1i}},3)\notag
  \\ 5 &&  & \vdash ( ( \mw{\varphi} \vee ( \mw{\varphi} \wedge \mw{\psi} ) ) \leftrightarrow \mw{\varphi} )&(\mbox{\ref{eq:mpbi}},1,4)\notag
\end{align}
\end{proof}

\begin{lemma}\blTitle{bl.PisPorQandP} Law of Or-Simplification: P is P or
(Q and P)
\begin{align}
\vdash ( \mw{\varphi} \leftrightarrow ( \mw{\varphi} \vee ( \mw{\psi} \wedge
    \mw{\varphi} ) ) )\label{eq:bl.PisPorQandP}\tag{bl.PisPorQandP}
\end{align}
\end{lemma}

\begin{proof}
\begin{align}
  1 &&  & \vdash ( \mw{\varphi} \leftrightarrow ( \mw{\varphi} \vee ( \mw{\varphi} \wedge \mw{\psi} ) ) )&(\mbox{\ref{eq:pm4.44}})\notag
  \\ 2 &&  & \vdash ( ( \mw{\varphi} \wedge \mw{\psi} ) \leftrightarrow ( \mw{\psi} \wedge \mw{\varphi} ) )&(\mbox{\ref{eq:ancom}})\notag
  \\ 3 &&  & \vdash ( ( \mw{\varphi} \vee ( \mw{\varphi} \wedge \mw{\psi} ) ) \leftrightarrow ( \mw{\varphi} \vee ( \mw{\psi} \wedge \mw{\varphi} ) ) )&(\mbox{\ref{eq:orbi2i}},2)\notag
  \\ 4 &&  & \vdash ( ( \mw{\varphi} \leftrightarrow ( \mw{\varphi} \vee ( \mw{\varphi} \wedge \mw{\psi} ) ) ) \leftrightarrow ( \mw{\varphi} \leftrightarrow ( \mw{\varphi} \vee ( \mw{\psi} \wedge \mw{\varphi} ) ) ) )&(\mbox{\ref{eq:bibi2i}},3)\notag
  \\ 5 &&  & \vdash ( \mw{\varphi} \leftrightarrow ( \mw{\varphi} \vee ( \mw{\psi} \wedge \mw{\varphi} ) ) )&(\mbox{\ref{eq:mpbi}},1,4)\notag
\end{align}
\end{proof}

\begin{metadefinition}\blTitle{df-bl.lan1} Just one wff in the wff-list for
conjunction.
\begin{align}
\vdash ( \bigwedge( \mw{\varphi} ) \leftrightarrow \mw{\varphi}
    )\label{eq:df-bl.lan1}\tag{df-bl.lan1}
\end{align}
\end{metadefinition}

\begin{metadefinition}\blTitle{df-bl.lan2wl} Conjunction of two wff-lists.
\begin{align}
\vdash ( \bigwedge( \mwl{l_1} \mwl{l_2} ) \leftrightarrow ( \bigwedge(
    \mwl{l_1} ) \wedge \bigwedge( \mwl{l_2} ) )
    )\label{eq:df-bl.lan2wl}\tag{df-bl.lan2wl}
\end{align}
\end{metadefinition}

\begin{metadefinition}\blTitle{df-bl.lor1} Just one wff in the wff-list for
disjunction.
\begin{align}
\vdash ( \bigvee( \mw{\varphi} ) \leftrightarrow \mw{\varphi}
    )\label{eq:df-bl.lor1}\tag{df-bl.lor1}
\end{align}
\end{metadefinition}

\begin{metadefinition}\blTitle{df-bl.lor2wl} Disjunction of wff-lists.
\begin{align}
\vdash ( \bigvee( \mwl{l_1} \mwl{l_2} ) \leftrightarrow ( \bigvee( \mwl{l_1} )
    \vee \bigvee( \mwl{l_2} ) ) )\label{eq:df-bl.lor2wl}\tag{df-bl.lor2wl}
\end{align}
\end{metadefinition}

\begin{lemma}\blTitle{bl.an2wl} Conjunction of two wffs.
\begin{align}
\vdash ( \bigwedge( \mw{\varphi} \mw{\psi} ) \leftrightarrow ( \mw{\varphi}
    \wedge \mw{\psi} ) )\label{eq:bl.an2wl}\tag{bl.an2wl}
\end{align}
\end{lemma}

\begin{proof}
\begin{align}
  1 &&  & \vdash ( \bigwedge( \mw{\varphi} \mw{\psi} ) \leftrightarrow ( \bigwedge( \mw{\varphi} ) \wedge \bigwedge( \mw{\psi} ) ) )&(\mbox{\ref{eq:df-bl.lan2wl}})\notag
  \\ 2 &&  & \vdash ( \bigwedge( \mw{\varphi} ) \leftrightarrow \mw{\varphi} )&(\mbox{\ref{eq:df-bl.lan1}})\notag
  \\ 3 &&  & \vdash ( \bigwedge( \mw{\psi} ) \leftrightarrow \mw{\psi} )&(\mbox{\ref{eq:df-bl.lan1}})\notag
  \\ 4 &&  & \vdash ( ( \bigwedge( \mw{\varphi} ) \wedge \bigwedge( \mw{\psi} ) ) \leftrightarrow ( \mw{\varphi} \wedge \mw{\psi} ) )&(\mbox{\ref{eq:anbi12i}},2,3)\notag
  \\ 5 &&  & \vdash ( \bigwedge( \mw{\varphi} \mw{\psi} ) \leftrightarrow ( \mw{\varphi} \wedge \mw{\psi} ) )&(\mbox{\ref{eq:bitri}},1,4)\notag
\end{align}
\end{proof}

\begin{lemma}\blTitle{bl.an3wl} Conjunction of three wffs.
\begin{align}
\vdash ( \bigwedge( \mw{\varphi} \mw{\psi} \mw{\chi} ) \leftrightarrow (
    \mw{\varphi} \wedge \mw{\psi} \wedge \mw{\chi} )
    )\label{eq:bl.an3wl}\tag{bl.an3wl}
\end{align}
\end{lemma}

\begin{proof}
\begin{align}
  1 &&  & \vdash ( \bigwedge( \mw{\varphi} \mw{\psi} \mw{\chi} ) \leftrightarrow ( \bigwedge( \mw{\varphi} ) \wedge \bigwedge( \mw{\psi} \mw{\chi} ) ) )&(\mbox{\ref{eq:df-bl.lan2wl}})\notag
  \\ 2 &&  & \vdash ( \bigwedge( \mw{\varphi} ) \leftrightarrow \mw{\varphi} )&(\mbox{\ref{eq:df-bl.lan1}})\notag
  \\ 3 &&  & \vdash ( \bigwedge( \mw{\psi} \mw{\chi} ) \leftrightarrow ( \mw{\psi} \wedge \mw{\chi} ) )&(\mbox{\ref{eq:bl.an2wl}})\notag
  \\ 4 &&  & \vdash ( ( \bigwedge( \mw{\varphi} ) \wedge \bigwedge( \mw{\psi} \mw{\chi} ) ) \leftrightarrow ( \mw{\varphi} \wedge ( \mw{\psi} \wedge \mw{\chi} ) ) )&(\mbox{\ref{eq:anbi12i}},2,3)\notag
  \\ 5 &&  & \vdash ( \bigwedge( \mw{\varphi} \mw{\psi} \mw{\chi} ) \leftrightarrow ( \mw{\varphi} \wedge ( \mw{\psi} \wedge \mw{\chi} ) ) )&(\mbox{\ref{eq:bitri}},1,4)\notag
  \\ 6 &&  & \vdash ( ( \mw{\varphi} \wedge \mw{\psi} \wedge \mw{\chi} ) \leftrightarrow ( \mw{\varphi} \wedge ( \mw{\psi} \wedge \mw{\chi} ) ) )&(\mbox{\ref{eq:3anass}})\notag
  \\ 7 &&  & \vdash ( \bigwedge( \mw{\varphi} \mw{\psi} \mw{\chi} ) \leftrightarrow ( \mw{\varphi} \wedge \mw{\psi} \wedge \mw{\chi} ) )&(\mbox{\ref{eq:bitr4i}},5,6)\notag
\end{align}
\end{proof}

\begin{lemma}\blTitle{bl.ancomphwl} Move first wff to end of conjunction
wff-list.
\begin{align}
\vdash ( \bigwedge( \mw{\varphi} \mwl{l_1} ) \leftrightarrow \bigwedge(
    \mwl{l_1} \mw{\varphi} ) )\label{eq:bl.ancomphwl}\tag{bl.ancomphwl}
\end{align}
\end{lemma}

\begin{proof}
\begin{align}
  1 &&  & \vdash ( \bigwedge( \mwl{l_1} \mw{\varphi} ) \leftrightarrow ( \bigwedge( \mwl{l_1} ) \wedge \bigwedge( \mw{\varphi} ) ) )&(\mbox{\ref{eq:df-bl.lan2wl}})\notag
  \\ 2 && @2:\  & \vdash ( \bigwedge( \mw{\varphi} ) \leftrightarrow \mw{\varphi} )&(\mbox{\ref{eq:df-bl.lan1}})\notag
  \\ 3 &&  & \vdash ( ( \bigwedge( \mwl{l_1} ) \wedge \bigwedge( \mw{\varphi} ) ) \leftrightarrow ( \bigwedge( \mwl{l_1} ) \wedge \mw{\varphi} ) )&(\mbox{\ref{eq:anbi2i}},2)\notag
  \\ 4 &&  & \vdash ( \bigwedge( \mw{\varphi} ) \leftrightarrow \mw{\varphi} )&(@2)\notag
  \\ 5 &&  & \vdash ( ( \bigwedge( \mw{\varphi} ) \wedge \bigwedge( \mwl{l_1} ) ) \leftrightarrow ( \mw{\varphi} \wedge \bigwedge( \mwl{l_1} ) ) )&(\mbox{\ref{eq:anbi1i}},4)\notag
  \\ 6 &&  & \vdash ( \bigwedge( \mw{\varphi} \mwl{l_1} ) \leftrightarrow ( \bigwedge( \mw{\varphi} ) \wedge \bigwedge( \mwl{l_1} ) ) )&(\mbox{\ref{eq:df-bl.lan2wl}})\notag
  \\ 7 &&  & \vdash ( ( \bigwedge( \mwl{l_1} ) \wedge \mw{\varphi} ) \leftrightarrow ( \mw{\varphi} \wedge \bigwedge( \mwl{l_1} ) ) )&(\mbox{\ref{eq:ancom}})\notag
  \\ 8 &&  & \vdash ( ( \bigwedge( \mwl{l_1} ) \wedge \mw{\varphi} ) \leftrightarrow \bigwedge( \mw{\varphi} \mwl{l_1} ) )&(\mbox{\ref{eq:3bitr4ri}},5,6,7)\notag
  \\ 9 &&  & \vdash ( ( \bigwedge( \mwl{l_1} ) \wedge \bigwedge( \mw{\varphi} ) ) \leftrightarrow \bigwedge( \mw{\varphi} \mwl{l_1} ) )&(\mbox{\ref{eq:bitri}},3,8)\notag
  \\ 10 &&  & \vdash ( \bigwedge( \mw{\varphi} \mwl{l_1} ) \leftrightarrow \bigwedge( \mwl{l_1} \mw{\varphi} ) )&(\mbox{\ref{eq:bitr2i}},1,9)\notag
\end{align}
\end{proof}

\begin{lemma}\blTitle{bl.ancomwlwl} Commute conjunction of two wff-lists.
\begin{align}
\vdash ( \bigwedge( \mwl{l_1} \mwl{l_2} ) \leftrightarrow \bigwedge( \mwl{l_2}
    \mwl{l_1} ) )\label{eq:bl.ancomwlwl}\tag{bl.ancomwlwl}
\end{align}
\end{lemma}

\begin{proof}
\begin{align}
  1 &&  & \vdash ( ( \bigwedge( \mwl{l_1} ) \wedge \bigwedge( \mwl{l_2} ) ) \leftrightarrow ( \bigwedge( \mwl{l_2} ) \wedge \bigwedge( \mwl{l_1} ) ) )&(\mbox{\ref{eq:ancom}})\notag
  \\ 2 &&  & \vdash ( \bigwedge( \mwl{l_1} \mwl{l_2} ) \leftrightarrow ( \bigwedge( \mwl{l_1} ) \wedge \bigwedge( \mwl{l_2} ) ) )&(\mbox{\ref{eq:df-bl.lan2wl}})\notag
  \\ 3 &&  & \vdash ( \bigwedge( \mwl{l_2} \mwl{l_1} ) \leftrightarrow ( \bigwedge( \mwl{l_2} ) \wedge \bigwedge( \mwl{l_1} ) ) )&(\mbox{\ref{eq:df-bl.lan2wl}})\notag
  \\ 4 &&  & \vdash ( \bigwedge( \mwl{l_1} \mwl{l_2} ) \leftrightarrow \bigwedge( \mwl{l_2} \mwl{l_1} ) )&(\mbox{\ref{eq:3bitr4i}},1,2,3)\notag
\end{align}
\end{proof}

\begin{lemma}\blTitle{bl.anpfw} Pull first wff from conjunction wff-list.
\begin{align}
\vdash ( \bigwedge( \mw{\varphi} \mwl{l_1} ) \leftrightarrow ( \mw{\varphi}
    \wedge \bigwedge( \mwl{l_1} ) ) )\label{eq:bl.anpfw}\tag{bl.anpfw}
\end{align}
\end{lemma}

\begin{proof}
\begin{align}
  1 &&  & \vdash ( \bigwedge( \mw{\varphi} \mwl{l_1} ) \leftrightarrow ( \bigwedge( \mw{\varphi} ) \wedge \bigwedge( \mwl{l_1} ) ) )&(\mbox{\ref{eq:df-bl.lan2wl}})\notag
  \\ 2 &&  & \vdash ( \bigwedge( \mw{\varphi} ) \leftrightarrow \mw{\varphi} )&(\mbox{\ref{eq:df-bl.lan1}})\notag
  \\ 3 &&  & \vdash ( ( \bigwedge( \mw{\varphi} ) \wedge \bigwedge( \mwl{l_1} ) ) \leftrightarrow ( \mw{\varphi} \wedge \bigwedge( \mwl{l_1} ) ) )&(\mbox{\ref{eq:anbi1i}},2)\notag
  \\ 4 &&  & \vdash ( \bigwedge( \mw{\varphi} \mwl{l_1} ) \leftrightarrow ( \mw{\varphi} \wedge \bigwedge( \mwl{l_1} ) ) )&(\mbox{\ref{eq:bitri}},1,3)\notag
\end{align}
\end{proof}

\begin{lemma}\blTitle{bl.anplw} Pull last wff from conjunction wff-list.
\begin{align}
\vdash ( \bigwedge( \mwl{l_1} \mw{\varphi} ) \leftrightarrow ( \bigwedge(
    \mwl{l_1} ) \wedge \mw{\varphi} ) )\label{eq:bl.anplw}\tag{bl.anplw}
\end{align}
\end{lemma}

\begin{proof}
\begin{align}
  1 &&  & \vdash ( \bigwedge( \mwl{l_1} \mw{\varphi} ) \leftrightarrow ( \bigwedge( \mwl{l_1} ) \wedge \bigwedge( \mw{\varphi} ) ) )&(\mbox{\ref{eq:df-bl.lan2wl}})\notag
  \\ 2 &&  & \vdash ( \bigwedge( \mw{\varphi} ) \leftrightarrow \mw{\varphi} )&(\mbox{\ref{eq:df-bl.lan1}})\notag
  \\ 3 &&  & \vdash ( ( \bigwedge( \mwl{l_1} ) \wedge \bigwedge( \mw{\varphi} ) ) \leftrightarrow ( \bigwedge( \mwl{l_1} ) \wedge \mw{\varphi} ) )&(\mbox{\ref{eq:anbi2i}},2)\notag
  \\ 4 &&  & \vdash ( \bigwedge( \mwl{l_1} \mw{\varphi} ) \leftrightarrow ( \bigwedge( \mwl{l_1} ) \wedge \mw{\varphi} ) )&(\mbox{\ref{eq:bitri}},1,3)\notag
\end{align}
\end{proof}

\begin{lemma}\blTitle{bl.anabpf} Absorb first term conjunction wff-list.
\begin{align}
\vdash ( \bigwedge( \bigwedge( \mwl{l_1} ) \mwl{l_2} ) \leftrightarrow
    \bigwedge( \mwl{l_1} \mwl{l_2} ) )\label{eq:bl.anabpf}\tag{bl.anabpf}
\end{align}
\end{lemma}

\begin{proof}
\begin{align}
  1 &&  & \vdash ( \bigwedge( \bigwedge( \mwl{l_1} ) \mwl{l_2} ) \leftrightarrow ( \bigwedge( \bigwedge( \mwl{l_1} ) ) \wedge \bigwedge( \mwl{l_2} ) ) )&(\mbox{\ref{eq:df-bl.lan2wl}})\notag
  \\ 2 &&  & \vdash ( \bigwedge( \mwl{l_1} \mwl{l_2} ) \leftrightarrow ( \bigwedge( \mwl{l_1} ) \wedge \bigwedge( \mwl{l_2} ) ) )&(\mbox{\ref{eq:df-bl.lan2wl}})\notag
  \\ 3 &&  & \vdash ( \bigwedge( \bigwedge( \mwl{l_1} ) ) \leftrightarrow \bigwedge( \mwl{l_1} ) )&(\mbox{\ref{eq:df-bl.lan1}})\notag
  \\ 4 &&  & \vdash ( ( \bigwedge( \bigwedge( \mwl{l_1} ) ) \wedge \bigwedge( \mwl{l_2} ) ) \leftrightarrow ( \bigwedge( \mwl{l_1} ) \wedge \bigwedge( \mwl{l_2} ) ) )&(\mbox{\ref{eq:anbi1i}},3)\notag
  \\ 5 &&  & \vdash ( \bigwedge( \mwl{l_1} \mwl{l_2} ) \leftrightarrow ( \bigwedge( \bigwedge( \mwl{l_1} ) ) \wedge \bigwedge( \mwl{l_2} ) ) )&(\mbox{\ref{eq:bitr4i}},2,4)\notag
  \\ 6 &&  & \vdash ( \bigwedge( \bigwedge( \mwl{l_1} ) \mwl{l_2} ) \leftrightarrow \bigwedge( \mwl{l_1} \mwl{l_2} ) )&(\mbox{\ref{eq:bitr4i}},1,5)\notag
\end{align}
\end{proof}

\begin{lemma}\blTitle{bl.anabpl} Absorb last term conjunction wff-list.
\begin{align}
\vdash ( \bigwedge( \mwl{l_1} \bigwedge( \mwl{l_2} ) ) \leftrightarrow
    \bigwedge( \mwl{l_1} \mwl{l_2} ) )\label{eq:bl.anabpl}\tag{bl.anabpl}
\end{align}
\end{lemma}

\begin{proof}
\begin{align}
  1 &&  & \vdash ( \bigwedge( \bigwedge( \mwl{l_2} ) ) \leftrightarrow \bigwedge( \mwl{l_2} ) )&(\mbox{\ref{eq:df-bl.lan1}})\notag
  \\ 2 &&  & \vdash ( ( \bigwedge( \mwl{l_1} ) \wedge \bigwedge( \bigwedge( \mwl{l_2} ) ) ) \leftrightarrow ( \bigwedge( \mwl{l_1} ) \wedge \bigwedge( \mwl{l_2} ) ) )&(\mbox{\ref{eq:anbi2i}},1)\notag
  \\ 3 &&  & \vdash ( \bigwedge( \mwl{l_1} \bigwedge( \mwl{l_2} ) ) \leftrightarrow ( \bigwedge( \mwl{l_1} ) \wedge \bigwedge( \bigwedge( \mwl{l_2} ) ) ) )&(\mbox{\ref{eq:df-bl.lan2wl}})\notag
  \\ 4 &&  & \vdash ( \bigwedge( \mwl{l_1} \mwl{l_2} ) \leftrightarrow ( \bigwedge( \mwl{l_1} ) \wedge \bigwedge( \mwl{l_2} ) ) )&(\mbox{\ref{eq:df-bl.lan2wl}})\notag
  \\ 5 &&  & \vdash ( \bigwedge( \mwl{l_1} \bigwedge( \mwl{l_2} ) ) \leftrightarrow \bigwedge( \mwl{l_1} \mwl{l_2} ) )&(\mbox{\ref{eq:3bitr4i}},2,3,4)\notag
\end{align}
\end{proof}

\begin{lemma}\blTitle{bl.or2wl} Disjunction of two wffs
\begin{align}
\vdash ( \bigvee( \mw{\varphi} \mw{\psi} ) \leftrightarrow ( \mw{\varphi} \vee
    \mw{\psi} ) )\label{eq:bl.or2wl}\tag{bl.or2wl}
\end{align}
\end{lemma}

\begin{proof}
\begin{align}
  1 &&  & \vdash ( \bigvee( \mw{\varphi} \mw{\psi} ) \leftrightarrow ( \bigvee( \mw{\varphi} ) \vee \bigvee( \mw{\psi} ) ) )&(\mbox{\ref{eq:df-bl.lor2wl}})\notag
  \\ 2 &&  & \vdash ( \bigvee( \mw{\varphi} ) \leftrightarrow \mw{\varphi} )&(\mbox{\ref{eq:df-bl.lor1}})\notag
  \\ 3 &&  & \vdash ( \bigvee( \mw{\psi} ) \leftrightarrow \mw{\psi} )&(\mbox{\ref{eq:df-bl.lor1}})\notag
  \\ 4 &&  & \vdash ( ( \bigvee( \mw{\varphi} ) \vee \bigvee( \mw{\psi} ) ) \leftrightarrow ( \mw{\varphi} \vee \mw{\psi} ) )&(\mbox{\ref{eq:orbi12i}},2,3)\notag
  \\ 5 &&  & \vdash ( \bigvee( \mw{\varphi} \mw{\psi} ) \leftrightarrow ( \mw{\varphi} \vee \mw{\psi} ) )&(\mbox{\ref{eq:bitri}},1,4)\notag
\end{align}
\end{proof}

\begin{lemma}\blTitle{bl.or3wl} Disjunction of three wffs
\begin{align}
\vdash ( \bigvee( \mw{\varphi} \mw{\psi} \mw{\chi} ) \leftrightarrow (
    \mw{\varphi} \vee \mw{\psi} \vee \mw{\chi} )
    )\label{eq:bl.or3wl}\tag{bl.or3wl}
\end{align}
\end{lemma}

\begin{proof}
\begin{align}
  1 &&  & \vdash ( \bigvee( \mw{\varphi} \mw{\psi} \mw{\chi} ) \leftrightarrow ( \bigvee( \mw{\varphi} \mw{\psi} ) \vee \bigvee( \mw{\chi} ) ) )&(\mbox{\ref{eq:df-bl.lor2wl}})\notag
  \\ 2 &&  & \vdash ( ( \bigvee( \mw{\varphi} ) \vee \bigvee( \mw{\psi} ) \vee \bigvee( \mw{\chi} ) ) \leftrightarrow ( ( \bigvee( \mw{\varphi} ) \vee \bigvee( \mw{\psi} ) ) \vee \bigvee( \mw{\chi} ) ) )&(\mbox{\ref{eq:df-3or}})\notag
  \\ 3 &&  & \vdash ( \bigvee( \mw{\varphi} \mw{\psi} ) \leftrightarrow ( \bigvee( \mw{\varphi} ) \vee \bigvee( \mw{\psi} ) ) )&(\mbox{\ref{eq:df-bl.lor2wl}})\notag
  \\ 4 &&  & \vdash ( ( \bigvee( \mw{\varphi} \mw{\psi} ) \vee \bigvee( \mw{\chi} ) ) \leftrightarrow ( ( \bigvee( \mw{\varphi} ) \vee \bigvee( \mw{\psi} ) ) \vee \bigvee( \mw{\chi} ) ) )&(\mbox{\ref{eq:orbi1i}},3)\notag
  \\ 5 &&  & \vdash ( ( \bigvee( \mw{\varphi} ) \vee \bigvee( \mw{\psi} ) \vee \bigvee( \mw{\chi} ) ) \leftrightarrow ( \bigvee( \mw{\varphi} \mw{\psi} ) \vee \bigvee( \mw{\chi} ) ) )&(\mbox{\ref{eq:bitr4i}},2,4)\notag
  \\ 6 &&  & \vdash ( \bigvee( \mw{\varphi} ) \leftrightarrow \mw{\varphi} )&(\mbox{\ref{eq:df-bl.lor1}})\notag
  \\ 7 &&  & \vdash ( \bigvee( \mw{\psi} ) \leftrightarrow \mw{\psi} )&(\mbox{\ref{eq:df-bl.lor1}})\notag
  \\ 8 &&  & \vdash ( \bigvee( \mw{\chi} ) \leftrightarrow \mw{\chi} )&(\mbox{\ref{eq:df-bl.lor1}})\notag
  \\ 9 &&  & \vdash ( ( \bigvee( \mw{\varphi} ) \vee \bigvee( \mw{\psi} ) \vee \bigvee( \mw{\chi} ) ) \leftrightarrow ( \mw{\varphi} \vee \mw{\psi} \vee \mw{\chi} ) )&(\mbox{\ref{eq:3orbi123i}},6,7,8)\notag
  \\ 10 &&  & \vdash ( ( \bigvee( \mw{\varphi} \mw{\psi} ) \vee \bigvee( \mw{\chi} ) ) \leftrightarrow ( \mw{\varphi} \vee \mw{\psi} \vee \mw{\chi} ) )&(\mbox{\ref{eq:bitr3i}},5,9)\notag
  \\ 11 &&  & \vdash ( \bigvee( \mw{\varphi} \mw{\psi} \mw{\chi} ) \leftrightarrow ( \mw{\varphi} \vee \mw{\psi} \vee \mw{\chi} ) )&(\mbox{\ref{eq:bitri}},1,10)\notag
\end{align}
\end{proof}

\begin{lemma}\blTitle{bl.orcomphwl} Move first wff to end of disjunction
wff-list.
\begin{align}
\vdash ( \bigvee( \mw{\varphi} \mwl{l_1} ) \leftrightarrow \bigvee( \mwl{l_1}
    \mw{\varphi} ) )\label{eq:bl.orcomphwl}\tag{bl.orcomphwl}
\end{align}
\end{lemma}

\begin{proof}
\begin{align}
  1 &&  & \vdash ( ( \bigvee( \mw{\varphi} ) \vee \bigvee( \mwl{l_1} ) ) \leftrightarrow ( \bigvee( \mwl{l_1} ) \vee \bigvee( \mw{\varphi} ) ) )&(\mbox{\ref{eq:orcom}})\notag
  \\ 2 &&  & \vdash ( \bigvee( \mw{\varphi} \mwl{l_1} ) \leftrightarrow ( \bigvee( \mw{\varphi} ) \vee \bigvee( \mwl{l_1} ) ) )&(\mbox{\ref{eq:df-bl.lor2wl}})\notag
  \\ 3 &&  & \vdash ( \bigvee( \mwl{l_1} \mw{\varphi} ) \leftrightarrow ( \bigvee( \mwl{l_1} ) \vee \bigvee( \mw{\varphi} ) ) )&(\mbox{\ref{eq:df-bl.lor2wl}})\notag
  \\ 4 &&  & \vdash ( \bigvee( \mw{\varphi} \mwl{l_1} ) \leftrightarrow \bigvee( \mwl{l_1} \mw{\varphi} ) )&(\mbox{\ref{eq:3bitr4i}},1,2,3)\notag
\end{align}
\end{proof}

\begin{lemma}\blTitle{bl.orcomwlwl} Commute disjunction of two wff-lists.
\begin{align}
\vdash ( \bigvee( \mwl{l_1} \mwl{l_2} ) \leftrightarrow \bigvee( \mwl{l_2}
    \mwl{l_1} ) )\label{eq:bl.orcomwlwl}\tag{bl.orcomwlwl}
\end{align}
\end{lemma}

\begin{proof}
\begin{align}
  1 &&  & \vdash ( ( \bigvee( \mwl{l_1} ) \vee \bigvee( \mwl{l_2} ) ) \leftrightarrow ( \bigvee( \mwl{l_2} ) \vee \bigvee( \mwl{l_1} ) ) )&(\mbox{\ref{eq:orcom}})\notag
  \\ 2 &&  & \vdash ( \bigvee( \mwl{l_1} \mwl{l_2} ) \leftrightarrow ( \bigvee( \mwl{l_1} ) \vee \bigvee( \mwl{l_2} ) ) )&(\mbox{\ref{eq:df-bl.lor2wl}})\notag
  \\ 3 &&  & \vdash ( \bigvee( \mwl{l_2} \mwl{l_1} ) \leftrightarrow ( \bigvee( \mwl{l_2} ) \vee \bigvee( \mwl{l_1} ) ) )&(\mbox{\ref{eq:df-bl.lor2wl}})\notag
  \\ 4 &&  & \vdash ( \bigvee( \mwl{l_1} \mwl{l_2} ) \leftrightarrow \bigvee( \mwl{l_2} \mwl{l_1} ) )&(\mbox{\ref{eq:3bitr4i}},1,2,3)\notag
\end{align}
\end{proof}

\begin{lemma}\blTitle{bl.orpfw} Pull first wff from disjunction wff-list.
\begin{align}
\vdash ( \bigvee( \mw{\varphi} \mwl{l_1} ) \leftrightarrow ( \mw{\varphi} \vee
    \bigvee( \mwl{l_1} ) ) )\label{eq:bl.orpfw}\tag{bl.orpfw}
\end{align}
\end{lemma}

\begin{proof}
\begin{align}
  1 &&  & \vdash ( \bigvee( \mw{\varphi} \mwl{l_1} ) \leftrightarrow ( \bigvee( \mw{\varphi} ) \vee \bigvee( \mwl{l_1} ) ) )&(\mbox{\ref{eq:df-bl.lor2wl}})\notag
  \\ 2 &&  & \vdash ( \bigvee( \mw{\varphi} ) \leftrightarrow \mw{\varphi} )&(\mbox{\ref{eq:df-bl.lor1}})\notag
  \\ 3 &&  & \vdash ( ( \bigvee( \mw{\varphi} ) \vee \bigvee( \mwl{l_1} ) ) \leftrightarrow ( \mw{\varphi} \vee \bigvee( \mwl{l_1} ) ) )&(\mbox{\ref{eq:orbi1i}},2)\notag
  \\ 4 &&  & \vdash ( \bigvee( \mw{\varphi} \mwl{l_1} ) \leftrightarrow ( \mw{\varphi} \vee \bigvee( \mwl{l_1} ) ) )&(\mbox{\ref{eq:bitri}},1,3)\notag
\end{align}
\end{proof}

\begin{lemma}\blTitle{bl.orplw} Pull last wff from disjunction wff-list.
\begin{align}
\vdash ( \bigvee( \mwl{l_1} \mw{\varphi} ) \leftrightarrow ( \bigvee( \mwl{l_1}
    ) \vee \mw{\varphi} ) )\label{eq:bl.orplw}\tag{bl.orplw}
\end{align}
\end{lemma}

\begin{proof}
\begin{align}
  1 &&  & \vdash ( \bigvee( \mwl{l_1} \mw{\varphi} ) \leftrightarrow ( \bigvee( \mwl{l_1} ) \vee \bigvee( \mw{\varphi} ) ) )&(\mbox{\ref{eq:df-bl.lor2wl}})\notag
  \\ 2 &&  & \vdash ( \bigvee( \mw{\varphi} ) \leftrightarrow \mw{\varphi} )&(\mbox{\ref{eq:df-bl.lor1}})\notag
  \\ 3 &&  & \vdash ( ( \bigvee( \mwl{l_1} ) \vee \bigvee( \mw{\varphi} ) ) \leftrightarrow ( \bigvee( \mwl{l_1} ) \vee \mw{\varphi} ) )&(\mbox{\ref{eq:orbi2i}},2)\notag
  \\ 4 &&  & \vdash ( \bigvee( \mwl{l_1} \mw{\varphi} ) \leftrightarrow ( \bigvee( \mwl{l_1} ) \vee \mw{\varphi} ) )&(\mbox{\ref{eq:bitri}},1,3)\notag
\end{align}
\end{proof}

\begin{lemma}\blTitle{bl.orabpf} Absorb first term disjunction wff-list.
\begin{align}
\vdash ( \bigvee( \bigvee( \mwl{l_1} ) \mwl{l_2} ) \leftrightarrow \bigvee(
    \mwl{l_1} \mwl{l_2} ) )\label{eq:bl.orabpf}\tag{bl.orabpf}
\end{align}
\end{lemma}

\begin{proof}
\begin{align}
  1 &&  & \vdash ( \bigvee( \bigvee( \mwl{l_1} ) ) \leftrightarrow \bigvee( \mwl{l_1} ) )&(\mbox{\ref{eq:df-bl.lor1}})\notag
  \\ 2 &&  & \vdash ( ( \bigvee( \bigvee( \mwl{l_1} ) ) \vee \bigvee( \mwl{l_2} ) ) \leftrightarrow ( \bigvee( \mwl{l_1} ) \vee \bigvee( \mwl{l_2} ) ) )&(\mbox{\ref{eq:orbi1i}},1)\notag
  \\ 3 &&  & \vdash ( \bigvee( \bigvee( \mwl{l_1} ) \mwl{l_2} ) \leftrightarrow ( \bigvee( \bigvee( \mwl{l_1} ) ) \vee \bigvee( \mwl{l_2} ) ) )&(\mbox{\ref{eq:df-bl.lor2wl}})\notag
  \\ 4 &&  & \vdash ( \bigvee( \mwl{l_1} \mwl{l_2} ) \leftrightarrow ( \bigvee( \mwl{l_1} ) \vee \bigvee( \mwl{l_2} ) ) )&(\mbox{\ref{eq:df-bl.lor2wl}})\notag
  \\ 5 &&  & \vdash ( \bigvee( \bigvee( \mwl{l_1} ) \mwl{l_2} ) \leftrightarrow \bigvee( \mwl{l_1} \mwl{l_2} ) )&(\mbox{\ref{eq:3bitr4i}},2,3,4)\notag
\end{align}
\end{proof}

\begin{lemma}\blTitle{bl.orabpl} Absorb last term disjunction wff-list.
\begin{align}
\vdash ( \bigvee( \mwl{l_1} \bigvee( \mwl{l_2} ) ) \leftrightarrow \bigvee(
    \mwl{l_1} \mwl{l_2} ) )\label{eq:bl.orabpl}\tag{bl.orabpl}
\end{align}
\end{lemma}

\begin{proof}
\begin{align}
  1 &&  & \vdash ( \bigvee( \bigvee( \mwl{l_2} ) ) \leftrightarrow \bigvee( \mwl{l_2} ) )&(\mbox{\ref{eq:df-bl.lor1}})\notag
  \\ 2 &&  & \vdash ( ( \bigvee( \mwl{l_1} ) \vee \bigvee( \bigvee( \mwl{l_2} ) ) ) \leftrightarrow ( \bigvee( \mwl{l_1} ) \vee \bigvee( \mwl{l_2} ) ) )&(\mbox{\ref{eq:orbi2i}},1)\notag
  \\ 3 &&  & \vdash ( \bigvee( \mwl{l_1} \bigvee( \mwl{l_2} ) ) \leftrightarrow ( \bigvee( \mwl{l_1} ) \vee \bigvee( \bigvee( \mwl{l_2} ) ) ) )&(\mbox{\ref{eq:df-bl.lor2wl}})\notag
  \\ 4 &&  & \vdash ( \bigvee( \mwl{l_1} \mwl{l_2} ) \leftrightarrow ( \bigvee( \mwl{l_1} ) \vee \bigvee( \mwl{l_2} ) ) )&(\mbox{\ref{eq:df-bl.lor2wl}})\notag
  \\ 5 &&  & \vdash ( \bigvee( \mwl{l_1} \bigvee( \mwl{l_2} ) ) \leftrightarrow \bigvee( \mwl{l_1} \mwl{l_2} ) )&(\mbox{\ref{eq:3bitr4i}},2,3,4)\notag
\end{align}
\end{proof}

\begin{lemma}\blTitle{bl.dbo2a} Distribution or-over-and two terms with
wff-lists
\begin{align}
\vdash ( \bigwedge( \mwl{l_1} ( \mw{\varphi} \vee \mw{\psi} ) ) \leftrightarrow
    ( ( \mw{\varphi} \wedge \bigwedge( \mwl{l_1} ) ) \vee ( \mw{\psi} \wedge
    \bigwedge( \mwl{l_1} ) ) ) )\label{eq:bl.dbo2a}\tag{bl.dbo2a}
\end{align}
\end{lemma}

\begin{proof}
\begin{align}
  1 &&  & \vdash ( \bigwedge( \mwl{l_1} ( \mw{\varphi} \vee \mw{\psi} ) ) \leftrightarrow \bigwedge( ( \mw{\varphi} \vee \mw{\psi} ) \mwl{l_1} ) )&(\mbox{\ref{eq:bl.ancomwlwl}})\notag
  \\ 2 &&  & \vdash ( \bigwedge( ( \mw{\varphi} \vee \mw{\psi} ) \mwl{l_1} ) \leftrightarrow ( \bigwedge( ( \mw{\varphi} \vee \mw{\psi} ) ) \wedge \bigwedge( \mwl{l_1} ) ) )&(\mbox{\ref{eq:df-bl.lan2wl}})\notag
  \\ 3 &&  & \vdash ( \bigwedge( \mwl{l_1} ( \mw{\varphi} \vee \mw{\psi} ) ) \leftrightarrow ( \bigwedge( ( \mw{\varphi} \vee \mw{\psi} ) ) \wedge \bigwedge( \mwl{l_1} ) ) )&(\mbox{\ref{eq:bitri}},1,2)\notag
  \\ 4 &&  & \vdash ( \bigwedge( ( \mw{\varphi} \vee \mw{\psi} ) ) \leftrightarrow ( \mw{\varphi} \vee \mw{\psi} ) )&(\mbox{\ref{eq:df-bl.lan1}})\notag
  \\ 5 &&  & \vdash ( ( \bigwedge( ( \mw{\varphi} \vee \mw{\psi} ) ) \wedge \bigwedge( \mwl{l_1} ) ) \leftrightarrow ( ( \mw{\varphi} \vee \mw{\psi} ) \wedge \bigwedge( \mwl{l_1} ) ) )&(\mbox{\ref{eq:anbi1i}},4)\notag
  \\ 6 &&  & \vdash ( ( ( \mw{\varphi} \vee \mw{\psi} ) \wedge \bigwedge( \mwl{l_1} ) ) \leftrightarrow ( ( \mw{\varphi} \wedge \bigwedge( \mwl{l_1} ) ) \vee ( \mw{\psi} \wedge \bigwedge( \mwl{l_1} ) ) ) )&(\mbox{\ref{eq:andir}})\notag
  \\ 7 &&  & \vdash ( \bigwedge( \mwl{l_1} ( \mw{\varphi} \vee \mw{\psi} ) ) \leftrightarrow ( ( \mw{\varphi} \wedge \bigwedge( \mwl{l_1} ) ) \vee ( \mw{\psi} \wedge \bigwedge( \mwl{l_1} ) ) ) )&(\mbox{\ref{eq:3bitri}},3,5,6)\notag
\end{align}
\end{proof}

\begin{lemma}\blTitle{bl.dba2o} Distribution and-over-or two terms with
wff-lists.
\begin{align}
\vdash ( \bigvee( \mwl{l_1} ( \mw{\varphi} \wedge \mw{\psi} ) ) \leftrightarrow
    ( ( \mw{\varphi} \vee \bigvee( \mwl{l_1} ) ) \wedge ( \mw{\psi} \vee
    \bigvee( \mwl{l_1} ) ) ) )\label{eq:bl.dba2o}\tag{bl.dba2o}
\end{align}
\end{lemma}

\begin{proof}
\begin{align}
  1 &&  & \vdash ( \bigvee( \mwl{l_1} ( \mw{\varphi} \wedge \mw{\psi} ) ) \leftrightarrow \bigvee( ( \mw{\varphi} \wedge \mw{\psi} ) \mwl{l_1} ) )&(\mbox{\ref{eq:bl.orcomwlwl}})\notag
  \\ 2 &&  & \vdash ( \bigvee( ( \mw{\varphi} \wedge \mw{\psi} ) \mwl{l_1} ) \leftrightarrow ( \bigvee( ( \mw{\varphi} \wedge \mw{\psi} ) ) \vee \bigvee( \mwl{l_1} ) ) )&(\mbox{\ref{eq:df-bl.lor2wl}})\notag
  \\ 3 &&  & \vdash ( \bigvee( \mwl{l_1} ( \mw{\varphi} \wedge \mw{\psi} ) ) \leftrightarrow ( \bigvee( ( \mw{\varphi} \wedge \mw{\psi} ) ) \vee \bigvee( \mwl{l_1} ) ) )&(\mbox{\ref{eq:bitri}},1,2)\notag
  \\ 4 &&  & \vdash ( \bigvee( ( \mw{\varphi} \wedge \mw{\psi} ) ) \leftrightarrow ( \mw{\varphi} \wedge \mw{\psi} ) )&(\mbox{\ref{eq:df-bl.lor1}})\notag
  \\ 5 &&  & \vdash ( ( \bigvee( ( \mw{\varphi} \wedge \mw{\psi} ) ) \vee \bigvee( \mwl{l_1} ) ) \leftrightarrow ( ( \mw{\varphi} \wedge \mw{\psi} ) \vee \bigvee( \mwl{l_1} ) ) )&(\mbox{\ref{eq:orbi1i}},4)\notag
  \\ 6 &&  & \vdash ( ( ( \mw{\varphi} \wedge \mw{\psi} ) \vee \bigvee( \mwl{l_1} ) ) \leftrightarrow ( ( \mw{\varphi} \vee \bigvee( \mwl{l_1} ) ) \wedge ( \mw{\psi} \vee \bigvee( \mwl{l_1} ) ) ) )&(\mbox{\ref{eq:ordir}})\notag
  \\ 7 &&  & \vdash ( \bigvee( \mwl{l_1} ( \mw{\varphi} \wedge \mw{\psi} ) ) \leftrightarrow ( ( \mw{\varphi} \vee \bigvee( \mwl{l_1} ) ) \wedge ( \mw{\psi} \vee \bigvee( \mwl{l_1} ) ) ) )&(\mbox{\ref{eq:3bitri}},3,5,6)\notag
\end{align}
\end{proof}

\begin{lemma}\blTitle{bl.dbo2awl} Distribution or-over-and with wff-lists.
\begin{align}
\vdash ( \bigwedge( \mwl{l_1} \bigvee( \mwl{l_2} \mwl{l_3} ) ) \leftrightarrow
    ( ( \bigvee( \mwl{l_2} ) \wedge \bigwedge( \mwl{l_1} ) ) \vee ( \bigvee(
    \mwl{l_3} ) \wedge \bigwedge( \mwl{l_1} ) ) )
    )\label{eq:bl.dbo2awl}\tag{bl.dbo2awl}
\end{align}
\end{lemma}

\begin{proof}
\begin{align}
  1 &&  & \vdash ( \bigwedge( \mwl{l_1} \bigvee( \mwl{l_2} \mwl{l_3} ) ) \leftrightarrow ( \bigwedge( \mwl{l_1} ) \wedge \bigvee( \mwl{l_2} \mwl{l_3} ) ) )&(\mbox{\ref{eq:bl.anplw}})\notag
  \\ 2 &&  & \vdash ( ( \bigwedge( \mwl{l_1} ) \wedge \bigvee( \mwl{l_2} \mwl{l_3} ) ) \leftrightarrow ( \bigvee( \mwl{l_2} \mwl{l_3} ) \wedge \bigwedge( \mwl{l_1} ) ) )&(\mbox{\ref{eq:ancom}})\notag
  \\ 3 &&  & \vdash ( ( \bigwedge( \mwl{l_1} \bigvee( \mwl{l_2} \mwl{l_3} ) ) \leftrightarrow ( \bigwedge( \mwl{l_1} ) \wedge \bigvee( \mwl{l_2} \mwl{l_3} ) ) ) \leftrightarrow ( \bigwedge( \mwl{l_1} \bigvee( \mwl{l_2} \mwl{l_3} ) ) \leftrightarrow ( \bigvee( \mwl{l_2} \mwl{l_3} ) \wedge \bigwedge( \mwl{l_1} ) ) ) )&(\mbox{\ref{eq:bibi2i}},2)\notag
  \\ 4 &&  & \vdash ( \bigwedge( \mwl{l_1} \bigvee( \mwl{l_2} \mwl{l_3} ) ) \leftrightarrow ( \bigvee( \mwl{l_2} \mwl{l_3} ) \wedge \bigwedge( \mwl{l_1} ) ) )&(\mbox{\ref{eq:mpbi}},1,3)\notag
  \\ 5 &&  & \vdash ( \bigvee( \mwl{l_2} \mwl{l_3} ) \leftrightarrow ( \bigvee( \mwl{l_2} ) \vee \bigvee( \mwl{l_3} ) ) )&(\mbox{\ref{eq:df-bl.lor2wl}})\notag
  \\ 6 &&  & \vdash ( ( \bigvee( \mwl{l_2} \mwl{l_3} ) \wedge \bigwedge( \mwl{l_1} ) ) \leftrightarrow ( ( \bigvee( \mwl{l_2} ) \vee \bigvee( \mwl{l_3} ) ) \wedge \bigwedge( \mwl{l_1} ) ) )&(\mbox{\ref{eq:anbi1i}},5)\notag
  \\ 7 &&  & \vdash ( ( ( \bigvee( \mwl{l_2} ) \vee \bigvee( \mwl{l_3} ) ) \wedge \bigwedge( \mwl{l_1} ) ) \leftrightarrow ( ( \bigvee( \mwl{l_2} ) \wedge \bigwedge( \mwl{l_1} ) ) \vee ( \bigvee( \mwl{l_3} ) \wedge \bigwedge( \mwl{l_1} ) ) ) )&(\mbox{\ref{eq:andir}})\notag
  \\ 8 &&  & \vdash ( \bigwedge( \mwl{l_1} \bigvee( \mwl{l_2} \mwl{l_3} ) ) \leftrightarrow ( ( \bigvee( \mwl{l_2} ) \wedge \bigwedge( \mwl{l_1} ) ) \vee ( \bigvee( \mwl{l_3} ) \wedge \bigwedge( \mwl{l_1} ) ) ) )&(\mbox{\ref{eq:3bitri}},4,6,7)\notag
\end{align}
\end{proof}

\begin{lemma}\blTitle{bl.dba2owl} Distribution and-over-or with wff lists.
\begin{align}
\vdash ( \bigvee( \mwl{l_1} \bigwedge( \mwl{l_2} \mwl{l_3} ) ) \leftrightarrow
    ( ( \bigwedge( \mwl{l_2} ) \vee \bigvee( \mwl{l_1} ) ) \wedge ( \bigwedge(
    \mwl{l_3} ) \vee \bigvee( \mwl{l_1} ) ) )
    )\label{eq:bl.dba2owl}\tag{bl.dba2owl}
\end{align}
\end{lemma}

\begin{proof}
\begin{align}
  1 &&  & \vdash ( \bigvee( \mwl{l_1} \bigwedge( \mwl{l_2} \mwl{l_3} ) ) \leftrightarrow ( \bigvee( \mwl{l_1} ) \vee \bigwedge( \mwl{l_2} \mwl{l_3} ) ) )&(\mbox{\ref{eq:bl.orplw}})\notag
  \\ 2 &&  & \vdash ( ( \bigvee( \mwl{l_1} ) \vee \bigwedge( \mwl{l_2} \mwl{l_3} ) ) \leftrightarrow ( \bigwedge( \mwl{l_2} \mwl{l_3} ) \vee \bigvee( \mwl{l_1} ) ) )&(\mbox{\ref{eq:orcom}})\notag
  \\ 3 &&  & \vdash ( ( \bigvee( \mwl{l_1} \bigwedge( \mwl{l_2} \mwl{l_3} ) ) \leftrightarrow ( \bigvee( \mwl{l_1} ) \vee \bigwedge( \mwl{l_2} \mwl{l_3} ) ) ) \leftrightarrow ( \bigvee( \mwl{l_1} \bigwedge( \mwl{l_2} \mwl{l_3} ) ) \leftrightarrow ( \bigwedge( \mwl{l_2} \mwl{l_3} ) \vee \bigvee( \mwl{l_1} ) ) ) )&(\mbox{\ref{eq:bibi2i}},2)\notag
  \\ 4 &&  & \vdash ( \bigvee( \mwl{l_1} \bigwedge( \mwl{l_2} \mwl{l_3} ) ) \leftrightarrow ( \bigwedge( \mwl{l_2} \mwl{l_3} ) \vee \bigvee( \mwl{l_1} ) ) )&(\mbox{\ref{eq:mpbi}},1,3)\notag
  \\ 5 &&  & \vdash ( \bigwedge( \mwl{l_2} \mwl{l_3} ) \leftrightarrow ( \bigwedge( \mwl{l_2} ) \wedge \bigwedge( \mwl{l_3} ) ) )&(\mbox{\ref{eq:df-bl.lan2wl}})\notag
  \\ 6 &&  & \vdash ( ( \bigwedge( \mwl{l_2} \mwl{l_3} ) \vee \bigvee( \mwl{l_1} ) ) \leftrightarrow ( ( \bigwedge( \mwl{l_2} ) \wedge \bigwedge( \mwl{l_3} ) ) \vee \bigvee( \mwl{l_1} ) ) )&(\mbox{\ref{eq:orbi1i}},5)\notag
  \\ 7 &&  & \vdash ( ( ( \bigwedge( \mwl{l_2} ) \wedge \bigwedge( \mwl{l_3} ) ) \vee \bigvee( \mwl{l_1} ) ) \leftrightarrow ( ( \bigwedge( \mwl{l_2} ) \vee \bigvee( \mwl{l_1} ) ) \wedge ( \bigwedge( \mwl{l_3} ) \vee \bigvee( \mwl{l_1} ) ) ) )&(\mbox{\ref{eq:ordir}})\notag
  \\ 8 &&  & \vdash ( \bigvee( \mwl{l_1} \bigwedge( \mwl{l_2} \mwl{l_3} ) ) \leftrightarrow ( ( \bigwedge( \mwl{l_2} ) \vee \bigvee( \mwl{l_1} ) ) \wedge ( \bigwedge( \mwl{l_3} ) \vee \bigvee( \mwl{l_1} ) ) ) )&(\mbox{\ref{eq:3bitri}},4,6,7)\notag
\end{align}
\end{proof}

\begin{metadefinition}\blTitle{df-bl.sbw} define substitution of a wff by a wff
in a wff
\begin{align}
\vdash ( \overline{[} \mw{\psi} / \mw{\chi} \overline{]} \mw{\varphi}
    \leftrightarrow ( ( \mw{\chi} \leftrightarrow \mw{\psi} ) \rightarrow
    \mw{\varphi} ) )\label{eq:df-bl.sbw}\tag{df-bl.sbw}
\end{align}
\end{metadefinition}

\begin{lemma}\blTitle{bl.sbwid} Substitution by the Same wff has no Effect
\begin{align}
\vdash ( \overline{[} \mw{\psi} / \mw{\psi} \overline{]} \mw{\varphi}
    \leftrightarrow \mw{\varphi} )\label{eq:bl.sbwid}\tag{bl.sbwid}
\end{align}
\end{lemma}

\begin{proof}
\begin{align}
  1 &&  & \vdash ( \overline{[} \mw{\psi} / \mw{\psi} \overline{]} \mw{\varphi} \leftrightarrow ( ( \mw{\psi} \leftrightarrow \mw{\psi} ) \rightarrow \mw{\varphi} ) )&(\mbox{\ref{eq:df-bl.sbw}})\notag
  \\ 2 &&  & \vdash ( \mw{\psi} \leftrightarrow \mw{\psi} )&(\mbox{\ref{eq:biid}})\notag
  \\ 3 &&  & \vdash ( \mw{\varphi} \leftrightarrow ( ( \mw{\psi} \leftrightarrow \mw{\psi} ) \rightarrow \mw{\varphi} ) )&(\mbox{\ref{eq:a1bi}},2)\notag
  \\ 4 &&  & \vdash ( ( ( \mw{\psi} \leftrightarrow \mw{\psi} ) \rightarrow \mw{\varphi} ) \leftrightarrow \mw{\varphi} )&(\mbox{\ref{eq:bicomi}},3)\notag
  \\ 5 &&  & \vdash ( \overline{[} \mw{\psi} / \mw{\psi} \overline{]} \mw{\varphi} \leftrightarrow \mw{\varphi} )&(\mbox{\ref{eq:bitri}},1,4)\notag
\end{align}
\end{proof}

\begin{lemma}\blTitle{bl.bisbwl} Equivalence of wff Substitution on Left
\begin{align}
\vdash ( \mw{\chi} \leftrightarrow \mw{\psi} )\quad\Rightarrow\quad \vdash (
    \overline{[} \mw{\psi} / \mw{\chi} \overline{]} \mw{\varphi}
    \leftrightarrow \mw{\varphi} )\label{eq:bl.bisbwl}\tag{bl.bisbwl}
\end{align}
\end{lemma}

\begin{proof}
\begin{align}
  1 &&  & \vdash ( \overline{[} \mw{\psi} / \mw{\chi} \overline{]} \mw{\varphi} \leftrightarrow ( ( \mw{\chi} \leftrightarrow \mw{\psi} ) \rightarrow \mw{\varphi} ) )&(\mbox{\ref{eq:df-bl.sbw}})\notag
  \\ 2 &&  & \vdash ( \mw{\chi} \leftrightarrow \mw{\psi} )&\text{Hyp~bisbwl.1}\notag%bl.bisbwl.1
  \\ 3 &&  & \vdash ( \mw{\varphi} \leftrightarrow ( ( \mw{\chi} \leftrightarrow \mw{\psi} ) \rightarrow \mw{\varphi} ) )&(\mbox{\ref{eq:a1bi}},2)\notag
  \\ 4 &&  & \vdash ( \overline{[} \mw{\psi} / \mw{\chi} \overline{]} \mw{\varphi} \leftrightarrow \mw{\varphi} )&(\mbox{\ref{eq:bitr4i}},1,3)\notag
\end{align}
\end{proof}

\begin{lemma}\blTitle{bl.bisbwr} Equivalence of wff Substitution on Right
\begin{align}
\vdash ( \mw{\chi} \leftrightarrow \mw{\psi} )\quad\Rightarrow\quad \vdash (
    \mw{\varphi} \leftrightarrow \overline{[} \mw{\psi} / \mw{\chi}
    \overline{]} \mw{\varphi} )\label{eq:bl.bisbwr}\tag{bl.bisbwr}
\end{align}
\end{lemma}

\begin{proof}
\begin{align}
  1 &&  & \vdash ( \mw{\chi} \leftrightarrow \mw{\psi} )&\text{Hyp~bisbwr.1}\notag%bl.bisbwr.1
  \\ 2 &&  & \vdash ( \overline{[} \mw{\psi} / \mw{\chi} \overline{]} \mw{\varphi} \leftrightarrow \mw{\varphi} )&(\mbox{\ref{eq:bl.bisbwl}},1)\notag
  \\ 3 &&  & \vdash ( \mw{\varphi} \leftrightarrow \overline{[} \mw{\psi} / \mw{\chi} \overline{]} \mw{\varphi} )&(\mbox{\ref{eq:bicomi}},2)\notag
\end{align}
\end{proof}

\begin{lemma}\blTitle{bl.sylsbw} Syllogism of wff Substitution
Introduction into Inference
\begin{align}
\vdash ( \mw{\varphi} \rightarrow \mw{\psi} )\quad\&\quad \vdash ( \mw{\chi}
    \leftrightarrow \mw{\theta} )\quad\Rightarrow\quad \vdash ( \overline{[}
    \mw{\theta} / \mw{\chi} \overline{]} \mw{\varphi} \rightarrow \overline{[}
    \mw{\theta} / \mw{\chi} \overline{]} \mw{\psi}
    )\label{eq:bl.sylsbw}\tag{bl.sylsbw}
\end{align}
\end{lemma}

\begin{proof}
\begin{align}
  1 &&  & \vdash ( \mw{\varphi} \rightarrow \mw{\psi} )&\text{Hyp~sylsbw.1}\notag%bl.sylsbw.1
  \\ 2 &&  & \vdash ( ( ( \mw{\chi} \leftrightarrow \mw{\theta} ) \rightarrow \mw{\varphi} ) \rightarrow ( ( \mw{\chi} \leftrightarrow \mw{\theta} ) \rightarrow \mw{\psi} ) )&(\mbox{\ref{eq:imim2i}},1)\notag
  \\ 3 &&  & \vdash ( \overline{[} \mw{\theta} / \mw{\chi} \overline{]} \mw{\varphi} \leftrightarrow ( ( \mw{\chi} \leftrightarrow \mw{\theta} ) \rightarrow \mw{\varphi} ) )&(\mbox{\ref{eq:df-bl.sbw}})\notag
  \\ 4 &&  & \vdash ( \overline{[} \mw{\theta} / \mw{\chi} \overline{]} \mw{\psi} \leftrightarrow ( ( \mw{\chi} \leftrightarrow \mw{\theta} ) \rightarrow \mw{\psi} ) )&(\mbox{\ref{eq:df-bl.sbw}})\notag
  \\ 5 &&  & \vdash ( \overline{[} \mw{\theta} / \mw{\chi} \overline{]} \mw{\varphi} \rightarrow \overline{[} \mw{\theta} / \mw{\chi} \overline{]} \mw{\psi} )&(\mbox{\ref{eq:3imtr4i}},2,3,4)\notag
\end{align}
\end{proof}

\begin{lemma}\blTitle{bl.sbwsyl} Syllogism of wff Substitution Elimination
from Inference
\begin{align}
\vdash ( \overline{[} \mw{\theta} / \mw{\chi} \overline{]} \mw{\varphi}
    \rightarrow \overline{[} \mw{\theta} / \mw{\chi} \overline{]} \mw{\psi}
    )\quad\&\quad \vdash ( \mw{\chi} \leftrightarrow \mw{\theta}
    )\quad\Rightarrow\quad \vdash ( \mw{\varphi} \rightarrow \mw{\psi}
    )\label{eq:bl.sbwsyl}\tag{bl.sbwsyl}
\end{align}
\end{lemma}

\begin{proof}
\begin{align}
  1 &&  & \vdash ( \mw{\chi} \leftrightarrow \mw{\theta} )&\text{Hyp~sbwsyl.2}\notag%bl.sbwsyl.2
  \\ 2 &&  & \vdash ( \overline{[} \mw{\theta} / \mw{\chi} \overline{]} \mw{\varphi} \rightarrow \overline{[} \mw{\theta} / \mw{\chi} \overline{]} \mw{\psi} )&\text{Hyp~sbwsyl.1}\notag%bl.sbwsyl.1
  \\ 3 &&  & \vdash ( \overline{[} \mw{\theta} / \mw{\chi} \overline{]} \mw{\varphi} \leftrightarrow ( ( \mw{\chi} \leftrightarrow \mw{\theta} ) \rightarrow \mw{\varphi} ) )&(\mbox{\ref{eq:df-bl.sbw}})\notag
  \\ 4 &&  & \vdash ( \overline{[} \mw{\theta} / \mw{\chi} \overline{]} \mw{\psi} \leftrightarrow ( ( \mw{\chi} \leftrightarrow \mw{\theta} ) \rightarrow \mw{\psi} ) )&(\mbox{\ref{eq:df-bl.sbw}})\notag
  \\ 5 &&  & \vdash ( ( ( \mw{\chi} \leftrightarrow \mw{\theta} ) \rightarrow \mw{\varphi} ) \rightarrow ( ( \mw{\chi} \leftrightarrow \mw{\theta} ) \rightarrow \mw{\psi} ) )&(\mbox{\ref{eq:3imtr3i}},2,3,4)\notag
  \\ 6 &&  & \vdash ( ( \mw{\chi} \leftrightarrow \mw{\theta} ) \rightarrow ( \mw{\varphi} \rightarrow \mw{\psi} ) )&(\mbox{\ref{eq:pm2.86i}},5)\notag
  \\ 7 &&  & \vdash ( \mw{\varphi} \rightarrow \mw{\psi} )&(\mbox{\ref{eq:ax-mp}},1,6)\notag
\end{align}
\end{proof}

\begin{lemma}\blTitle{bl.ctao} Common Term Between And-List and Or-List
(right)
\begin{align}
\vdash ( \bigwedge( \mwl{l_1} \mw{\varphi} ) \rightarrow \bigvee( \mwl{l_2}
    \mw{\varphi} ) )\label{eq:bl.ctao}\tag{bl.ctao}
\end{align}
\end{lemma}

\begin{proof}
\begin{align}
  1 &&  & \vdash ( ( \bigwedge( \mwl{l_1} ) \wedge \mw{\varphi} ) \rightarrow ( \bigvee( \mwl{l_2} ) \vee \mw{\varphi} ) )&(\mbox{\ref{eq:animorr}})\notag
  \\ 2 &&  & \vdash ( \bigwedge( \mwl{l_1} \mw{\varphi} ) \leftrightarrow ( \bigwedge( \mwl{l_1} ) \wedge \mw{\varphi} ) )&(\mbox{\ref{eq:bl.anplw}})\notag
  \\ 3 &&  & \vdash ( \bigvee( \mwl{l_2} \mw{\varphi} ) \leftrightarrow ( \bigvee( \mwl{l_2} ) \vee \mw{\varphi} ) )&(\mbox{\ref{eq:bl.orplw}})\notag
  \\ 4 &&  & \vdash ( \bigwedge( \mwl{l_1} \mw{\varphi} ) \rightarrow \bigvee( \mwl{l_2} \mw{\varphi} ) )&(\mbox{\ref{eq:3imtr4i}},1,2,3)\notag
\end{align}
\end{proof}

\begin{lemma}\blTitle{bl.ctaol} Common Term Between And-List and Or-List
(left)
\begin{align}
\vdash ( \bigwedge( \mw{\varphi} \mwl{l_1} ) \rightarrow \bigvee( \mw{\varphi}
    \mwl{l_2} ) )\label{eq:bl.ctaol}\tag{bl.ctaol}
\end{align}
\end{lemma}

\begin{proof}
\begin{align}
  1 &&  & \vdash ( \bigwedge( \mwl{l_1} \mwl{l_2} ) \leftrightarrow ( \bigwedge( \mwl{l_1} ) \wedge \bigwedge( \mwl{l_2} ) ) )&(\mbox{\ref{eq:df-bl.lan2wl}})\notag
  \\ 2 &&  & \vdash ( ( \bigwedge( \mwl{l_1} ) \wedge \bigwedge( \mwl{l_2} ) ) \rightarrow ( \bigvee( \mwl{l_3} ) \vee \bigwedge( \mwl{l_2} ) ) )&(\mbox{\ref{eq:animorr}})\notag
  \\ 3 &&  & \vdash ( \bigwedge( \mwl{l_1} \mwl{l_2} ) \rightarrow ( \bigvee( \mwl{l_3} ) \vee \bigwedge( \mwl{l_2} ) ) )&(\mbox{\ref{eq:sylbi}},1,2)\notag
  \\ 4 &&  & \vdash ( \bigvee( \mwl{l_3} \bigwedge( \mwl{l_2} ) ) \leftrightarrow ( \bigvee( \mwl{l_3} ) \vee \bigwedge( \mwl{l_2} ) ) )&(\mbox{\ref{eq:bl.orplw}})\notag
  \\ 5 &&  & \vdash ( \bigwedge( \mwl{l_1} \mwl{l_2} ) \rightarrow \bigvee( \mwl{l_3} \bigwedge( \mwl{l_2} ) ) )&(\mbox{\ref{eq:sylibr}},3,4)\notag
\end{align}
\end{proof}

\begin{lemma}\blTitle{bl.animporanl} Premise Has All Terms of Conjunction
Within Disjunction (left)
\begin{align}
\vdash ( \bigwedge( \mwl{l_2} \mwl{l_1} ) \rightarrow \bigvee( \bigwedge(
    \mwl{l_2} ) \mwl{l_3} ) )\label{eq:bl.animporanl}\tag{bl.animporanl}
\end{align}
\end{lemma}

\begin{proof}
\begin{align}
  1 &&  & \vdash ( ( \bigwedge( \mwl{l_2} ) \wedge \bigwedge( \mwl{l_1} ) ) \rightarrow ( \bigwedge( \mwl{l_2} ) \vee \bigvee( \mwl{l_3} ) ) )&(\mbox{\ref{eq:animorl}})\notag
  \\ 2 &&  & \vdash ( \bigwedge( \mwl{l_2} \mwl{l_1} ) \leftrightarrow ( \bigwedge( \mwl{l_2} ) \wedge \bigwedge( \mwl{l_1} ) ) )&(\mbox{\ref{eq:df-bl.lan2wl}})\notag
  \\ 3 &&  & \vdash ( \bigvee( \bigwedge( \mwl{l_2} ) \mwl{l_3} ) \leftrightarrow ( \bigwedge( \mwl{l_2} ) \vee \bigvee( \mwl{l_3} ) ) )&(\mbox{\ref{eq:bl.orpfw}})\notag
  \\ 4 &&  & \vdash ( \bigwedge( \mwl{l_2} \mwl{l_1} ) \rightarrow \bigvee( \bigwedge( \mwl{l_2} ) \mwl{l_3} ) )&(\mbox{\ref{eq:3imtr4i}},1,2,3)\notag
\end{align}
\end{proof}

\begin{lemma}\blTitle{bl.orcwl} Or-Introduction Schema of wff (right)
\begin{align}
\vdash ( \mw{\varphi} \rightarrow \bigvee( \mwl{l_1} \mw{\varphi} )
    )\label{eq:bl.orcwl}\tag{bl.orcwl}
\end{align}
\end{lemma}

\begin{proof}
\begin{align}
  1 &&  & \vdash ( \mw{\varphi} \rightarrow ( \bigvee( \mwl{l_1} ) \vee \mw{\varphi} ) )&(\mbox{\ref{eq:olc}})\notag
  \\ 2 &&  & \vdash ( \bigvee( \mwl{l_1} \mw{\varphi} ) \leftrightarrow ( \bigvee( \mwl{l_1} ) \vee \mw{\varphi} ) )&(\mbox{\ref{eq:bl.orplw}})\notag
  \\ 3 &&  & \vdash ( \mw{\varphi} \rightarrow \bigvee( \mwl{l_1} \mw{\varphi} ) )&(\mbox{\ref{eq:sylibr}},1,2)\notag
\end{align}
\end{proof}

\begin{lemma}\blTitle{bl.orcwll} Or-Introduction Schema of wff (left)
\begin{align}
\vdash ( \mw{\varphi} \rightarrow \bigvee( \mw{\varphi} \mwl{l_1} )
    )\label{eq:bl.orcwll}\tag{bl.orcwll}
\end{align}
\end{lemma}

\begin{proof}
\begin{align}
  1 &&  & \vdash ( \mw{\varphi} \rightarrow ( \mw{\varphi} \vee \bigvee( \mwl{l_1} ) ) )&(\mbox{\ref{eq:orc}})\notag
  \\ 2 &&  & \vdash ( \bigvee( \mw{\varphi} \mwl{l_1} ) \leftrightarrow ( \mw{\varphi} \vee \bigvee( \mwl{l_1} ) ) )&(\mbox{\ref{eq:bl.orpfw}})\notag
  \\ 3 &&  & \vdash ( \mw{\varphi} \rightarrow \bigvee( \mw{\varphi} \mwl{l_1} ) )&(\mbox{\ref{eq:sylibr}},1,2)\notag
\end{align}
\end{proof}

\begin{lemma}\blTitle{bl.aisl} And-Introduction Schema of wff (left)
\begin{align}
\vdash ( \bigwedge( \mwl{l_1} \mw{\varphi} ) \rightarrow \mw{\varphi}
    )\label{eq:bl.aisl}\tag{bl.aisl}
\end{align}
\end{lemma}

\begin{proof}
\begin{align}
  1 &&  & \vdash ( \bigwedge( \mwl{l_1} \mw{\varphi} ) \leftrightarrow ( \bigwedge( \mwl{l_1} ) \wedge \mw{\varphi} ) )&(\mbox{\ref{eq:bl.anplw}})\notag
  \\ 2 &&  & \vdash ( ( \bigwedge( \mwl{l_1} ) \wedge \mw{\varphi} ) \rightarrow \mw{\varphi} )&(\mbox{\ref{eq:simpr}})\notag
  \\ 3 &&  & \vdash ( \bigwedge( \mwl{l_1} \mw{\varphi} ) \rightarrow \mw{\varphi} )&(\mbox{\ref{eq:sylbi}},1,2)\notag
\end{align}
\end{proof}

\begin{lemma}\blTitle{bl.aisr} And-Introduction Schema of wff (right)
\begin{align}
\vdash ( \bigwedge( \mw{\varphi} \mwl{l_1} ) \rightarrow \mw{\varphi}
    )\label{eq:bl.aisr}\tag{bl.aisr}
\end{align}
\end{lemma}

\begin{proof}
\begin{align}
  1 &&  & \vdash ( \bigwedge( \mw{\varphi} \mwl{l_1} ) \leftrightarrow ( \mw{\varphi} \wedge \bigwedge( \mwl{l_1} ) ) )&(\mbox{\ref{eq:bl.anpfw}})\notag
  \\ 2 &&  & \vdash ( ( \mw{\varphi} \wedge \bigwedge( \mwl{l_1} ) ) \rightarrow \mw{\varphi} )&(\mbox{\ref{eq:simpl}})\notag
  \\ 3 &&  & \vdash ( \bigwedge( \mw{\varphi} \mwl{l_1} ) \rightarrow \mw{\varphi} )&(\mbox{\ref{eq:sylbi}},1,2)\notag
\end{align}
\end{proof}

\begin{lemma}\blTitle{bl.aiswll} And-Introduction Schema of wff-List
(left)
\begin{align}
\vdash ( \bigwedge( \mwl{l_1} \mwl{l_2} ) \rightarrow \bigwedge( \mwl{l_2} )
    )\label{eq:bl.aiswll}\tag{bl.aiswll}
\end{align}
\end{lemma}

\begin{proof}
\begin{align}
  1 &&  & \vdash ( \bigwedge( \mwl{l_1} \mwl{l_2} ) \leftrightarrow ( \bigwedge( \mwl{l_1} ) \wedge \bigwedge( \mwl{l_2} ) ) )&(\mbox{\ref{eq:df-bl.lan2wl}})\notag
  \\ 2 &&  & \vdash ( ( \bigwedge( \mwl{l_1} ) \wedge \bigwedge( \mwl{l_2} ) ) \rightarrow \bigwedge( \mwl{l_2} ) )&(\mbox{\ref{eq:simpr}})\notag
  \\ 3 &&  & \vdash ( \bigwedge( \mwl{l_1} \mwl{l_2} ) \rightarrow \bigwedge( \mwl{l_2} ) )&(\mbox{\ref{eq:sylbi}},1,2)\notag
\end{align}
\end{proof}

\begin{lemma}\blTitle{bl.aiswlr} And-Introduction Schema of wff-List
(left)
\begin{align}
\vdash ( \bigwedge( \mwl{l_1} \mwl{l_2} ) \rightarrow \bigwedge( \mwl{l_1} )
    )\label{eq:bl.aiswlr}\tag{bl.aiswlr}
\end{align}
\end{lemma}

\begin{proof}
\begin{align}
  1 &&  & \vdash ( \bigwedge( \mwl{l_1} \mwl{l_2} ) \leftrightarrow ( \bigwedge( \mwl{l_1} ) \wedge \bigwedge( \mwl{l_2} ) ) )&(\mbox{\ref{eq:df-bl.lan2wl}})\notag
  \\ 2 &&  & \vdash ( ( \bigwedge( \mwl{l_1} ) \wedge \bigwedge( \mwl{l_2} ) ) \rightarrow \bigwedge( \mwl{l_1} ) )&(\mbox{\ref{eq:simpl}})\notag
  \\ 3 &&  & \vdash ( \bigwedge( \mwl{l_1} \mwl{l_2} ) \rightarrow \bigwedge( \mwl{l_1} ) )&(\mbox{\ref{eq:sylbi}},1,2)\notag
\end{align}
\end{proof}

\begin{metadefinition}\blTitle{bl.cla} a single class is a classlist
\begin{align}
\mathsf{CL} A\label{eq:bl.cla}\tag{bl.cla}
\end{align}
\end{metadefinition}

\begin{metadefinition}\blTitle{bl.cln} a classlist followed by a comma and a
class is a class list
\begin{align}
\mathsf{CL} \msc{csl_1} , \msc{csl_2}\label{eq:bl.cln}\tag{bl.cln}
\end{align}
\end{metadefinition}

\begin{metadefinition}\blTitle{df-bl.cl} make a class list into a class
\begin{align}
{\rm class} \langle \msc{csl_1} \rangle\label{eq:df-bl.cl}\tag{df-bl.cl}
\end{align}
\end{metadefinition}

\begin{metadefinition}\blTitle{df-bl.cl0} allow 0-Tuple
\begin{align}
{\rm class} \langle \rangle\label{eq:df-bl.cl0}\tag{df-bl.cl0}
\end{align}
\end{metadefinition}

\begin{metadefinition}\blTitle{df-bl.2teqop} correspondence of 2-Tuple to
ordered pair
\begin{align}
\vdash \langle A , B \rangle = \langle A , B
    \rangle\label{eq:df-bl.2teqop}\tag{df-bl.2teqop}
\end{align}
\end{metadefinition}

\begin{metadefinition}\blTitle{bl.cadd} Multi-term Addition is a Class
\begin{align}
{\rm class} +( \msc{csl_1} )\label{eq:bl.cadd}\tag{bl.cadd}
\end{align}
\end{metadefinition}

\begin{metadefinition}\blTitle{df-bl.add1} Multi-term Addition of Single Value
is that Value.
\begin{align}
\vdash ( A \in \mathbb{C} \rightarrow +( A ) = A
    )\label{eq:df-bl.add1}\tag{df-bl.add1}
\end{align}
\end{metadefinition}

\begin{metadefinition}\blTitle{df-bl.add2} Multi-term Addition of Two Class
Lists is their Sum.
\begin{align}
\vdash +( \msc{csl_1} , \msc{csl_2} ) = ( +( \msc{csl_1} ) + +( \msc{csl_2} )
    )\label{eq:df-bl.add2}\tag{df-bl.add2}
\end{align}
\end{metadefinition}

\begin{lemma}\blTitle{bl.addcom} Multi-term Addition Commutes
\begin{align}
\vdash +( \msc{csl_1} , \msc{csl_2} ) = +( \msc{csl_2} , \msc{csl_1}
    )\label{eq:bl.addcom}\tag{bl.addcom}
\end{align}
\end{lemma}

\begin{proof}
\begin{align}
  1 &&  & \vdash ( +( \msc{csl_1} ) + +( \msc{csl_2} ) ) = ( +( \msc{csl_2} ) + +( \msc{csl_1} ) )&(\mbox{\ref{eq:addcomgi}})\notag
  \\ 2 &&  & \vdash +( \msc{csl_1} , \msc{csl_2} ) = ( +( \msc{csl_1} ) + +( \msc{csl_2} ) )&(\mbox{\ref{eq:df-bl.add2}})\notag
  \\ 3 &&  & \vdash +( \msc{csl_2} , \msc{csl_1} ) = ( +( \msc{csl_2} ) + +( \msc{csl_1} ) )&(\mbox{\ref{eq:df-bl.add2}})\notag
  \\ 4 &&  & \vdash +( \msc{csl_1} , \msc{csl_2} ) = +( \msc{csl_2} , \msc{csl_1} )&(\mbox{\ref{eq:3eqtr4i}},1,2,3)\notag
\end{align}
\end{proof}

\begin{lemma}\blTitle{bl.adeqi} adding equals
\begin{align}
\vdash A = B\quad\&\quad \vdash C = D\quad\Rightarrow\quad \vdash ( A + C ) = (
    B + D )\label{eq:bl.adeqi}\tag{bl.adeqi}
\end{align}
\end{lemma}

\begin{proof}
\begin{align}
  1 &&  & \vdash A = B&\text{Hyp~adeqi.1}\notag%bl.adeqi.1
  \\ 2 &&  & \vdash C = D&\text{Hyp~adeqi.2}\notag%bl.adeqi.2
  \\ 3 &&  & \vdash ( A + C ) = ( B + D )&(\mbox{\ref{eq:oveq12i}},1,2)\notag
\end{align}
\end{proof}

\begin{lemma}\blTitle{bl.add2i} Multi-term Addition of Two Values Lists is
their Sum.
\begin{align}
\vdash A \in \mathbb{C}\quad\&\quad \vdash B \in
    \mathbb{C}\quad\Rightarrow\quad \vdash +( A , B ) = ( A + B
    )\label{eq:bl.add2i}\tag{bl.add2i}
\end{align}
\end{lemma}

\begin{proof}
\begin{align}
  1 &&  & \vdash +( A , B ) = ( +( A ) + +( B ) )&(\mbox{\ref{eq:df-bl.add2}})\notag
  \\ 2 &&  & \vdash A \in \mathbb{C}&\text{Hyp~add2i.1}\notag%bl.add2i.1
  \\ 3 &&  & \vdash ( A \in \mathbb{C} \rightarrow +( A ) = A )&(\mbox{\ref{eq:df-bl.add1}})\notag
  \\ 4 &&  & \vdash +( A ) = A&(\mbox{\ref{eq:ax-mp}},2,3)\notag
  \\ 5 &&  & \vdash B \in \mathbb{C}&\text{Hyp~add2i.2}\notag%bl.add2i.2
  \\ 6 &&  & \vdash ( B \in \mathbb{C} \rightarrow +( B ) = B )&(\mbox{\ref{eq:df-bl.add1}})\notag
  \\ 7 &&  & \vdash +( B ) = B&(\mbox{\ref{eq:ax-mp}},5,6)\notag
  \\ 8 &&  & \vdash ( +( A ) + +( B ) ) = ( A + B )&(\mbox{\ref{eq:bl.adeqi}},4,7)\notag
  \\ 9 &&  & \vdash +( A , B ) = ( A + B )&(\mbox{\ref{eq:eqtri}},1,8)\notag
\end{align}
\end{proof}

\begin{lemma}\blTitle{bl.add3i} Multi-term Addition of Three Values
\begin{align}
\vdash A \in \mathbb{C}\quad\&\quad \vdash B \in \mathbb{C}\quad\&\quad \vdash
    C \in \mathbb{C}\quad\Rightarrow\quad \vdash +( A , B , C ) = ( A + ( B + C
    ) )\label{eq:bl.add3i}\tag{bl.add3i}
\end{align}
\end{lemma}

\begin{proof}
\begin{align}
  1 &&  & \vdash +( A , B , C ) = ( +( A ) + +( B , C ) )&(\mbox{\ref{eq:df-bl.add2}})\notag
  \\ 2 &&  & \vdash A \in \mathbb{C}&\text{Hyp~add3i.1}\notag%bl.add3i.1
  \\ 3 &&  & \vdash ( A \in \mathbb{C} \rightarrow +( A ) = A )&(\mbox{\ref{eq:df-bl.add1}})\notag
  \\ 4 &&  & \vdash +( A ) = A&(\mbox{\ref{eq:ax-mp}},2,3)\notag
  \\ 5 &&  & \vdash B \in \mathbb{C}&\text{Hyp~add3i.2}\notag%bl.add3i.2
  \\ 6 &&  & \vdash C \in \mathbb{C}&\text{Hyp~add3i.3}\notag%bl.add3i.3
  \\ 7 &&  & \vdash +( B , C ) = ( B + C )&(\mbox{\ref{eq:bl.add2i}},5,6)\notag
  \\ 8 &&  & \vdash ( +( A ) + +( B , C ) ) = ( A + ( B + C ) )&(\mbox{\ref{eq:bl.adeqi}},4,7)\notag
  \\ 9 &&  & \vdash +( A , B , C ) = ( A + ( B + C ) )&(\mbox{\ref{eq:eqtri}},1,8)\notag
\end{align}
\end{proof}

\begin{lemma}\blTitle{bl.add4i} Multi-term Addition of Four Values
\begin{align}
\vdash A \in \mathbb{C}\quad\&\quad \vdash B \in \mathbb{C}\quad\&\quad \vdash
    C \in \mathbb{C}\quad\&\quad \vdash D \in \mathbb{C}\quad\Rightarrow\quad
    \vdash +( A , B , C , D ) = ( ( A + B ) + ( C + D )
    )\label{eq:bl.add4i}\tag{bl.add4i}
\end{align}
\end{lemma}

\begin{proof}
\begin{align}
  1 &&  & \vdash +( A , B , C , D ) = ( +( A , B ) + +( C , D ) )&(\mbox{\ref{eq:df-bl.add2}})\notag
  \\ 2 &&  & \vdash A \in \mathbb{C}&\text{Hyp~add4i.1}\notag%bl.add4i.1
  \\ 3 &&  & \vdash B \in \mathbb{C}&\text{Hyp~add4i.2}\notag%bl.add4i.2
  \\ 4 &&  & \vdash +( A , B ) = ( A + B )&(\mbox{\ref{eq:bl.add2i}},2,3)\notag
  \\ 5 &&  & \vdash C \in \mathbb{C}&\text{Hyp~add4i.3}\notag%bl.add4i.3
  \\ 6 &&  & \vdash D \in \mathbb{C}&\text{Hyp~add4i.4}\notag%bl.add4i.4
  \\ 7 &&  & \vdash +( C , D ) = ( C + D )&(\mbox{\ref{eq:bl.add2i}},5,6)\notag
  \\ 8 &&  & \vdash ( +( A , B ) + +( C , D ) ) = ( ( A + B ) + ( C + D ) )&(\mbox{\ref{eq:bl.adeqi}},4,7)\notag
  \\ 9 &&  & \vdash +( A , B , C , D ) = ( ( A + B ) + ( C + D ) )&(\mbox{\ref{eq:eqtri}},1,8)\notag
\end{align}
\end{proof}

\begin{lemma}\blTitle{bl.addcom2} Addition Two Values Commutes
\begin{align}
\vdash A \in \mathbb{C}\quad\&\quad \vdash B \in
    \mathbb{C}\quad\Rightarrow\quad \vdash +( A , B ) = +( B , A
    )\label{eq:bl.addcom2}\tag{bl.addcom2}
\end{align}
\end{lemma}

\begin{proof}
\begin{align}
  1 &&  & \vdash +( A , B ) = +( B , A )&(\mbox{\ref{eq:bl.addcom}})\notag
\end{align}
\end{proof}

%\begin{lemma}\blTitle{bl.add3i} Multi-term Addition of Three Values
%\begin{align}
%\vdash A \in \mathbb{C}\quad\&\quad \vdash B \in \mathbb{C}\quad\&\quad \vdash
%    C \in \mathbb{C}\quad\Rightarrow\quad \vdash +( A , B , C ) = ( A + ( B + C
%    ) )\label{eq:bl.add3i}\tag{bl.add3i}
%\end{align}
%\end{lemma}

%\begin{proof}
%\begin{align}
%  1 &&  & \vdash +( A , B , C ) = ( +( A ) + +( B , C ) )&(\mbox{\ref{eq:df-bl.add2}})\notag
%  \\ 2 &&  & \vdash A \in \mathbb{C}&\text{Hyp~add3i.1}\notag%bl.add3i.1
%  \\ 3 &&  & \vdash ( A \in \mathbb{C} \rightarrow +( A ) = A )&(\mbox{\ref{eq:df-bl.add1}})\notag
%  \\ 4 &&  & \vdash +( A ) = A&(\mbox{\ref{eq:ax-mp}},2,3)\notag
%  \\ 5 &&  & \vdash B \in \mathbb{C}&\text{Hyp~add3i.2}\notag%bl.add3i.2
%  \\ 6 &&  & \vdash C \in \mathbb{C}&\text{Hyp~add3i.3}\notag%bl.add3i.3
%  \\ 7 &&  & \vdash +( B , C ) = ( B + C )&(\mbox{\ref{eq:bl.add2i}},5,6)\notag
%  \\ 8 &&  & \vdash ( +( A ) + +( B , C ) ) = ( A + ( B + C ) )&(\mbox{\ref{eq:bl.adeqi}},4,7)\notag
%  \\ 9 &&  & \vdash +( A , B , C ) = ( A + ( B + C ) )&(\mbox{\ref{eq:eqtri}},1,8)\notag
%\end{align}
%\end{proof}

\begin{lemma}\blTitle{bl.addcom3acb} Addition of Three Values Commutes
\begin{align}
\vdash A \in \mathbb{C}\quad\&\quad \vdash B \in \mathbb{C}\quad\&\quad \vdash
    C \in \mathbb{C}\quad\Rightarrow\quad \vdash +( A , B , C ) = +( A , C , B
    )\label{eq:bl.addcom3acb}\tag{bl.addcom3acb}
\end{align}
\end{lemma}

\begin{proof}
\begin{align}
  1 &&  & \vdash ( B + C ) = ( C + B )&(\mbox{\ref{eq:addcomgi}})\notag
  \\ 2 &&  & \vdash ( A + ( B + C ) ) = ( A + ( C + B ) )&(\mbox{\ref{eq:oveq2i}},1)\notag
  \\ 3 &&  & \vdash A \in
\mathbb{C}&\text{Hyp~addcom3acb.1}\notag%bl.addcom3acb.1
  \\ 4 &&  & \vdash B \in
\mathbb{C}&\text{Hyp~addcom3acb.2}\notag%bl.addcom3acb.2
  \\ 5 &&  & \vdash C \in
\mathbb{C}&\text{Hyp~addcom3acb.3}\notag%bl.addcom3acb.3
  \\ 6 &&  & \vdash +( A , B , C ) = ( A + ( B + C ) )&(\mbox{\ref{eq:bl.add3i}},3,4,5)\notag
  \\ 7 &&  & \vdash A \in
\mathbb{C}&\text{Hyp~addcom3acb.1}\notag%bl.addcom3acb.1
  \\ 8 &&  & \vdash C \in
\mathbb{C}&\text{Hyp~addcom3acb.3}\notag%bl.addcom3acb.3
  \\ 9 &&  & \vdash B \in
\mathbb{C}&\text{Hyp~addcom3acb.2}\notag%bl.addcom3acb.2
  \\ 10 &&  & \vdash +( A , C , B ) = ( A + ( C + B ) )&(\mbox{\ref{eq:bl.add3i}},7,8,9)\notag
  \\ 11 &&  & \vdash +( A , B , C ) = +( A , C , B )&(\mbox{\ref{eq:3eqtr4i}},2,6,10)\notag
\end{align}
\end{proof}

\begin{lemma}\blTitle{bl.addcom3ac} Addition of Three Values Commutes
\begin{align}
\vdash A \in \mathbb{C}\quad\&\quad \vdash B \in \mathbb{C}\quad\&\quad \vdash
    C \in \mathbb{C}\quad\Rightarrow\quad \vdash +( A , B , C ) = +( C , B , A
    )\label{eq:bl.addcom3ac}\tag{bl.addcom3ac}
\end{align}
\end{lemma}

\begin{proof}
\begin{align}
  1 &&  & \vdash ( B + C ) = ( C + B )&(\mbox{\ref{eq:addcomgi}})\notag
  \\ 2 &&  & \vdash ( A + ( B + C ) ) = ( A + ( C + B ) )&(\mbox{\ref{eq:oveq2i}},1)\notag
  \\ 3 &&  & \vdash A \in
\mathbb{C}&\text{Hyp~addcom3cba.1}\notag%bl.addcom3cba.1
  \\ 4 &&  & \vdash C \in
\mathbb{C}&\text{Hyp~addcom3cba.3}\notag%bl.addcom3cba.3
  \\ 5 &&  & \vdash B \in
\mathbb{C}&\text{Hyp~addcom3cba.2}\notag%bl.addcom3cba.2
  \\ 6 &&  & \vdash ( A \in \mathbb{C} \wedge C \in \mathbb{C} \wedge B \in \mathbb{C} )&(\mbox{\ref{eq:3pm3.2i}},3,4,5)\notag
  \\ 7 &&  & \vdash ( ( A \in \mathbb{C} \wedge C \in \mathbb{C} \wedge B \in \mathbb{C} ) \rightarrow ( A + ( C + B ) ) = ( C + ( A + B ) ) )&(\mbox{\ref{eq:add12}})\notag
  \\ 8 &&  & \vdash ( A + ( C + B ) ) = ( C + ( A + B ) )&(\mbox{\ref{eq:ax-mp}},6,7)\notag
  \\ 9 &&  & \vdash ( A + ( B + C ) ) = ( C + ( A + B ) )&(\mbox{\ref{eq:eqtri}},2,8)\notag
  \\ 10 &&  & \vdash ( A + B ) = ( B + A )&(\mbox{\ref{eq:addcomgi}})\notag
  \\ 11 &&  & \vdash ( C + ( A + B ) ) = ( C + ( B + A ) )&(\mbox{\ref{eq:oveq2i}},10)\notag
  \\ 12 &&  & \vdash ( A + ( B + C ) ) = ( C + ( B + A ) )&(\mbox{\ref{eq:eqtri}},9,11)\notag
  \\ 13 &&  & \vdash A \in
\mathbb{C}&\text{Hyp~addcom3cba.1}\notag%bl.addcom3cba.1
  \\ 14 &&  & \vdash B \in
\mathbb{C}&\text{Hyp~addcom3cba.2}\notag%bl.addcom3cba.2
  \\ 15 &&  & \vdash C \in
\mathbb{C}&\text{Hyp~addcom3cba.3}\notag%bl.addcom3cba.3
  \\ 16 &&  & \vdash +( A , B , C ) = ( A + ( B + C ) )&(\mbox{\ref{eq:bl.add3i}},13,14,15)\notag
  \\ 17 &&  & \vdash C \in
\mathbb{C}&\text{Hyp~addcom3cba.3}\notag%bl.addcom3cba.3
  \\ 18 &&  & \vdash B \in
\mathbb{C}&\text{Hyp~addcom3cba.2}\notag%bl.addcom3cba.2
  \\ 19 &&  & \vdash A \in
\mathbb{C}&\text{Hyp~addcom3cba.1}\notag%bl.addcom3cba.1
  \\ 20 &&  & \vdash +( C , B , A ) = ( C + ( B + A ) )&(\mbox{\ref{eq:bl.add3i}},17,18,19)\notag
  \\ 21 &&  & \vdash +( A , B , C ) = +( C , B , A )&(\mbox{\ref{eq:3eqtr4i}},12,16,20)\notag
\end{align}
\end{proof}

\begin{lemma}\blTitle{bl.addcom3cab} Addition of Three Values Commutes
\begin{align}
\vdash A \in \mathbb{C}\quad\&\quad \vdash B \in \mathbb{C}\quad\&\quad \vdash
    C \in \mathbb{C}\quad\Rightarrow\quad \vdash +( A , B , C ) = +( C , A , B
    )\label{eq:bl.addcom3cab}\tag{bl.addcom3cab}
\end{align}
\end{lemma}

\begin{proof}
\begin{align}
  1 &&  & \vdash ( B + C ) = ( C + B )&(\mbox{\ref{eq:addcomgi}})\notag
  \\ 2 &&  & \vdash ( A + ( B + C ) ) = ( A + ( C + B ) )&(\mbox{\ref{eq:oveq2i}},1)\notag
  \\ 3 &&  & \vdash A \in
\mathbb{C}&\text{Hyp~addcom3cab.1}\notag%bl.addcom3cab.1
  \\ 4 &&  & \vdash C \in
\mathbb{C}&\text{Hyp~addcom3cab.3}\notag%bl.addcom3cab.3
  \\ 5 &&  & \vdash B \in
\mathbb{C}&\text{Hyp~addcom3cab.2}\notag%bl.addcom3cab.2
  \\ 6 &&  & \vdash ( A \in \mathbb{C} \wedge C \in \mathbb{C} \wedge B \in \mathbb{C} )&(\mbox{\ref{eq:3pm3.2i}},3,4,5)\notag
  \\ 7 &&  & \vdash ( ( A \in \mathbb{C} \wedge C \in \mathbb{C} \wedge B \in \mathbb{C} ) \rightarrow ( A + ( C + B ) ) = ( C + ( A + B ) ) )&(\mbox{\ref{eq:add12}})\notag
  \\ 8 &&  & \vdash ( A + ( C + B ) ) = ( C + ( A + B ) )&(\mbox{\ref{eq:ax-mp}},6,7)\notag
  \\ 9 &&  & \vdash ( A + ( B + C ) ) = ( C + ( A + B ) )&(\mbox{\ref{eq:eqtri}},2,8)\notag
  \\ 10 &&  & \vdash A \in
\mathbb{C}&\text{Hyp~addcom3cab.1}\notag%bl.addcom3cab.1
  \\ 11 &&  & \vdash B \in
\mathbb{C}&\text{Hyp~addcom3cab.2}\notag%bl.addcom3cab.2
  \\ 12 &&  & \vdash C \in
\mathbb{C}&\text{Hyp~addcom3cab.3}\notag%bl.addcom3cab.3
  \\ 13 &&  & \vdash +( A , B , C ) = ( A + ( B + C ) )&(\mbox{\ref{eq:bl.add3i}},10,11,12)\notag
  \\ 14 &&  & \vdash C \in
\mathbb{C}&\text{Hyp~addcom3cab.3}\notag%bl.addcom3cab.3
  \\ 15 &&  & \vdash A \in
\mathbb{C}&\text{Hyp~addcom3cab.1}\notag%bl.addcom3cab.1
  \\ 16 &&  & \vdash B \in
\mathbb{C}&\text{Hyp~addcom3cab.2}\notag%bl.addcom3cab.2
  \\ 17 &&  & \vdash +( C , A , B ) = ( C + ( A + B ) )&(\mbox{\ref{eq:bl.add3i}},14,15,16)\notag
  \\ 18 &&  & \vdash +( A , B , C ) = +( C , A , B )&(\mbox{\ref{eq:3eqtr4i}},9,13,17)\notag
\end{align}
\end{proof}

\begin{metadefinition}\blTitle{bl.cmul} Multi-term Multiplication is a Class
\begin{align}
{\rm class} \cdot( \msc{csl_1} )\label{eq:bl.cmul}\tag{bl.cmul}
\end{align}
\end{metadefinition}

\begin{metadefinition}\blTitle{df-bl.mul1} Multi-term Multiplication of Single
Value is that Value.
\begin{align}
\vdash ( A \in \mathbb{C} \rightarrow \cdot( A ) = A
    )\label{eq:df-bl.mul1}\tag{df-bl.mul1}
\end{align}
\end{metadefinition}

\begin{metadefinition}\blTitle{df-bl.mul2} Multi-term Multiplication of Two
Class Lists is their Product.
\begin{align}
\vdash \cdot( \msc{csl_1} , \msc{csl_2} ) = ( \cdot( \msc{csl_1} ) \cdot \cdot(
    \msc{csl_2} ) )\label{eq:df-bl.mul2}\tag{df-bl.mul2}
\end{align}
\end{metadefinition}

\begin{lemma}\blTitle{bl.mulcom} Multi-term Multiplication Commutes
\begin{align}
\vdash \cdot( \msc{csl_1} ) \in \mathbb{C}\quad\&\quad \vdash \cdot(
    \msc{csl_2} ) \in \mathbb{C}\quad\Rightarrow\quad \vdash \cdot( \msc{csl_1}
    , \msc{csl_2} ) = \cdot( \msc{csl_2} , \msc{csl_1}
    )\label{eq:bl.mulcom}\tag{bl.mulcom}
\end{align}
\end{lemma}

\begin{proof}
\begin{align}
  1 &&  & \vdash \cdot( \msc{csl_1} ) \in \mathbb{C}&\text{Hyp~mulcom.1}\notag%bl.mulcom.1
  \\ 2 &&  & \vdash \cdot( \msc{csl_2} ) \in \mathbb{C}&\text{Hyp~mulcom.2}\notag%bl.mulcom.2
  \\ 3 &&  & \vdash ( \cdot( \msc{csl_1} ) \cdot \cdot( \msc{csl_2} ) ) = ( \cdot( \msc{csl_2} ) \cdot \cdot( \msc{csl_1} ) )&(\mbox{\ref{eq:mulcomi}},1,2)\notag
  \\ 4 &&  & \vdash \cdot( \msc{csl_1} , \msc{csl_2} ) = ( \cdot( \msc{csl_1} ) \cdot \cdot( \msc{csl_2} ) )&(\mbox{\ref{eq:df-bl.mul2}})\notag
  \\ 5 &&  & \vdash \cdot( \msc{csl_2} , \msc{csl_1} ) = ( \cdot( \msc{csl_2} ) \cdot \cdot( \msc{csl_1} ) )&(\mbox{\ref{eq:df-bl.mul2}})\notag
  \\ 6 &&  & \vdash \cdot( \msc{csl_1} , \msc{csl_2} ) = \cdot( \msc{csl_2} , \msc{csl_1} )&(\mbox{\ref{eq:3eqtr4i}},3,4,5)\notag
\end{align}
\end{proof}

\begin{lemma}\blTitle{bl.muleqi} Multiplying Equals
\begin{align}
\vdash A = B\quad\&\quad \vdash C = D\quad\Rightarrow\quad \vdash ( A \cdot C )
    = ( B \cdot D )\label{eq:bl.muleqi}\tag{bl.muleqi}
\end{align}
\end{lemma}

\begin{proof}
\begin{align}
  1 &&  & \vdash A = B&\text{Hyp~muleqi.1}\notag%bl.muleqi.1
  \\ 2 &&  & \vdash C = D&\text{Hyp~muleqi.2}\notag%bl.muleqi.2
  \\ 3 &&  & \vdash ( A \cdot C ) = ( B \cdot D )&(\mbox{\ref{eq:oveq12i}},1,2)\notag
\end{align}
\end{proof}

\begin{lemma}\blTitle{bl.mul2i} Multi-term Multiplication of Two Values is
their Product.
\begin{align}
\vdash A \in \mathbb{C}\quad\&\quad \vdash B \in
    \mathbb{C}\quad\Rightarrow\quad \vdash \cdot( A , B ) = ( A \cdot B
    )\label{eq:bl.mul2i}\tag{bl.mul2i}
\end{align}
\end{lemma}

\begin{proof}
\begin{align}
  1 &&  & \vdash \cdot( A , B ) = ( \cdot( A ) \cdot \cdot( B ) )&(\mbox{\ref{eq:df-bl.mul2}})\notag
  \\ 2 &&  & \vdash A \in \mathbb{C}&\text{Hyp~mul2i.1}\notag%bl.mul2i.1
  \\ 3 &&  & \vdash ( A \in \mathbb{C} \rightarrow \cdot( A ) = A )&(\mbox{\ref{eq:df-bl.mul1}})\notag
  \\ 4 &&  & \vdash \cdot( A ) = A&(\mbox{\ref{eq:ax-mp}},2,3)\notag
  \\ 5 &&  & \vdash B \in \mathbb{C}&\text{Hyp~mul2i.2}\notag%bl.mul2i.2
  \\ 6 &&  & \vdash ( B \in \mathbb{C} \rightarrow \cdot( B ) = B )&(\mbox{\ref{eq:df-bl.mul1}})\notag
  \\ 7 &&  & \vdash \cdot( B ) = B&(\mbox{\ref{eq:ax-mp}},5,6)\notag
  \\ 8 &&  & \vdash ( \cdot( A ) \cdot \cdot( B ) ) = ( A \cdot B )&(\mbox{\ref{eq:bl.muleqi}},4,7)\notag
  \\ 9 &&  & \vdash \cdot( A , B ) = ( A \cdot B )&(\mbox{\ref{eq:eqtri}},1,8)\notag
\end{align}
\end{proof}

\begin{lemma}\blTitle{bl.mul3i} Multi-term Multiplication of Three Values
\begin{align}
\vdash A \in \mathbb{C}\quad\&\quad \vdash B \in \mathbb{C}\quad\&\quad \vdash
    C \in \mathbb{C}\quad\Rightarrow\quad \vdash \cdot( A , B , C ) = ( A \cdot
    ( B \cdot C ) )\label{eq:bl.mul3i}\tag{bl.mul3i}
\end{align}
\end{lemma}

\begin{proof}
\begin{align}
  1 &&  & \vdash \cdot( A , B , C ) = ( \cdot( A ) \cdot \cdot( B , C ) )&(\mbox{\ref{eq:df-bl.mul2}})\notag
  \\ 2 &&  & \vdash A \in \mathbb{C}&\text{Hyp~mul3i.1}\notag%bl.mul3i.1
  \\ 3 &&  & \vdash ( A \in \mathbb{C} \rightarrow \cdot( A ) = A )&(\mbox{\ref{eq:df-bl.mul1}})\notag
  \\ 4 &&  & \vdash \cdot( A ) = A&(\mbox{\ref{eq:ax-mp}},2,3)\notag
  \\ 5 &&  & \vdash B \in \mathbb{C}&\text{Hyp~mul3i.2}\notag%bl.mul3i.2
  \\ 6 &&  & \vdash C \in \mathbb{C}&\text{Hyp~mul3i.3}\notag%bl.mul3i.3
  \\ 7 &&  & \vdash \cdot( B , C ) = ( B \cdot C )&(\mbox{\ref{eq:bl.mul2i}},5,6)\notag
  \\ 8 &&  & \vdash ( \cdot( A ) \cdot \cdot( B , C ) ) = ( A \cdot ( B \cdot C ) )&(\mbox{\ref{eq:bl.muleqi}},4,7)\notag
  \\ 9 &&  & \vdash \cdot( A , B , C ) = ( A \cdot ( B \cdot C ) )&(\mbox{\ref{eq:eqtri}},1,8)\notag
\end{align}
\end{proof}

\begin{lemma}\blTitle{bl.mul4i} Multi-term Multiplication of Four Values
\begin{align}
\vdash A \in \mathbb{C}\quad\&\quad \vdash B \in \mathbb{C}\quad\&\quad \vdash
    C \in \mathbb{C}\quad\&\quad \vdash D \in \mathbb{C}\quad\Rightarrow\quad
    \vdash \cdot( A , B , C , D ) = ( ( A \cdot B ) \cdot ( C \cdot D )
    )\label{eq:bl.mul4i}\tag{bl.mul4i}
\end{align}
\end{lemma}

\begin{proof}
\begin{align}
  1 &&  & \vdash \cdot( A , B , C , D ) = ( \cdot( A , B ) \cdot \cdot( C , D ) )&(\mbox{\ref{eq:df-bl.mul2}})\notag
  \\ 2 &&  & \vdash A \in \mathbb{C}&\text{Hyp~mul4i.1}\notag%bl.mul4i.1
  \\ 3 &&  & \vdash B \in \mathbb{C}&\text{Hyp~mul4i.2}\notag%bl.mul4i.2
  \\ 4 &&  & \vdash \cdot( A , B ) = ( A \cdot B )&(\mbox{\ref{eq:bl.mul2i}},2,3)\notag
  \\ 5 &&  & \vdash C \in \mathbb{C}&\text{Hyp~mul4i.3}\notag%bl.mul4i.3
  \\ 6 &&  & \vdash D \in \mathbb{C}&\text{Hyp~mul4i.4}\notag%bl.mul4i.4
  \\ 7 &&  & \vdash \cdot( C , D ) = ( C \cdot D )&(\mbox{\ref{eq:bl.mul2i}},5,6)\notag
  \\ 8 &&  & \vdash ( \cdot( A , B ) \cdot \cdot( C , D ) ) = ( ( A \cdot B ) \cdot ( C \cdot D ) )&(\mbox{\ref{eq:bl.muleqi}},4,7)\notag
  \\ 9 &&  & \vdash \cdot( A , B , C , D ) = ( ( A \cdot B ) \cdot ( C \cdot D ) )&(\mbox{\ref{eq:eqtri}},1,8)\notag
\end{align}
\end{proof}

\begin{lemma}\blTitle{bl.mulcom2} Multiplication of Two Values Commutes
\begin{align}
\vdash A \in \mathbb{C}\quad\&\quad \vdash B \in
    \mathbb{C}\quad\Rightarrow\quad \vdash \cdot( A , B ) = \cdot( B , A
    )\label{eq:bl.mulcom2}\tag{bl.mulcom2}
\end{align}
\end{lemma}

\begin{proof}
\begin{align}
  1 &&  & \vdash A \in \mathbb{C}&\text{Hyp~mulcom2.1}\notag%bl.mulcom2.1
  \\ 2 &&  & \vdash ( A \in \mathbb{C} \rightarrow \cdot( A ) = A )&(\mbox{\ref{eq:df-bl.mul1}})\notag
  \\ 3 &&  & \vdash \cdot( A ) = A&(\mbox{\ref{eq:ax-mp}},1,2)\notag
  \\ 4 &&  & \vdash A \in \mathbb{C}&\text{Hyp~mulcom2.1}\notag%bl.mulcom2.1
  \\ 5 &&  & \vdash \cdot( A ) \in \mathbb{C}&(\mbox{\ref{eq:eqeltri}},3,4)\notag
  \\ 6 &&  & \vdash B \in \mathbb{C}&\text{Hyp~mulcom2.2}\notag%bl.mulcom2.2
  \\ 7 &&  & \vdash ( B \in \mathbb{C} \rightarrow \cdot( B ) = B )&(\mbox{\ref{eq:df-bl.mul1}})\notag
  \\ 8 &&  & \vdash \cdot( B ) = B&(\mbox{\ref{eq:ax-mp}},6,7)\notag
  \\ 9 &&  & \vdash B \in \mathbb{C}&\text{Hyp~mulcom2.2}\notag%bl.mulcom2.2
  \\ 10 &&  & \vdash \cdot( B ) \in \mathbb{C}&(\mbox{\ref{eq:eqeltri}},8,9)\notag
  \\ 11 &&  & \vdash ( \cdot( A ) \cdot \cdot( B ) ) = ( \cdot( B ) \cdot \cdot( A ) )&(\mbox{\ref{eq:mulcomi}},5,10)\notag
  \\ 12 &&  & \vdash \cdot( A , B ) = ( \cdot( A ) \cdot \cdot( B ) )&(\mbox{\ref{eq:df-bl.mul2}})\notag
  \\ 13 &&  & \vdash \cdot( B , A ) = ( \cdot( B ) \cdot \cdot( A ) )&(\mbox{\ref{eq:df-bl.mul2}})\notag
  \\ 14 &&  & \vdash \cdot( A , B ) = \cdot( B , A )&(\mbox{\ref{eq:3eqtr4i}},11,12,13)\notag
\end{align}
\end{proof}

\begin{lemma}\blTitle{bl.mulcom3bac} Multiplication of Three Values
Commutes
\begin{align}
\vdash A \in \mathbb{C}\quad\&\quad \vdash B \in \mathbb{C}\quad\&\quad \vdash
    C \in \mathbb{C}\quad\Rightarrow\quad \vdash \cdot( A , B , C ) = \cdot( B
    , A , C )\label{eq:bl.mulcom3bac}\tag{bl.mulcom3bac}
\end{align}
\end{lemma}

\begin{proof}
\begin{align}
  1 &&  & \vdash A \in \mathbb{C}&\text{Hyp~mulcom3bac.1}\notag%bl.mulcom3bac.1
  \\ 2 &&  & \vdash B \in
\mathbb{C}&\text{Hyp~mulcom3bac.2}\notag%bl.mulcom3bac.2
  \\ 3 &&  & \vdash C \in
\mathbb{C}&\text{Hyp~mulcom3bac.3}\notag%bl.mulcom3bac.3
  \\ 4 &&  & \vdash ( A \cdot ( B \cdot C ) ) = ( B \cdot ( A \cdot C ) )&(\mbox{\ref{eq:mul12i}},1,2,3)\notag
  \\ 5 &&  & \vdash A \in
\mathbb{C}&\text{Hyp~mulcom3bac.1}\notag%bl.mulcom3bac.1
  \\ 6 &&  & \vdash B \in
\mathbb{C}&\text{Hyp~mulcom3bac.2}\notag%bl.mulcom3bac.2
  \\ 7 &&  & \vdash C \in
\mathbb{C}&\text{Hyp~mulcom3bac.3}\notag%bl.mulcom3bac.3
  \\ 8 &&  & \vdash \cdot( A , B , C ) = ( A \cdot ( B \cdot C ) )&(\mbox{\ref{eq:bl.mul3i}},5,6,7)\notag
  \\ 9 &&  & \vdash B \in
\mathbb{C}&\text{Hyp~mulcom3bac.2}\notag%bl.mulcom3bac.2
  \\ 10 &&  & \vdash A \in
\mathbb{C}&\text{Hyp~mulcom3bac.1}\notag%bl.mulcom3bac.1
  \\ 11 &&  & \vdash C \in
\mathbb{C}&\text{Hyp~mulcom3bac.3}\notag%bl.mulcom3bac.3
  \\ 12 &&  & \vdash \cdot( B , A , C ) = ( B \cdot ( A \cdot C ) )&(\mbox{\ref{eq:bl.mul3i}},9,10,11)\notag
  \\ 13 &&  & \vdash \cdot( A , B , C ) = \cdot( B , A , C )&(\mbox{\ref{eq:3eqtr4i}},4,8,12)\notag
\end{align}
\end{proof}

\begin{lemma}\blTitle{bl.mulcom3acb} Multiplication of Three Values
Commutes
\begin{align}
\vdash A \in \mathbb{C}\quad\&\quad \vdash B \in \mathbb{C}\quad\&\quad \vdash
    C \in \mathbb{C}\quad\Rightarrow\quad \vdash \cdot( A , B , C ) = \cdot( A
    , C , B )\label{eq:bl.mulcom3acb}\tag{bl.mulcom3acb}
\end{align}
\end{lemma}

\begin{proof}
\begin{align}
  1 &&  & \vdash B \in \mathbb{C}&\text{Hyp~mulcom3acb.2}\notag%bl.mulcom3acb.2
  \\ 2 &&  & \vdash C \in
\mathbb{C}&\text{Hyp~mulcom3acb.3}\notag%bl.mulcom3acb.3
  \\ 3 &&  & \vdash ( B \cdot C ) = ( C \cdot B )&(\mbox{\ref{eq:mulcomi}},1,2)\notag
  \\ 4 &&  & \vdash ( A \cdot ( B \cdot C ) ) = ( A \cdot ( C \cdot B ) )&(\mbox{\ref{eq:oveq2i}},3)\notag
  \\ 5 &&  & \vdash A \in
\mathbb{C}&\text{Hyp~mulcom3acb.1}\notag%bl.mulcom3acb.1
  \\ 6 &&  & \vdash B \in
\mathbb{C}&\text{Hyp~mulcom3acb.2}\notag%bl.mulcom3acb.2
  \\ 7 &&  & \vdash C \in
\mathbb{C}&\text{Hyp~mulcom3acb.3}\notag%bl.mulcom3acb.3
  \\ 8 &&  & \vdash \cdot( A , B , C ) = ( A \cdot ( B \cdot C ) )&(\mbox{\ref{eq:bl.mul3i}},5,6,7)\notag
  \\ 9 &&  & \vdash A \in
\mathbb{C}&\text{Hyp~mulcom3acb.1}\notag%bl.mulcom3acb.1
  \\ 10 &&  & \vdash C \in
\mathbb{C}&\text{Hyp~mulcom3acb.3}\notag%bl.mulcom3acb.3
  \\ 11 &&  & \vdash B \in
\mathbb{C}&\text{Hyp~mulcom3acb.2}\notag%bl.mulcom3acb.2
  \\ 12 &&  & \vdash \cdot( A , C , B ) = ( A \cdot ( C \cdot B ) )&(\mbox{\ref{eq:bl.mul3i}},9,10,11)\notag
  \\ 13 &&  & \vdash \cdot( A , B , C ) = \cdot( A , C , B )&(\mbox{\ref{eq:3eqtr4i}},4,8,12)\notag
\end{align}
\end{proof}

\begin{lemma}\blTitle{bl.mulcom3cba} Multiplication of Three Values
Commutes
\begin{align}
\vdash A \in \mathbb{C}\quad\&\quad \vdash B \in \mathbb{C}\quad\&\quad \vdash
    C \in \mathbb{C}\quad\Rightarrow\quad \vdash \cdot( A , B , C ) = \cdot( C
    , B , A )\label{eq:bl.mulcom3cba}\tag{bl.mulcom3cba}
\end{align}
\end{lemma}

\begin{proof}
\begin{align}
  1 &&  & \vdash A \in \mathbb{C}&\text{Hyp~mulcom3cba.1}\notag%bl.mulcom3cba.1
  \\ 2 &&  & \vdash B \in
\mathbb{C}&\text{Hyp~mulcom3cba.2}\notag%bl.mulcom3cba.2
  \\ 3 &&  & \vdash C \in
\mathbb{C}&\text{Hyp~mulcom3cba.3}\notag%bl.mulcom3cba.3
  \\ 4 &&  & \vdash ( B \in \mathbb{C} \wedge C \in \mathbb{C} )&(\mbox{\ref{eq:pm3.2i}},2,3)\notag
  \\ 5 &&  & \vdash ( ( B \in \mathbb{C} \wedge C \in \mathbb{C} ) \rightarrow ( B \cdot C ) \in \mathbb{C} )&(\mbox{\ref{eq:axmulcl}})\notag
  \\ 6 &&  & \vdash ( B \cdot C ) \in \mathbb{C}&(\mbox{\ref{eq:ax-mp}},4,5)\notag
  \\ 7 &&  & \vdash ( A \in \mathbb{C} \wedge ( B \cdot C ) \in \mathbb{C} )&(\mbox{\ref{eq:pm3.2i}},1,6)\notag
  \\ 8 &&  & \vdash ( ( A \in \mathbb{C} \wedge ( B \cdot C ) \in \mathbb{C} ) \rightarrow ( A \cdot ( B \cdot C ) ) = ( ( B \cdot C ) \cdot A ) )&(\mbox{\ref{eq:axmulcom}})\notag
  \\ 9 &&  & \vdash ( A \cdot ( B \cdot C ) ) = ( ( B \cdot C ) \cdot A )&(\mbox{\ref{eq:ax-mp}},7,8)\notag
  \\ 10 &&  & \vdash C \in
\mathbb{C}&\text{Hyp~mulcom3cba.3}\notag%bl.mulcom3cba.3
  \\ 11 &&  & \vdash B \in
\mathbb{C}&\text{Hyp~mulcom3cba.2}\notag%bl.mulcom3cba.2
  \\ 12 &&  & \vdash ( C \in \mathbb{C} \wedge B \in \mathbb{C} )&(\mbox{\ref{eq:pm3.2i}},10,11)\notag
  \\ 13 &&  & \vdash ( ( C \in \mathbb{C} \wedge B \in \mathbb{C} ) \rightarrow ( C \cdot B ) = ( B \cdot C ) )&(\mbox{\ref{eq:axmulcom}})\notag
  \\ 14 &&  & \vdash ( C \cdot B ) = ( B \cdot C )&(\mbox{\ref{eq:ax-mp}},12,13)\notag
  \\ 15 &&  & \vdash ( ( C \cdot B ) \cdot A ) = ( ( B \cdot C ) \cdot A )&(\mbox{\ref{eq:oveq1i}},14)\notag
  \\ 16 &&  & \vdash C \in
\mathbb{C}&\text{Hyp~mulcom3cba.3}\notag%bl.mulcom3cba.3
  \\ 17 &&  & \vdash B \in
\mathbb{C}&\text{Hyp~mulcom3cba.2}\notag%bl.mulcom3cba.2
  \\ 18 &&  & \vdash A \in
\mathbb{C}&\text{Hyp~mulcom3cba.1}\notag%bl.mulcom3cba.1
  \\ 19 &&  & \vdash ( C \in \mathbb{C} \wedge B \in \mathbb{C} \wedge A \in \mathbb{C} )&(\mbox{\ref{eq:3pm3.2i}},16,17,18)\notag
  \\ 20 &&  & \vdash ( ( C \in \mathbb{C} \wedge B \in \mathbb{C} \wedge A \in \mathbb{C} ) \rightarrow ( ( C \cdot B ) \cdot A ) = ( C \cdot ( B \cdot A ) ) )&(\mbox{\ref{eq:axmulass}})\notag
  \\ 21 &&  & \vdash ( ( C \cdot B ) \cdot A ) = ( C \cdot ( B \cdot A ) )&(\mbox{\ref{eq:ax-mp}},19,20)\notag
  \\ 22 &&  & \vdash ( A \cdot ( B \cdot C ) ) = ( C \cdot ( B \cdot A ) )&(\mbox{\ref{eq:3eqtr2i}},9,15,21)\notag
  \\ 23 &&  & \vdash A \in
\mathbb{C}&\text{Hyp~mulcom3cba.1}\notag%bl.mulcom3cba.1
  \\ 24 &&  & \vdash B \in
\mathbb{C}&\text{Hyp~mulcom3cba.2}\notag%bl.mulcom3cba.2
  \\ 25 &&  & \vdash C \in
\mathbb{C}&\text{Hyp~mulcom3cba.3}\notag%bl.mulcom3cba.3
  \\ 26 &&  & \vdash \cdot( A , B , C ) = ( A \cdot ( B \cdot C ) )&(\mbox{\ref{eq:bl.mul3i}},23,24,25)\notag
  \\ 27 &&  & \vdash C \in
\mathbb{C}&\text{Hyp~mulcom3cba.3}\notag%bl.mulcom3cba.3
  \\ 28 &&  & \vdash B \in
\mathbb{C}&\text{Hyp~mulcom3cba.2}\notag%bl.mulcom3cba.2
  \\ 29 &&  & \vdash A \in
\mathbb{C}&\text{Hyp~mulcom3cba.1}\notag%bl.mulcom3cba.1
  \\ 30 &&  & \vdash \cdot( C , B , A ) = ( C \cdot ( B \cdot A ) )&(\mbox{\ref{eq:bl.mul3i}},27,28,29)\notag
  \\ 31 &&  & \vdash \cdot( A , B , C ) = \cdot( C , B , A )&(\mbox{\ref{eq:3eqtr4i}},22,26,30)\notag
\end{align}
\end{proof}

\begin{lemma}\blTitle{bl.mulcom3cab} Multiplication of Three Values
Commutes
\begin{align}
\vdash A \in \mathbb{C}\quad\&\quad \vdash B \in \mathbb{C}\quad\&\quad \vdash
    C \in \mathbb{C}\quad\Rightarrow\quad \vdash \cdot( A , B , C ) = \cdot( C
    , A , B )\label{eq:bl.mulcom3cab}\tag{bl.mulcom3cab}
\end{align}
\end{lemma}

\begin{proof}
\begin{align}
  1 &&  & \vdash A \in \mathbb{C}&\text{Hyp~mulcom3cab.1}\notag%bl.mulcom3cab.1
  \\ 2 &&  & \vdash B \in
\mathbb{C}&\text{Hyp~mulcom3cab.2}\notag%bl.mulcom3cab.2
  \\ 3 &&  & \vdash C \in
\mathbb{C}&\text{Hyp~mulcom3cab.3}\notag%bl.mulcom3cab.3
  \\ 4 &&  & \vdash ( A \in \mathbb{C} \wedge B \in \mathbb{C} \wedge C \in \mathbb{C} )&(\mbox{\ref{eq:3pm3.2i}},1,2,3)\notag
  \\ 5 &&  & \vdash ( ( A \in \mathbb{C} \wedge B \in \mathbb{C} \wedge C \in \mathbb{C} ) \rightarrow ( ( A \cdot B ) \cdot C ) = ( A \cdot ( B \cdot C ) ) )&(\mbox{\ref{eq:axmulass}})\notag
  \\ 6 &&  & \vdash ( ( A \cdot B ) \cdot C ) = ( A \cdot ( B \cdot C ) )&(\mbox{\ref{eq:ax-mp}},4,5)\notag
  \\ 7 &&  & \vdash A \in
\mathbb{C}&\text{Hyp~mulcom3cab.1}\notag%bl.mulcom3cab.1
  \\ 8 &&  & \vdash B \in
\mathbb{C}&\text{Hyp~mulcom3cab.2}\notag%bl.mulcom3cab.2
  \\ 9 &&  & \vdash ( A \cdot B ) \in \mathbb{C}&(\mbox{\ref{eq:mulcli}},7,8)\notag
  \\ 10 &&  & \vdash C \in
\mathbb{C}&\text{Hyp~mulcom3cab.3}\notag%bl.mulcom3cab.3
  \\ 11 &&  & \vdash ( ( A \cdot B ) \in \mathbb{C} \wedge C \in \mathbb{C} )&(\mbox{\ref{eq:pm3.2i}},9,10)\notag
  \\ 12 &&  & \vdash ( ( ( A \cdot B ) \in \mathbb{C} \wedge C \in \mathbb{C} ) \rightarrow ( ( A \cdot B ) \cdot C ) = ( C \cdot ( A \cdot B ) ) )&(\mbox{\ref{eq:axmulcom}})\notag
  \\ 13 &&  & \vdash ( ( A \cdot B ) \cdot C ) = ( C \cdot ( A \cdot B ) )&(\mbox{\ref{eq:ax-mp}},11,12)\notag
  \\ 14 &&  & \vdash ( A \cdot ( B \cdot C ) ) = ( C \cdot ( A \cdot B ) )&(\mbox{\ref{eq:eqtr3i}},6,13)\notag
  \\ 15 &&  & \vdash A \in
\mathbb{C}&\text{Hyp~mulcom3cab.1}\notag%bl.mulcom3cab.1
  \\ 16 &&  & \vdash B \in
\mathbb{C}&\text{Hyp~mulcom3cab.2}\notag%bl.mulcom3cab.2
  \\ 17 &&  & \vdash C \in
\mathbb{C}&\text{Hyp~mulcom3cab.3}\notag%bl.mulcom3cab.3
  \\ 18 &&  & \vdash \cdot( A , B , C ) = ( A \cdot ( B \cdot C ) )&(\mbox{\ref{eq:bl.mul3i}},15,16,17)\notag
  \\ 19 &&  & \vdash C \in
\mathbb{C}&\text{Hyp~mulcom3cab.3}\notag%bl.mulcom3cab.3
  \\ 20 &&  & \vdash A \in
\mathbb{C}&\text{Hyp~mulcom3cab.1}\notag%bl.mulcom3cab.1
  \\ 21 &&  & \vdash B \in
\mathbb{C}&\text{Hyp~mulcom3cab.2}\notag%bl.mulcom3cab.2
  \\ 22 &&  & \vdash \cdot( C , A , B ) = ( C \cdot ( A \cdot B ) )&(\mbox{\ref{eq:bl.mul3i}},19,20,21)\notag
  \\ 23 &&  & \vdash \cdot( A , B , C ) = \cdot( C , A , B )&(\mbox{\ref{eq:3eqtr4i}},14,18,22)\notag
\end{align}
\end{proof}

\begin{metadefinition}\blTitle{cnow} TO DO: PUT DESCRIPTION HERE
\begin{align}
{\rm class} \mathsf{\msc{now}}\label{eq:cnow}\tag{cnow}
\end{align}
\end{metadefinition}

\begin{metadefinition}\blTitle{df-bl.time} Define domain of time $\mathbf{T}$
from $t = 0$ to $\mathsf{\msc{now}}$
\begin{align}
\vdash \mathbf{T} = \{ \msc{t}~\vert~( \msc{t} \in \mathbb{R} \wedge 0 \le \msc{t}
    \wedge \msc{t} \le \mathsf{\msc{now}} )
    \}\label{eq:df-bl.time}\tag{df-bl.time}
\end{align}
\end{metadefinition}

\begin{metadefinition}\blTitle{df-bl.before} Temporal Precedence (before)
\begin{align}
\vdash ( \msc{t_1} \prec \msc{t_2} \leftrightarrow \msc{t_1} < \msc{t_2}
    )\label{eq:df-bl.before}\tag{df-bl.before}
\end{align}
\end{metadefinition}

\begin{metadefinition}\blTitle{df-bl.beforeeq} Temporal Precedence (before or
coincident)
\begin{align}
\vdash ( \msc{t_1} \preccurlyeq \msc{t_2} \leftrightarrow ( \msc{t_1} <
    \msc{t_2} \vee \msc{t_1} = \msc{t_2} )
    )\label{eq:df-bl.beforeeq}\tag{df-bl.beforeeq}
\end{align}
\end{metadefinition}

\begin{metadefinition}\blTitle{df-bl.prop0} Value of proposition at time
$\msc{t}$
\begin{align}
\vdash ( \mathfrak{M} , \msc{t} \models p \leftrightarrow \msc{t} \in ( V ` p )
    )\label{eq:df-bl.prop0}\tag{df-bl.prop0}
\end{align}
\end{metadefinition}

\begin{metadefinition}\blTitle{df-bl.not} Logical Complement
\begin{align}
\vdash ( \mathfrak{M} , \msc{t} \models \lnot \mw{\varphi} \leftrightarrow
    \mathfrak{M} , \msc{t} \nmodels \mw{\varphi}
    )\label{eq:df-bl.not}\tag{df-bl.not}
\end{align}
\end{metadefinition}

\begin{metadefinition}\blTitle{clat0} Define $( A @ \msc{t} )$ as a class
\begin{align}
{\rm class} ( A @ \msc{t} )\label{eq:clat0}\tag{clat0}
\end{align}
\end{metadefinition}

\begin{metadefinition}\blTitle{clatb} TO DO: PUT DESCRIPTION HERE
\begin{align}
{\rm class} ( A @ B )\label{eq:clatb}\tag{clatb}
\end{align}
\end{metadefinition}

\begin{metadefinition}\blTitle{df-bl.at} @ specifies time of evaluation
\begin{align}
\vdash ( \mathfrak{M} , \msc{t} \models ( \mw{\varphi} @ \msc{t_1} )
    \rightarrow \mathfrak{M} , \msc{t_1} \models \mw{\varphi}
    )\label{eq:df-bl.at}\tag{df-bl.at}
\end{align}
\end{metadefinition}

\begin{metadefinition}\blTitle{df-bl.at2} Second application of @ is
inconsequential
\begin{align}
\vdash ( \mathfrak{M} , \msc{t} \models ( ( \mw{\varphi} @ \msc{t_1} ) @
    \msc{t_2} ) \rightarrow \mathfrak{M} , \msc{t_1} \models \mw{\varphi}
    )\label{eq:df-bl.at2}\tag{df-bl.at2}
\end{align}
\end{metadefinition}

\begin{metadefinition}\blTitle{cldd} Define $[ \msc{t_1}  ..  \msc{t_2} ]$ to be
a class
\begin{align}
{\rm class} [ \msc{t_1} .. \msc{t_2} ]\label{eq:cldd}\tag{cldd}
\end{align}
\end{metadefinition}

\begin{metadefinition}\blTitle{df-bl.dd} Define closed time interval
\begin{align}
\vdash ( \msc{t} \in [ \msc{t_1} .. \msc{t_2} ] \leftrightarrow ( \msc{t_1}
    \preccurlyeq \msc{t} \wedge \msc{t} \preccurlyeq \msc{t_2} )
    )\label{eq:df-bl.dd}\tag{df-bl.dd}
\end{align}
\end{metadefinition}

\begin{lemma}\blTitle{bl.alldd} Universal quantification over closed time
interval
\begin{align}
\vdash ( \forall \msc{t} ( \msc{t} \in [ \msc{t_1} .. \msc{t_2} ] \wedge
    \mw{\varphi} ) \leftrightarrow \forall \msc{t} ( ( \msc{t_1} \preccurlyeq
    \msc{t} \wedge \msc{t} \preccurlyeq \msc{t_2} ) \wedge \mw{\varphi} )
    )\label{eq:bl.alldd}\tag{bl.alldd}
\end{align}
\end{lemma}

\begin{proof}
\begin{align}
  1 &&  & \vdash ( \msc{t} \in [ \msc{t_1} .. \msc{t_2} ] \leftrightarrow ( \msc{t_1} \preccurlyeq \msc{t} \wedge \msc{t} \preccurlyeq \msc{t_2} ) )&(\mbox{\ref{eq:df-bl.dd}})\notag
  \\ 2 &&  & \vdash ( \mw{\varphi} \leftrightarrow \mw{\varphi} )&(\mbox{\ref{eq:biid}})\notag
  \\ 3 &&  & \vdash ( ( \msc{t} \in [ \msc{t_1} .. \msc{t_2} ] \wedge \mw{\varphi} ) \leftrightarrow ( ( \msc{t_1} \preccurlyeq \msc{t} \wedge \msc{t} \preccurlyeq \msc{t_2} ) \wedge \mw{\varphi} ) )&(\mbox{\ref{eq:anbi12i}},1,2)\notag
  \\ 4 &&  & \vdash ( \forall \msc{t} ( \msc{t} \in [ \msc{t_1} .. \msc{t_2} ] \wedge \mw{\varphi} ) \leftrightarrow \forall \msc{t} ( ( \msc{t_1} \preccurlyeq \msc{t} \wedge \msc{t} \preccurlyeq \msc{t_2} ) \wedge \mw{\varphi} ) )&(\mbox{\ref{eq:albii}},3)\notag
\end{align}
\end{proof}

\begin{metadefinition}\blTitle{clacd} Define $[ \msc{t_1}  ,.  \msc{t_2} ]$ to
be a class
\begin{align}
{\rm class} [ \msc{t_1} ,. \msc{t_2} ]\label{eq:clacd}\tag{clacd}
\end{align}
\end{metadefinition}

\begin{metadefinition}\blTitle{df-bl.cd} Define open-left time interval
\begin{align}
\vdash ( \msc{t} \in [ \msc{t_1} ,. \msc{t_2} ] \leftrightarrow ( \msc{t_1}
    \prec \msc{t} \wedge \msc{t} \preccurlyeq \msc{t_2} )
    )\label{eq:df-bl.cd}\tag{df-bl.cd}
\end{align}
\end{metadefinition}

\begin{lemma}\blTitle{bl.allcd} Universal quantification over open-left
time interval
\begin{align}
\vdash ( \forall \msc{t} ( \msc{t} \in [ \msc{t_1} ,. \msc{t_2} ] \wedge
    \mw{\varphi} ) \leftrightarrow \forall \msc{t} ( ( \msc{t_1} \prec \msc{t}
    \wedge \msc{t} \preccurlyeq \msc{t_2} ) \wedge \mw{\varphi} )
    )\label{eq:bl.allcd}\tag{bl.allcd}
\end{align}
\end{lemma}

\begin{proof}
\begin{align}
  1 &&  & \vdash ( \msc{t} \in [ \msc{t_1} ,. \msc{t_2} ] \leftrightarrow ( \msc{t_1} \prec \msc{t} \wedge \msc{t} \preccurlyeq \msc{t_2} ) )&(\mbox{\ref{eq:df-bl.cd}})\notag
  \\ 2 &&  & \vdash ( \mw{\varphi} \leftrightarrow \mw{\varphi} )&(\mbox{\ref{eq:biid}})\notag
  \\ 3 &&  & \vdash ( ( \msc{t} \in [ \msc{t_1} ,. \msc{t_2} ] \wedge \mw{\varphi} ) \leftrightarrow ( ( \msc{t_1} \prec \msc{t} \wedge \msc{t} \preccurlyeq \msc{t_2} ) \wedge \mw{\varphi} ) )&(\mbox{\ref{eq:anbi12i}},1,2)\notag
  \\ 4 &&  & \vdash ( \forall \msc{t} ( \msc{t} \in [ \msc{t_1} ,. \msc{t_2} ] \wedge \mw{\varphi} ) \leftrightarrow \forall \msc{t} ( ( \msc{t_1} \prec \msc{t} \wedge \msc{t} \preccurlyeq \msc{t_2} ) \wedge \mw{\varphi} ) )&(\mbox{\ref{eq:albii}},3)\notag
\end{align}
\end{proof}

\begin{metadefinition}\blTitle{cldc} Define $[ \msc{t_1}  .,  \msc{t_2} ]$ to be
a class
\begin{align}
{\rm class} [ \msc{t_1} ., \msc{t_2} ]\label{eq:cldc}\tag{cldc}
\end{align}
\end{metadefinition}

\begin{metadefinition}\blTitle{df-bl.dc} Define open-right time interval
\begin{align}
\vdash ( \msc{t} \in [ \msc{t_1} ., \msc{t_2} ] \leftrightarrow ( \msc{t_1}
    \preccurlyeq \msc{t} \wedge \msc{t} \prec \msc{t_2} )
    )\label{eq:df-bl.dc}\tag{df-bl.dc}
\end{align}
\end{metadefinition}

\begin{lemma}\blTitle{bl.alldc} Universal quantification over open-right
time interval
\begin{align}
\vdash ( \forall \msc{t} ( \msc{t} \in [ \msc{t_1} ., \msc{t_2} ] \wedge
    \mw{\varphi} ) \leftrightarrow \forall \msc{t} ( ( \msc{t_1} \preccurlyeq
    \msc{t} \wedge \msc{t} \prec \msc{t_2} ) \wedge \mw{\varphi} )
    )\label{eq:bl.alldc}\tag{bl.alldc}
\end{align}
\end{lemma}

\begin{proof}
\begin{align}
  1 &&  & \vdash ( \msc{t} \in [ \msc{t_1} ., \msc{t_2} ] \leftrightarrow ( \msc{t_1} \preccurlyeq \msc{t} \wedge \msc{t} \prec \msc{t_2} ) )&(\mbox{\ref{eq:df-bl.dc}})\notag
  \\ 2 &&  & \vdash ( \mw{\varphi} \leftrightarrow \mw{\varphi} )&(\mbox{\ref{eq:biid}})\notag
  \\ 3 &&  & \vdash ( ( \msc{t} \in [ \msc{t_1} ., \msc{t_2} ] \wedge \mw{\varphi} ) \leftrightarrow ( ( \msc{t_1} \preccurlyeq \msc{t} \wedge \msc{t} \prec \msc{t_2} ) \wedge \mw{\varphi} ) )&(\mbox{\ref{eq:anbi12i}},1,2)\notag
  \\ 4 &&  & \vdash ( \forall \msc{t} ( \msc{t} \in [ \msc{t_1} ., \msc{t_2} ] \wedge \mw{\varphi} ) \leftrightarrow \forall \msc{t} ( ( \msc{t_1} \preccurlyeq \msc{t} \wedge \msc{t} \prec \msc{t_2} ) \wedge \mw{\varphi} ) )&(\mbox{\ref{eq:albii}},3)\notag
\end{align}
\end{proof}

\begin{metadefinition}\blTitle{clcc} Define $[ \msc{t_1} \,\mathrm{{\tt -}c}
\msc{t_2} ]$ to be a class
\begin{align}
{\rm class} [ \msc{t_1} ,, \msc{t_2} ]\label{eq:clcc}\tag{clcc}
\end{align}
\end{metadefinition}

\begin{metadefinition}\blTitle{df-bl.cc} Define open time interval
\begin{align}
\vdash ( \msc{t} \in [ \msc{t_1} ,, \msc{t_2} ] \leftrightarrow ( \msc{t_1}
    \prec \msc{t} \wedge \msc{t} \prec \msc{t_2} )
    )\label{eq:df-bl.cc}\tag{df-bl.cc}
\end{align}
\end{metadefinition}

\begin{lemma}\blTitle{bl.allcc} Universal quantification over open time
interval
\begin{align}
\vdash ( \forall \msc{t} ( \msc{t} \in [ \msc{t_1} ,, \msc{t_2} ] \wedge
    \mw{\varphi} ) \leftrightarrow \forall \msc{t} ( ( \msc{t_1} \prec \msc{t}
    \wedge \msc{t} \prec \msc{t_2} ) \wedge \mw{\varphi} )
    )\label{eq:bl.allcc}\tag{bl.allcc}
\end{align}
\end{lemma}

\begin{proof}
\begin{align}
  1 &&  & \vdash ( \msc{t} \in [ \msc{t_1} ,, \msc{t_2} ] \leftrightarrow ( \msc{t_1} \prec \msc{t} \wedge \msc{t} \prec \msc{t_2} ) )&(\mbox{\ref{eq:df-bl.cc}})\notag
  \\ 2 &&  & \vdash ( \mw{\varphi} \leftrightarrow \mw{\varphi} )&(\mbox{\ref{eq:biid}})\notag
  \\ 3 &&  & \vdash ( ( \msc{t} \in [ \msc{t_1} ,, \msc{t_2} ] \wedge \mw{\varphi} ) \leftrightarrow ( ( \msc{t_1} \prec \msc{t} \wedge \msc{t} \prec \msc{t_2} ) \wedge \mw{\varphi} ) )&(\mbox{\ref{eq:anbi12i}},1,2)\notag
  \\ 4 &&  & \vdash ( \forall \msc{t} ( \msc{t} \in [ \msc{t_1} ,, \msc{t_2} ] \wedge \mw{\varphi} ) \leftrightarrow \forall \msc{t} ( ( \msc{t_1} \prec \msc{t} \wedge \msc{t} \prec \msc{t_2} ) \wedge \mw{\varphi} ) )&(\mbox{\ref{eq:albii}},3)\notag
\end{align}
\end{proof}

\begin{lemma}\blTitle{bl.exdd} Existential quantification over closed
interval
\begin{align}
\vdash ( \exists \msc{t} ( \msc{t} \in [ \msc{t_1} .. \msc{t_2} ] \wedge
    \mw{\varphi} ) \leftrightarrow \lnot \forall \msc{t} \lnot ( \msc{t} \in [
    \msc{t_1} .. \msc{t_2} ] \wedge \mw{\varphi} )
    )\label{eq:bl.exdd}\tag{bl.exdd}
\end{align}
\end{lemma}

\begin{proof}
\begin{align}
  1 &&  & \vdash ( \exists \msc{t} ( \msc{t} \in [ \msc{t_1} .. \msc{t_2} ] \wedge \mw{\varphi} ) \leftrightarrow \lnot \forall \msc{t} \lnot ( \msc{t} \in [ \msc{t_1} .. \msc{t_2} ] \wedge \mw{\varphi} ) )&(\mbox{\ref{eq:df-ex}})\notag
\end{align}
\end{proof}

\begin{lemma}\blTitle{bl.excc} Existential quantification over open
interval
\begin{align}
\vdash ( \exists \msc{t} ( \msc{t} \in [ \msc{t_1} ,, \msc{t_2} ] \wedge
    \mw{\varphi} ) \leftrightarrow \lnot \forall \msc{t} \lnot ( \msc{t} \in [
    \msc{t_1} ,, \msc{t_2} ] \wedge \mw{\varphi} )
    )\label{eq:bl.excc}\tag{bl.excc}
\end{align}
\end{lemma}

\begin{proof}
\begin{align}
  1 &&  & \vdash ( \exists \msc{t} ( \msc{t} \in [ \msc{t_1} ,, \msc{t_2} ] \wedge \mw{\varphi} ) \leftrightarrow \lnot \forall \msc{t} \lnot ( \msc{t} \in [ \msc{t_1} ,, \msc{t_2} ] \wedge \mw{\varphi} ) )&(\mbox{\ref{eq:df-ex}})\notag
\end{align}
\end{proof}

\begin{lemma}\blTitle{bl.excd} Existential quantification over open-left
interval
\begin{align}
\vdash ( \exists \msc{t} ( \msc{t} \in [ \msc{t_1} ,. \msc{t_2} ] \wedge
    \mw{\varphi} ) \leftrightarrow \lnot \forall \msc{t} \lnot ( \msc{t} \in [
    \msc{t_1} ,. \msc{t_2} ] \wedge \mw{\varphi} )
    )\label{eq:bl.excd}\tag{bl.excd}
\end{align}
\end{lemma}

\begin{proof}
\begin{align}
  1 &&  & \vdash ( \exists \msc{t} ( \msc{t} \in [ \msc{t_1} ,. \msc{t_2} ] \wedge \mw{\varphi} ) \leftrightarrow \lnot \forall \msc{t} \lnot ( \msc{t} \in [ \msc{t_1} ,. \msc{t_2} ] \wedge \mw{\varphi} ) )&(\mbox{\ref{eq:df-ex}})\notag
\end{align}
\end{proof}

\begin{lemma}\blTitle{bl.exdc} Existential quantification over open-right
interval
\begin{align}
\vdash ( \exists \msc{t} ( \msc{t} \in [ \msc{t_1} ., \msc{t_2} ] \wedge
    \mw{\varphi} ) \leftrightarrow \lnot \forall \msc{t} \lnot ( \msc{t} \in [
    \msc{t_1} ., \msc{t_2} ] \wedge \mw{\varphi} )
    )\label{eq:bl.exdc}\tag{bl.exdc}
\end{align}
\end{lemma}

\begin{proof}
\begin{align}
  1 &&  & \vdash ( \exists \msc{t} ( \msc{t} \in [ \msc{t_1} ., \msc{t_2} ] \wedge \mw{\varphi} ) \leftrightarrow \lnot \forall \msc{t} \lnot ( \msc{t} \in [ \msc{t_1} ., \msc{t_2} ] \wedge \mw{\varphi} ) )&(\mbox{\ref{eq:df-ex}})\notag
\end{align}
\end{proof}

\begin{metadefinition}\blTitle{df-bl.atad2} Distribute @ over addition, two
terms
\begin{align}
\vdash ( ( \msc{t} \in \mathbf{T} \wedge B \in \mathbb{C} \wedge C \in
    \mathbb{C} ) \rightarrow ( ( B + C ) @ \msc{t} ) = ( ( B @ \msc{t} ) + ( C
    @ \msc{t} ) ) )\label{eq:df-bl.atad2}\tag{df-bl.atad2}
\end{align}
\end{metadefinition}

\begin{metadefinition}\blTitle{df-bl.atmul2} Distribute @ over multiplication,
two terms
\begin{align}
\vdash ( ( \msc{t} \in \mathbf{T} \wedge B \in \mathbb{C} \wedge C \in
    \mathbb{C} ) \rightarrow ( ( B \cdot C ) @ \msc{t} ) = ( ( B @ \msc{t} )
    \cdot ( C @ \msc{t} ) ) )\label{eq:df-bl.atmul2}\tag{df-bl.atmul2}
\end{align}
\end{metadefinition}

\begin{metadefinition}\blTitle{df-bl.atsub} Distribute @ over subtraction
\begin{align}
\vdash ( ( \msc{t} \in \mathbf{T} \wedge B \in \mathbb{C} \wedge C \in
    \mathbb{C} ) \rightarrow ( ( B \minus C ) @ \msc{t} ) = ( ( B @ \msc{t} )
    \minus ( C @ \msc{t} ) ) )\label{eq:df-bl.atsub}\tag{df-bl.atsub}
\end{align}
\end{metadefinition}

\begin{metadefinition}\blTitle{df-bl.atdiv} Distribute @ over division
\begin{align}
\vdash ( ( \msc{t} \in \mathbf{T} \wedge B \in \mathbb{C} \wedge C \in
    \mathbb{C} ) \rightarrow ( ( B / C ) @ \msc{t} ) = ( ( B @ \msc{t} ) / ( C
    @ \msc{t} ) ) )\label{eq:df-bl.atdiv}\tag{df-bl.atdiv}
\end{align}
\end{metadefinition}

\begin{metadefinition}\blTitle{df-bl.atan2} Distribute @ over conjunction, two
terms
\begin{align}
\vdash ( \msc{t} \in \mathbf{T} \rightarrow ( ( ( \mw{\varphi} \wedge \mw{\psi}
    ) @ \msc{t} ) \leftrightarrow ( ( \mw{\varphi} @ \msc{t} ) \wedge (
    \mw{\psi} @ \msc{t} ) ) ) )\label{eq:df-bl.atan2}\tag{df-bl.atan2}
\end{align}
\end{metadefinition}

\begin{metadefinition}\blTitle{df-bl.ator2} Distribute @ over disjuntion, two
terms
\begin{align}
\vdash ( \msc{t} \in \mathbf{T} \rightarrow ( ( ( \mw{\varphi} \vee \mw{\psi} )
    @ \msc{t} ) \leftrightarrow ( ( \mw{\varphi} @ \msc{t} ) \vee ( \mw{\psi} @
    \msc{t} ) ) ) )\label{eq:df-bl.ator2}\tag{df-bl.ator2}
\end{align}
\end{metadefinition}

\begin{lemma}\blTitle{bl.atad3} Distribute @ over addition, three terms $(
A + ( B + C ) ) @ \msc{t} )$.
\begin{align}
\vdash \msc{t} \in \mathbf{T}\quad\&\quad \vdash A \in \mathbb{C}\quad\&\quad
    \vdash B \in \mathbb{C}\quad\&\quad \vdash C \in
    \mathbb{C}\quad\Rightarrow\quad \vdash ( ( A + ( B + C ) ) @ \msc{t} ) = (
    ( A @ \msc{t} ) + ( ( B @ \msc{t} ) + ( C @ \msc{t} ) )
    )\label{eq:bl.atad3}\tag{bl.atad3}
\end{align}
\end{lemma}

\begin{proof}
\begin{align}
  1 &&  & \vdash \msc{t} \in \mathbf{T}&\text{Hyp~atad3.1}\notag%bl.atad3.1
  \\ 2 &&  & \vdash A \in \mathbb{C}&\text{Hyp~atad3.2}\notag%bl.atad3.2
  \\ 3 &&  & \vdash B \in \mathbb{C}&\text{Hyp~atad3.3}\notag%bl.atad3.3
  \\ 4 &&  & \vdash C \in \mathbb{C}&\text{Hyp~atad3.4}\notag%bl.atad3.4
  \\ 5 &&  & \vdash ( B + C ) \in \mathbb{C}&(\mbox{\ref{eq:addcli}},3,4)\notag
  \\ 6 &&  & \vdash ( \msc{t} \in \mathbf{T} \wedge A \in \mathbb{C} \wedge ( B + C ) \in \mathbb{C} )&(\mbox{\ref{eq:3pm3.2i}},1,2,5)\notag
  \\ 7 &&  & \vdash ( ( \msc{t} \in \mathbf{T} \wedge A \in \mathbb{C} \wedge ( B + C ) \in \mathbb{C} ) \rightarrow ( ( A + ( B + C ) ) @ \msc{t} ) = ( ( A @ \msc{t} ) + ( ( B + C ) @ \msc{t} ) ) )&(\mbox{\ref{eq:df-bl.atad2}})\notag
  \\ 8 &&  & \vdash ( ( A + ( B + C ) ) @ \msc{t} ) = ( ( A @ \msc{t} ) + ( ( B + C ) @ \msc{t} ) )&(\mbox{\ref{eq:ax-mp}},6,7)\notag
  \\ 9 &&  & \vdash \msc{t} \in \mathbf{T}&\text{Hyp~atad3.1}\notag%bl.atad3.1
  \\ 10 &&  & \vdash B \in \mathbb{C}&\text{Hyp~atad3.3}\notag%bl.atad3.3
  \\ 11 &&  & \vdash C \in \mathbb{C}&\text{Hyp~atad3.4}\notag%bl.atad3.4
  \\ 12 &&  & \vdash ( \msc{t} \in \mathbf{T} \wedge B \in \mathbb{C} \wedge C \in \mathbb{C} )&(\mbox{\ref{eq:3pm3.2i}},9,10,11)\notag
  \\ 13 &&  & \vdash ( ( \msc{t} \in \mathbf{T} \wedge B \in \mathbb{C} \wedge C \in \mathbb{C} ) \rightarrow ( ( B + C ) @ \msc{t} ) = ( ( B @ \msc{t} ) + ( C @ \msc{t} ) ) )&(\mbox{\ref{eq:df-bl.atad2}})\notag
  \\ 14 &&  & \vdash ( ( B + C ) @ \msc{t} ) = ( ( B @ \msc{t} ) + ( C @ \msc{t} ) )&(\mbox{\ref{eq:ax-mp}},12,13)\notag
  \\ 15 &&  & \vdash ( ( A @ \msc{t} ) + ( ( B + C ) @ \msc{t} ) ) = ( ( A @ \msc{t} ) + ( ( B @ \msc{t} ) + ( C @ \msc{t} ) ) )&(\mbox{\ref{eq:oveq2i}},14)\notag
  \\ 16 &&  & \vdash ( ( A + ( B + C ) ) @ \msc{t} ) = ( ( A @ \msc{t} ) + ( ( B @ \msc{t} ) + ( C @ \msc{t} ) ) )&(\mbox{\ref{eq:eqtri}},8,15)\notag
\end{align}
\end{proof}

\begin{lemma}\blTitle{bl.atad3l} Distribute @ over addition, three terms,
left $( ( A + B ) + C ) @ \msc{t} )$.
\begin{align}
\vdash \msc{t} \in \mathbf{T}\quad\&\quad \vdash A \in \mathbb{C}\quad\&\quad
    \vdash B \in \mathbb{C}\quad\&\quad \vdash C \in
    \mathbb{C}\quad\Rightarrow\quad \vdash ( ( ( A + B ) + C ) @ \msc{t} ) = (
    ( ( A @ \msc{t} ) + ( B @ \msc{t} ) ) + ( C @ \msc{t} )
    )\label{eq:bl.atad3l}\tag{bl.atad3l}
\end{align}
\end{lemma}

\begin{proof}
\begin{align}
  1 &&  & \vdash \msc{t} \in \mathbf{T}&\text{Hyp~atad3l.1}\notag%bl.atad3l.1
  \\ 2 &&  & \vdash A \in \mathbb{C}&\text{Hyp~atad3l.2}\notag%bl.atad3l.2
  \\ 3 &&  & \vdash B \in \mathbb{C}&\text{Hyp~atad3l.3}\notag%bl.atad3l.3
  \\ 4 &&  & \vdash ( A + B ) \in \mathbb{C}&(\mbox{\ref{eq:addcli}},2,3)\notag
  \\ 5 &&  & \vdash C \in \mathbb{C}&\text{Hyp~atad3l.4}\notag%bl.atad3l.4
  \\ 6 &&  & \vdash ( \msc{t} \in \mathbf{T} \wedge ( A + B ) \in \mathbb{C} \wedge C \in \mathbb{C} )&(\mbox{\ref{eq:3pm3.2i}},1,4,5)\notag
  \\ 7 &&  & \vdash ( ( \msc{t} \in \mathbf{T} \wedge ( A + B ) \in \mathbb{C} \wedge C \in \mathbb{C} ) \rightarrow ( ( ( A + B ) + C ) @ \msc{t} ) = ( ( ( A + B ) @ \msc{t} ) + ( C @ \msc{t} ) ) )&(\mbox{\ref{eq:df-bl.atad2}})\notag
  \\ 8 &&  & \vdash ( ( ( A + B ) + C ) @ \msc{t} ) = ( ( ( A + B ) @ \msc{t} ) + ( C @ \msc{t} ) )&(\mbox{\ref{eq:ax-mp}},6,7)\notag
  \\ 9 &&  & \vdash \msc{t} \in \mathbf{T}&\text{Hyp~atad3l.1}\notag%bl.atad3l.1
  \\ 10 &&  & \vdash A \in \mathbb{C}&\text{Hyp~atad3l.2}\notag%bl.atad3l.2
  \\ 11 &&  & \vdash B \in \mathbb{C}&\text{Hyp~atad3l.3}\notag%bl.atad3l.3
  \\ 12 &&  & \vdash ( \msc{t} \in \mathbf{T} \wedge A \in \mathbb{C} \wedge B \in \mathbb{C} )&(\mbox{\ref{eq:3pm3.2i}},9,10,11)\notag
  \\ 13 &&  & \vdash ( ( \msc{t} \in \mathbf{T} \wedge A \in \mathbb{C} \wedge B \in \mathbb{C} ) \rightarrow ( ( A + B ) @ \msc{t} ) = ( ( A @ \msc{t} ) + ( B @ \msc{t} ) ) )&(\mbox{\ref{eq:df-bl.atad2}})\notag
  \\ 14 &&  & \vdash ( ( A + B ) @ \msc{t} ) = ( ( A @ \msc{t} ) + ( B @ \msc{t} ) )&(\mbox{\ref{eq:ax-mp}},12,13)\notag
  \\ 15 &&  & \vdash ( ( ( A + B ) @ \msc{t} ) + ( C @ \msc{t} ) ) = ( ( ( A @ \msc{t} ) + ( B @ \msc{t} ) ) + ( C @ \msc{t} ) )&(\mbox{\ref{eq:oveq1i}},14)\notag
  \\ 16 &&  & \vdash ( ( ( A + B ) + C ) @ \msc{t} ) = ( ( ( A @ \msc{t} ) + ( B @ \msc{t} ) ) + ( C @ \msc{t} ) )&(\mbox{\ref{eq:eqtri}},8,15)\notag
\end{align}
\end{proof}

\begin{lemma}\blTitle{bl.atmul3} Distribute @ over multiplication, three
terms $( A \cdot ( B \cdot C ) ) @ \msc{t} )$.
\begin{align}
\vdash \msc{t} \in \mathbf{T}\quad\&\quad \vdash A \in \mathbb{C}\quad\&\quad
    \vdash B \in \mathbb{C}\quad\&\quad \vdash C \in
    \mathbb{C}\quad\Rightarrow\quad \vdash ( ( A \cdot ( B \cdot C ) ) @
    \msc{t} ) = ( ( A @ \msc{t} ) \cdot ( ( B @ \msc{t} ) \cdot ( C @ \msc{t} )
    ) )\label{eq:bl.atmul3}\tag{bl.atmul3}
\end{align}
\end{lemma}

\begin{proof}
\begin{align}
  1 &&  & \vdash \msc{t} \in \mathbf{T}&\text{Hyp~atmul3.1}\notag%bl.atmul3.1
  \\ 2 &&  & \vdash A \in \mathbb{C}&\text{Hyp~atmul3.2}\notag%bl.atmul3.2
  \\ 3 &&  & \vdash B \in \mathbb{C}&\text{Hyp~atmul3.3}\notag%bl.atmul3.3
  \\ 4 &&  & \vdash C \in \mathbb{C}&\text{Hyp~atmul3.4}\notag%bl.atmul3.4
  \\ 5 &&  & \vdash ( B \cdot C ) \in \mathbb{C}&(\mbox{\ref{eq:mulcli}},3,4)\notag
  \\ 6 &&  & \vdash ( \msc{t} \in \mathbf{T} \wedge A \in \mathbb{C} \wedge ( B \cdot C ) \in \mathbb{C} )&(\mbox{\ref{eq:3pm3.2i}},1,2,5)\notag
  \\ 7 &&  & \vdash ( ( \msc{t} \in \mathbf{T} \wedge A \in \mathbb{C} \wedge ( B \cdot C ) \in \mathbb{C} ) \rightarrow ( ( A \cdot ( B \cdot C ) ) @ \msc{t} ) = ( ( A @ \msc{t} ) \cdot ( ( B \cdot C ) @ \msc{t} ) ) )&(\mbox{\ref{eq:df-bl.atmul2}})\notag
  \\ 8 &&  & \vdash ( ( A \cdot ( B \cdot C ) ) @ \msc{t} ) = ( ( A @ \msc{t} ) \cdot ( ( B \cdot C ) @ \msc{t} ) )&(\mbox{\ref{eq:ax-mp}},6,7)\notag
  \\ 9 &&  & \vdash \msc{t} \in \mathbf{T}&\text{Hyp~atmul3.1}\notag%bl.atmul3.1
  \\ 10 &&  & \vdash B \in \mathbb{C}&\text{Hyp~atmul3.3}\notag%bl.atmul3.3
  \\ 11 &&  & \vdash C \in \mathbb{C}&\text{Hyp~atmul3.4}\notag%bl.atmul3.4
  \\ 12 &&  & \vdash ( \msc{t} \in \mathbf{T} \wedge B \in \mathbb{C} \wedge C \in \mathbb{C} )&(\mbox{\ref{eq:3pm3.2i}},9,10,11)\notag
  \\ 13 &&  & \vdash ( ( \msc{t} \in \mathbf{T} \wedge B \in \mathbb{C} \wedge C \in \mathbb{C} ) \rightarrow ( ( B \cdot C ) @ \msc{t} ) = ( ( B @ \msc{t} ) \cdot ( C @ \msc{t} ) ) )&(\mbox{\ref{eq:df-bl.atmul2}})\notag
  \\ 14 &&  & \vdash ( ( B \cdot C ) @ \msc{t} ) = ( ( B @ \msc{t} ) \cdot ( C @ \msc{t} ) )&(\mbox{\ref{eq:ax-mp}},12,13)\notag
  \\ 15 &&  & \vdash ( ( A @ \msc{t} ) \cdot ( ( B \cdot C ) @ \msc{t} ) ) = ( ( A @ \msc{t} ) \cdot ( ( B @ \msc{t} ) \cdot ( C @ \msc{t} ) ) )&(\mbox{\ref{eq:oveq2i}},14)\notag
  \\ 16 &&  & \vdash ( ( A \cdot ( B \cdot C ) ) @ \msc{t} ) = ( ( A @ \msc{t} ) \cdot ( ( B @ \msc{t} ) \cdot ( C @ \msc{t} ) ) )&(\mbox{\ref{eq:eqtri}},8,15)\notag
\end{align}
\end{proof}

\begin{lemma}\blTitle{bl.atmul3l} Distribute @ over multiplication, three
terms, left $( ( A \cdot B ) \cdot C ) @ \msc{t} )$.
\begin{align}
\vdash \msc{t} \in \mathbf{T}\quad\&\quad \vdash A \in \mathbb{C}\quad\&\quad
    \vdash B \in \mathbb{C}\quad\&\quad \vdash C \in
    \mathbb{C}\quad\Rightarrow\quad \vdash ( ( ( A \cdot B ) \cdot C ) @
    \msc{t} ) = ( ( ( A @ \msc{t} ) \cdot ( B @ \msc{t} ) ) \cdot ( C @ \msc{t}
    ) )\label{eq:bl.atmul3l}\tag{bl.atmul3l}
\end{align}
\end{lemma}

\begin{proof}
\begin{align}
  1 &&  & \vdash \msc{t} \in \mathbf{T}&\text{Hyp~atmul3l.1}\notag%bl.atmul3l.1
  \\ 2 &&  & \vdash A \in \mathbb{C}&\text{Hyp~atmul3l.2}\notag%bl.atmul3l.2
  \\ 3 &&  & \vdash B \in \mathbb{C}&\text{Hyp~atmul3l.3}\notag%bl.atmul3l.3
  \\ 4 &&  & \vdash ( A \cdot B ) \in \mathbb{C}&(\mbox{\ref{eq:mulcli}},2,3)\notag
  \\ 5 &&  & \vdash C \in \mathbb{C}&\text{Hyp~atmul3l.4}\notag%bl.atmul3l.4
  \\ 6 &&  & \vdash ( \msc{t} \in \mathbf{T} \wedge ( A \cdot B ) \in \mathbb{C} \wedge C \in \mathbb{C} )&(\mbox{\ref{eq:3pm3.2i}},1,4,5)\notag
  \\ 7 &&  & \vdash ( ( \msc{t} \in \mathbf{T} \wedge ( A \cdot B ) \in \mathbb{C} \wedge C \in \mathbb{C} ) \rightarrow ( ( ( A \cdot B ) \cdot C ) @ \msc{t} ) = ( ( ( A \cdot B ) @ \msc{t} ) \cdot ( C @ \msc{t} ) ) )&(\mbox{\ref{eq:df-bl.atmul2}})\notag
  \\ 8 &&  & \vdash ( ( ( A \cdot B ) \cdot C ) @ \msc{t} ) = ( ( ( A \cdot B ) @ \msc{t} ) \cdot ( C @ \msc{t} ) )&(\mbox{\ref{eq:ax-mp}},6,7)\notag
  \\ 9 &&  & \vdash \msc{t} \in \mathbf{T}&\text{Hyp~atmul3l.1}\notag%bl.atmul3l.1
  \\ 10 &&  & \vdash A \in \mathbb{C}&\text{Hyp~atmul3l.2}\notag%bl.atmul3l.2
  \\ 11 &&  & \vdash B \in \mathbb{C}&\text{Hyp~atmul3l.3}\notag%bl.atmul3l.3
  \\ 12 &&  & \vdash ( \msc{t} \in \mathbf{T} \wedge A \in \mathbb{C} \wedge B \in \mathbb{C} )&(\mbox{\ref{eq:3pm3.2i}},9,10,11)\notag
  \\ 13 &&  & \vdash ( ( \msc{t} \in \mathbf{T} \wedge A \in \mathbb{C} \wedge B \in \mathbb{C} ) \rightarrow ( ( A \cdot B ) @ \msc{t} ) = ( ( A @ \msc{t} ) \cdot ( B @ \msc{t} ) ) )&(\mbox{\ref{eq:df-bl.atmul2}})\notag
  \\ 14 &&  & \vdash ( ( A \cdot B ) @ \msc{t} ) = ( ( A @ \msc{t} ) \cdot ( B @ \msc{t} ) )&(\mbox{\ref{eq:ax-mp}},12,13)\notag
  \\ 15 &&  & \vdash ( ( ( A \cdot B ) @ \msc{t} ) \cdot ( C @ \msc{t} ) ) = ( ( ( A @ \msc{t} ) \cdot ( B @ \msc{t} ) ) \cdot ( C @ \msc{t} ) )&(\mbox{\ref{eq:oveq1i}},14)\notag
  \\ 16 &&  & \vdash ( ( ( A \cdot B ) \cdot C ) @ \msc{t} ) = ( ( ( A @ \msc{t} ) \cdot ( B @ \msc{t} ) ) \cdot ( C @ \msc{t} ) )&(\mbox{\ref{eq:eqtri}},8,15)\notag
\end{align}
\end{proof}

\begin{lemma}\blTitle{bl.atan3} Distribute @ over conjunction, three terms
$\mw{\varphi} \wedge ( \mw{\psi} \wedge \mw{\chi} ) \,\mathrm{{\tt '}}$
\begin{align}
\vdash ( \msc{t} \in \mathbf{T} \rightarrow ( ( ( \mw{\varphi} \wedge (
    \mw{\psi} \wedge \mw{\chi} ) ) @ \msc{t} ) \leftrightarrow ( ( \mw{\varphi}
    @ \msc{t} ) \wedge ( ( \mw{\psi} @ \msc{t} ) \wedge ( \mw{\chi} @ \msc{t} )
    ) ) ) )\label{eq:bl.atan3}\tag{bl.atan3}
\end{align}
\end{lemma}

\begin{proof}
\begin{align}
  1 &&  & \vdash ( \msc{t} \in \mathbf{T} \rightarrow ( ( ( \mw{\varphi} \wedge ( \mw{\psi} \wedge \mw{\chi} ) ) @ \msc{t} ) \leftrightarrow ( ( \mw{\varphi} @ \msc{t} ) \wedge ( ( \mw{\psi} \wedge \mw{\chi} ) @ \msc{t} ) ) ) )&(\mbox{\ref{eq:df-bl.atan2}})\notag
  \\ 2 &&  & \vdash ( \msc{t} \in \mathbf{T} \rightarrow ( ( ( \mw{\psi} \wedge \mw{\chi} ) @ \msc{t} ) \leftrightarrow ( ( \mw{\psi} @ \msc{t} ) \wedge ( \mw{\chi} @ \msc{t} ) ) ) )&(\mbox{\ref{eq:df-bl.atan2}})\notag
  \\ 3 &&  & \vdash ( \msc{t} \in \mathbf{T} \rightarrow ( ( ( \mw{\varphi} @ \msc{t} ) \wedge ( ( \mw{\psi} \wedge \mw{\chi} ) @ \msc{t} ) ) \leftrightarrow ( ( \mw{\varphi} @ \msc{t} ) \wedge ( ( \mw{\psi} @ \msc{t} ) \wedge ( \mw{\chi} @ \msc{t} ) ) ) ) )&(\mbox{\ref{eq:anbi2d}},2)\notag
  \\ 4 &&  & \vdash ( \msc{t} \in \mathbf{T} \rightarrow ( ( ( \mw{\varphi} \wedge ( \mw{\psi} \wedge \mw{\chi} ) ) @ \msc{t} ) \leftrightarrow ( ( \mw{\varphi} @ \msc{t} ) \wedge ( ( \mw{\psi} @ \msc{t} ) \wedge ( \mw{\chi} @ \msc{t} ) ) ) ) )&(\mbox{\ref{eq:bitrd}},1,3)\notag
\end{align}
\end{proof}

\begin{lemma}\blTitle{bl.ator3} Distribute @ over disjunction, three terms
$\mw{\varphi} \vee ( \mw{\psi} \vee \mw{\chi} ) \,\mathrm{{\tt '}}$
\begin{align}
\vdash ( \msc{t} \in \mathbf{T} \rightarrow ( ( ( \mw{\varphi} \vee ( \mw{\psi}
    \vee \mw{\chi} ) ) @ \msc{t} ) \leftrightarrow ( ( \mw{\varphi} @ \msc{t} )
    \vee ( ( \mw{\psi} @ \msc{t} ) \vee ( \mw{\chi} @ \msc{t} ) ) ) )
    )\label{eq:bl.ator3}\tag{bl.ator3}
\end{align}
\end{lemma}

\begin{proof}
\begin{align}
  1 &&  & \vdash ( \msc{t} \in \mathbf{T} \rightarrow ( ( ( \mw{\varphi} \vee ( \mw{\psi} \vee \mw{\chi} ) ) @ \msc{t} ) \leftrightarrow ( ( \mw{\varphi} @ \msc{t} ) \vee ( ( \mw{\psi} \vee \mw{\chi} ) @ \msc{t} ) ) ) )&(\mbox{\ref{eq:df-bl.ator2}})\notag
  \\ 2 &&  & \vdash ( \msc{t} \in \mathbf{T} \rightarrow ( ( ( \mw{\psi} \vee \mw{\chi} ) @ \msc{t} ) \leftrightarrow ( ( \mw{\psi} @ \msc{t} ) \vee ( \mw{\chi} @ \msc{t} ) ) ) )&(\mbox{\ref{eq:df-bl.ator2}})\notag
  \\ 3 &&  & \vdash ( \msc{t} \in \mathbf{T} \rightarrow ( ( ( \mw{\varphi} @ \msc{t} ) \vee ( ( \mw{\psi} \vee \mw{\chi} ) @ \msc{t} ) ) \leftrightarrow ( ( \mw{\varphi} @ \msc{t} ) \vee ( ( \mw{\psi} @ \msc{t} ) \vee ( \mw{\chi} @ \msc{t} ) ) ) ) )&(\mbox{\ref{eq:orbi2d}},2)\notag
  \\ 4 &&  & \vdash ( \msc{t} \in \mathbf{T} \rightarrow ( ( ( \mw{\varphi} \vee ( \mw{\psi} \vee \mw{\chi} ) ) @ \msc{t} ) \leftrightarrow ( ( \mw{\varphi} @ \msc{t} ) \vee ( ( \mw{\psi} @ \msc{t} ) \vee ( \mw{\chi} @ \msc{t} ) ) ) ) )&(\mbox{\ref{eq:bitrd}},1,3)\notag
\end{align}
\end{proof}

\begin{metadefinition}\blTitle{clts} Define $( A ~\string^~ B )$ as a class
\begin{align}
{\rm class} ( A ~\string^~ B )\label{eq:clts}\tag{clts}
\end{align}
\end{metadefinition}

\begin{metadefinition}\blTitle{df-bl.ts} Time-Shift Predicate ($~\string^~$ not
in $\mw{\varphi}$)
\begin{align}
\vdash ( ( A \in \mathbb{Z} \wedge \delta \in \mathbb{R} ) \rightarrow ( (
    \mw{\varphi} ~\string^~ A ) \leftrightarrow ( \mw{\varphi} @ (
    \mathsf{\msc{now}} + ( \delta \cdot A ) ) ) )
    )\label{eq:df-bl.ts}\tag{df-bl.ts}
\end{align}
\end{metadefinition}

\begin{lemma}\blTitle{bl.tsi} Time-Shift Predicate inference ($~\string^~$
not in $\mw{\varphi}$)
\begin{align}
\vdash A \in \mathbb{Z}\quad\&\quad \vdash \delta \in
    \mathbb{R}\quad\Rightarrow\quad \vdash ( ( \mw{\varphi} ~\string^~ A )
    \leftrightarrow ( \mw{\varphi} @ ( \mathsf{\msc{now}} + ( \delta \cdot A )
    ) ) )\label{eq:bl.tsi}\tag{bl.tsi}
\end{align}
\end{lemma}

\begin{proof}
\begin{align}
  1 &&  & \vdash A \in \mathbb{Z}&\text{Hyp~tsi.1}\notag%bl.tsi.1
  \\ 2 &&  & \vdash \delta \in \mathbb{R}&\text{Hyp~tsi.2}\notag%bl.tsi.2
  \\ 3 &&  & \vdash ( ( A \in \mathbb{Z} \wedge \delta \in \mathbb{R} ) \rightarrow ( ( \mw{\varphi} ~\string^~ A ) \leftrightarrow ( \mw{\varphi} @ ( \mathsf{\msc{now}} + ( \delta \cdot A ) ) ) ) )&(\mbox{\ref{eq:df-bl.ts}})\notag
  \\ 4 &&  & \vdash ( ( \mw{\varphi} ~\string^~ A ) \leftrightarrow ( \mw{\varphi} @ ( \mathsf{\msc{now}} + ( \delta \cdot A ) ) ) )&(\mbox{\ref{eq:mp2an}},1,2,3)\notag
\end{align}
\end{proof}

\begin{metadefinition}\blTitle{df-bl.tsc} Time-Shift Expression ($~\string^~$
not in $A$)
\begin{align}
\vdash ( ( A \in \mathbb{C} \wedge B \in \mathbb{Z} \wedge \delta \in
    \mathbb{R} ) \rightarrow ( A ~\string^~ B ) = ( A @ ( \mathsf{\msc{now}} +
    ( \delta \cdot B ) ) ) )\label{eq:df-bl.tsc}\tag{df-bl.tsc}
\end{align}
\end{metadefinition}

\begin{lemma}\blTitle{bl.tsci} Time-Shift Expression inference
($~\string^~$ not in $A$)
\begin{align}
\vdash A \in \mathbb{C}\quad\&\quad \vdash B \in \mathbb{Z}\quad\&\quad \vdash
    \delta \in \mathbb{R}\quad\Rightarrow\quad \vdash ( A ~\string^~ B ) = ( A
    @ ( \mathsf{\msc{now}} + ( \delta \cdot B ) )
    )\label{eq:bl.tsci}\tag{bl.tsci}
\end{align}
\end{lemma}

\begin{proof}
\begin{align}
  1 &&  & \vdash A \in \mathbb{C}&\text{Hyp~tsci.1}\notag%bl.tsci.1
  \\ 2 &&  & \vdash B \in \mathbb{Z}&\text{Hyp~tsci.2}\notag%bl.tsci.2
  \\ 3 &&  & \vdash \delta \in \mathbb{R}&\text{Hyp~tsci.3}\notag%bl.tsci.3
  \\ 4 &&  & \vdash ( ( A \in \mathbb{C} \wedge B \in \mathbb{Z} \wedge \delta \in \mathbb{R} ) \rightarrow ( A ~\string^~ B ) = ( A @ ( \mathsf{\msc{now}} + ( \delta \cdot B ) ) ) )&(\mbox{\ref{eq:df-bl.tsc}})\notag
  \\ 5 &&  & \vdash ( A ~\string^~ B ) = ( A @ ( \mathsf{\msc{now}} + ( \delta \cdot B ) ) )&(\mbox{\ref{eq:mp3an}},1,2,3,4)\notag
\end{align}
\end{proof}

\begin{metadefinition}\blTitle{df-bl.tsdis} $~\string^~$ distributes over
predicates
\begin{align}
\vdash ( ( A \in \mathbb{Z} \wedge B \in \mathbb{Z} ) \rightarrow ( ( (
    \mw{\varphi} ~\string^~ A ) ~\string^~ B ) \leftrightarrow ( \mw{\varphi}
    ~\string^~ ( A + B ) ) ) )\label{eq:df-bl.tsdis}\tag{df-bl.tsdis}
\end{align}
\end{metadefinition}

\begin{lemma}\blTitle{bl.tsdisi} $~\string^~$ distributes over predicates
\begin{align}
\vdash ( A \in \mathbb{Z} \wedge B \in \mathbb{Z} )\quad\Rightarrow\quad \vdash
    ( ( ( \mw{\varphi} ~\string^~ A ) ~\string^~ B ) \leftrightarrow (
    \mw{\varphi} ~\string^~ ( A + B ) ) )\label{eq:bl.tsdisi}\tag{bl.tsdisi}
\end{align}
\end{lemma}

\begin{proof}
\begin{align}
  1 &&  & \vdash ( A \in \mathbb{Z} \wedge B \in \mathbb{Z} )&\text{Hyp~tsdisi.1}\notag%bl.tsdisi.1
  \\ 2 &&  & \vdash ( ( A \in \mathbb{Z} \wedge B \in \mathbb{Z} ) \rightarrow ( ( ( \mw{\varphi} ~\string^~ A ) ~\string^~ B ) \leftrightarrow ( \mw{\varphi} ~\string^~ ( A + B ) ) ) )&(\mbox{\ref{eq:df-bl.tsdis}})\notag
  \\ 3 &&  & \vdash ( ( ( \mw{\varphi} ~\string^~ A ) ~\string^~ B ) \leftrightarrow ( \mw{\varphi} ~\string^~ ( A + B ) ) )&(\mbox{\ref{eq:ax-mp}},1,2)\notag
\end{align}
\end{proof}

\begin{lemma}\blTitle{bl.tsdisci} $~\string^~$ distributes over values
\begin{align}
\vdash ( A \in \mathbb{Z} \wedge B \in \mathbb{Z} \wedge C \in \mathbb{C}
    )\quad\Rightarrow\quad \vdash ( ( C ~\string^~ A ) ~\string^~ B ) = ( C
    ~\string^~ ( A + B ) )\label{eq:bl.tsdisci}\tag{bl.tsdisci}
\end{align}
\end{lemma}

\begin{proof}
\begin{align}
  1 &&  & \vdash ( A \in \mathbb{Z} \wedge B \in \mathbb{Z} \wedge C \in \mathbb{C} )&\text{Hyp~tsdisci.1}\notag%bl.tsdisci.1
  \\ 2 &&  & \vdash ( ( A \in \mathbb{Z} \wedge B \in \mathbb{Z} \wedge C \in \mathbb{C} ) \rightarrow ( ( C ~\string^~ A ) ~\string^~ B ) = ( C ~\string^~ ( A + B ) ) )&(\mbox{\ref{eq:df-bl.tsdisc}})\notag
  \\ 3 &&  & \vdash ( ( C ~\string^~ A ) ~\string^~ B ) = ( C ~\string^~ ( A + B ) )&(\mbox{\ref{eq:ax-mp}},1,2)\notag
\end{align}
\end{proof}

\begin{lemma}\blTitle{bl.ts0} zero time shift, predicate
\begin{align}
\vdash ( ( \mw{\varphi} ~\string^~ 0 ) \leftrightarrow ( \mw{\varphi} @
    \mathsf{\msc{now}} ) )\label{eq:bl.ts0}\tag{bl.ts0}
\end{align}
\end{lemma}

\begin{proof}
\begin{align}
  1 &&  & &(\mbox{\ref{eq:?}})\notag
\end{align}
\end{proof}

\begin{lemma}\blTitle{bl.cd} TO DO: PUT DESCRIPTION HERE
\begin{align}
\vdash ( \mw{\varphi} \rightarrow \mw{\psi} )\quad\&\quad \vdash ( \mw{\chi}
    \rightarrow \mw{\theta} )\quad\&\quad \vdash ( \mw{\varphi} \vee \mw{\chi}
    )\quad\Rightarrow\quad \vdash ( \mw{\psi} \vee \mw{\theta}
    )\label{eq:bl.cd}\tag{bl.cd}
\end{align}
\end{lemma}

\begin{proof}
\begin{align}
  1 &&  & \vdash ( \mw{\varphi} \vee \mw{\chi} )&\text{Hyp~cd.3}\notag%bl.cd.3
  \\ 2 &&  & \vdash ( \mw{\varphi} \rightarrow \mw{\psi} )&\text{Hyp~cd.1}\notag%bl.cd.1
  \\ 3 &&  & \vdash ( \mw{\chi} \rightarrow \mw{\theta} )&\text{Hyp~cd.2}\notag%bl.cd.2
  \\ 4 &&  & \vdash ( ( \mw{\varphi} \vee \mw{\chi} ) \rightarrow ( \mw{\psi} \vee \mw{\theta} ) )&(\mbox{\ref{eq:orim12i}},2,3)\notag
  \\ 5 &&  & \vdash ( \mw{\psi} \vee \mw{\theta} )&(\mbox{\ref{eq:ax-mp}},1,4)\notag
\end{align}
\end{proof}

\begin{metadefinition}\blTitle{df-bl.id} Create Identifier Class; first
character must be a letter,
others may be letter or digit; underscore not allowed

\begin{align}
\vdash \,\mathrm{ID} = ( \,\mathrm{Word } \,\mathrm{Letter} \,\mathrm{ ++ }
    \,\mathrm{Word } ( \,\mathrm{Letter} \cup \,\mathrm{Digit} )
    )\label{eq:df-bl.id}\tag{df-bl.id}
\end{align}
\end{metadefinition}

\begin{metadefinition}\blTitle{df-bl.enum} Create Class Enum which contains all
identifiers;
Every enumeration will be a subset of Enum: $x \subsetneq \,\mathrm{Enum}$

\begin{align}
\vdash \,\mathrm{Enum} = \{ x~\vert~x \in \,\mathrm{ID}
    \}\label{eq:df-bl.enum}\tag{df-bl.enum}
\end{align}
\end{metadefinition}

\begin{metadefinition}\blTitle{df-bl.qt} A quantity is a pair of a number with a
unit.  A quantity type is a set of
pairs, where the unit is fixed, and the number may be restricted in range or
granularity.
$x$ is a quantity type iff $x$ is a pair of a real number and a unit
identifier.

\begin{align}
\vdash ( x \in \,\mathrm{QT} \leftrightarrow \exists u \in \,\mathrm{ID} x \in
    \{ \langle y , u \rangle~\vert~y \in \mathbb{C} \}
    )\label{eq:df-bl.qt}\tag{df-bl.qt}
\end{align}
\end{metadefinition}


\begin{lemma}\blTitle{bl.arst} An array is an extensible structure.
\begin{align}
\vdash ( F \,\mathrm{Array} \langle M , N \rangle \rightarrow F {\rm Struct}
    \langle M , N \rangle )\label{eq:bl.arst}\tag{bl.arst}
\end{align}
\end{lemma}

\begin{proof}
\begin{align}
  1 &&  & \vdash ( F \,\mathrm{Array} \langle M , N \rangle \leftrightarrow ( ( M \in \mathbb{N} \wedge N \in \mathbb{N} \wedge M \le N ) \wedge ( {\rm Fun} ( F \setminus \{ \varnothing \} ) \wedge {\rm dom} F \subseteq ( M \ldots N ) ) \wedge {\rm ran} F \subseteq \mathsf{TYPE} ) )&(\mbox{\ref{eq:df-bl.array}})\notag
  \\ 2 &&  & \vdash ( ( ( M \in \mathbb{N} \wedge N \in \mathbb{N} \wedge M \le N ) \wedge ( {\rm Fun} ( F \setminus \{ \varnothing \} ) \wedge {\rm dom} F \subseteq ( M \ldots N ) ) \wedge {\rm ran} F \subseteq \mathsf{TYPE} ) \rightarrow ( ( M \in \mathbb{N} \wedge N \in \mathbb{N} \wedge M \le N ) \wedge ( {\rm Fun} ( F \setminus \{ \varnothing \} ) \wedge {\rm dom} F \subseteq ( M \ldots N ) ) ) )&(\mbox{\ref{eq:3simpa}})\notag
  \\ 3 &&  & \vdash ( F \,\mathrm{Array} \langle M , N \rangle \rightarrow ( ( M \in \mathbb{N} \wedge N \in \mathbb{N} \wedge M \le N ) \wedge ( {\rm Fun} ( F \setminus \{ \varnothing \} ) \wedge {\rm dom} F \subseteq ( M \ldots N ) ) ) )&(\mbox{\ref{eq:sylbi}},1,2)\notag
  \\ 4 &&  & \vdash ( ( ( M \in \mathbb{N} \wedge N \in \mathbb{N} \wedge M \le N ) \wedge {\rm Fun} ( F \setminus \{ \varnothing \} ) \wedge {\rm dom} F \subseteq ( M \ldots N ) ) \leftrightarrow ( ( M \in \mathbb{N} \wedge N \in \mathbb{N} \wedge M \le N ) \wedge ( {\rm Fun} ( F \setminus \{ \varnothing \} ) \wedge {\rm dom} F \subseteq ( M \ldots N ) ) ) )&(\mbox{\ref{eq:3anass}})\notag
  \\ 5 &&  & \vdash ( F \,\mathrm{Array} \langle M , N \rangle \rightarrow ( ( M \in \mathbb{N} \wedge N \in \mathbb{N} \wedge M \le N ) \wedge {\rm Fun} ( F \setminus \{ \varnothing \} ) \wedge {\rm dom} F \subseteq ( M \ldots N ) ) )&(\mbox{\ref{eq:sylibr}},3,4)\notag
  \\ 6 &&  & \vdash ( F {\rm Struct} \langle M , N \rangle \leftrightarrow ( ( M \in \mathbb{N} \wedge N \in \mathbb{N} \wedge M \le N ) \wedge {\rm Fun} ( F \setminus \{ \varnothing \} ) \wedge {\rm dom} F \subseteq ( M \ldots N ) ) )&(\mbox{\ref{eq:isstruct}})\notag
  \\ 7 &&  & \vdash ( F \,\mathrm{Array} \langle M , N \rangle \rightarrow F {\rm Struct} \langle M , N \rangle )&(\mbox{\ref{eq:sylibr}},5,6)\notag
\end{align}
\end{proof}

%\begin{metadefinition}\blTitle{df-bl.record} Record type is a function from
%identifiers to types.
%\begin{align}
%\vdash ( x \in \,\mathrm{Record} \leftrightarrow x : \,\mathrm{ID}
%    \longrightarrow \mathsf{TYPE} )\label{eq:df-bl.record}\tag{df-bl.record}
%\end{align}
%\end{metadefinition}

%\begin{metadefinition}\blTitle{df-bl.variant} Variant type uses a discriminant
%to determine which record field holds its type.
%\begin{align}
%\vdash ( x \in \,\mathrm{Variant} \leftrightarrow x \in \{ \langle y , z
%    \rangle~\vert~( z \in \,\mathrm{Record} \wedge y \in {\rm dom} z ) \}
%    )\label{eq:df-bl.variant}\tag{df-bl.variant}
%\end{align}
%\end{metadefinition}

%\begin{metadefinition}\blTitle{df-bl.string} String type is a Word of QWERTY
%characters
%\begin{align}
%\vdash x \in \,\mathrm{Word }
%    \,\mathrm{QWERTY}\label{eq:df-bl.string}\tag{df-bl.string}
%\end{align}
%\end{metadefinition}

%\begin{metadefinition}\blTitle{df-bl.type} Define constructed type: \\~{\tt
%df-bl.qt:} $x \in \,\mathrm{QT}$ or $x \subsetneq \ldots$ for whole numbers,
%\\~{\tt df-bl.array:} $\mathbb{N} \longrightarrow \mathsf{TYPE}$ \\~{\tt
%df-bl.string:} $\,\mathrm{Word } \,\mathrm{QWERTY}$ \\~{\tt df-bl.record:}
%$\,\mathrm{ID} \longrightarrow \mathsf{TYPE}$, \\~{\tt df-bl.variant:}
%$\langle \,\mathrm{ID} , \,\mathrm{ID} \longrightarrow \mathsf{TYPE} \rangle
%\,\mathrm{{\tt \char`\\}{\tt \char`\\}{\tt \char`\~}{\tt \char`\~}{\tt
%\char`\~}} \,\mathrm{df{\tt -}bl{\tt .}enum{\tt :}}$ x C- \{ y / y e. ID \} $,
%\,\mathrm{{\tt \char`\\}{\tt \char`\\}{\tt \char`\~}{\tt \char`\~}{\tt
%\char`\~}} \,\mathrm{df{\tt -}bl{\tt .}bool{\tt :}}$ \{ true , false \} $$
%\begin{align}
%\vdash ( x \in \mathsf{TYPE} \leftrightarrow \bigvee( x \in \,\mathrm{QT} x
%    \subsetneq \ldots ( x \in F \wedge F \,\mathrm{Array} \langle M , N \rangle
%    ) x \in \,\mathrm{Word } \,\mathrm{QWERTY} x \in \,\mathrm{Record} x \in
%    \,\mathrm{Variant} x \subseteq \,\mathrm{Enum} x \in \,\mathrm{Bool} )
%    )\label{eq:df-bl.type}\tag{df-bl.type}
%\end{align}
%\end{metadefinition}

%\begin{metadefinition}\blTitle{df-bl.al} at least
%\begin{align}
%\vdash ( A \geq B \leftrightarrow B \le A )\label{eq:df-bl.al}\tag{df-bl.al}
%\end{align}
%\end{metadefinition}

%\begin{metadefinition}\blTitle{df-bl.gt} greater than
%\begin{align}
%\vdash ( A > B \leftrightarrow B < A )\label{eq:df-bl.gt}\tag{df-bl.gt}
%\end{align}
%\end{metadefinition}

%\begin{lemma}\blTitle{bl.npucani} Add/Subtract Same: $a \minus ( a + b ) =
%\textrm{-} b$
%\begin{align}
%\vdash A \in \mathbb{C}\quad\&\quad \vdash B \in
%    \mathbb{C}\quad\Rightarrow\quad \vdash ( A \minus ( A + B ) ) = \textrm{-}
%    B\label{eq:bl.npucani}\tag{bl.npucani}
%\end{align}
%\end{lemma}

%\begin{proof}
%\begin{align}
%  1 &&  & \vdash A \in \mathbb{C}&\text{Hyp~npucani.1}\notag%bl.npucani.1
%  \\ 2 &&  & \vdash A \in \mathbb{C}&\text{Hyp~npucani.1}\notag%bl.npucani.1
%  \\ 3 &&  & \vdash B \in \mathbb{C}&\text{Hyp~npucani.2}\notag%bl.npucani.2
%  \\ 4 &&  & \vdash ( A \in \mathbb{C} \wedge A \in \mathbb{C} \wedge B \in
%       \mathbb{C} )&(\mbox{\ref{eq:3pm3.2i}},1,2,3)\notag
%  \\ 5 &&  & \vdash ( ( A \in \mathbb{C} \wedge A \in \mathbb{C} \wedge B \in
%       \mathbb{C} ) \rightarrow ( ( A \minus A ) \minus B )
%       = ( A \minus ( A + B ) )
%       )&(\mbox{\ref{eq:subsub4}})\notag
%  \\ 6 &&  & \vdash ( ( A \minus A ) \minus B ) = ( A \minus ( A + B )
%       )&(\mbox{\ref{eq:ax-mp}},4,5)\notag
%  \\ 7 &&  & \vdash A \in \mathbb{C}&\text{Hyp~npucani.1}\notag%bl.npucani.1
%  \\ 8 &&  & \vdash ( A \in \mathbb{C} \rightarrow ( A \minus A ) = 0
%       )&(\mbox{\ref{eq:subid}})\notag
%  \\ 9 &&  & \vdash ( A \minus A ) = 0&(\mbox{\ref{eq:ax-mp}},7,8)\notag
%  \\ 10 &&  & \vdash A \in \mathbb{C}&\text{Hyp~npucani.1}\notag%bl.npucani.1
%  \\ 11 &&  & \vdash A \in \mathbb{C}&\text{Hyp~npucani.1}\notag%bl.npucani.1
%  \\ 12 &&  & \vdash ( A \in \mathbb{C} \wedge A \in \mathbb{C}
%       )&(\mbox{\ref{eq:pm3.2i}},10,11)\notag
%  \\ 13 &&  & \vdash ( ( A \in \mathbb{C} \wedge A \in \mathbb{C} ) \rightarrow
%       ( A \minus A ) \in \mathbb{C}
%       )&(\mbox{\ref{eq:subcl}})\notag
%  \\ 14 &&  & \vdash ( A \minus A ) \in
%       \mathbb{C}&(\mbox{\ref{eq:ax-mp}},12,13)\notag
%  \\ 15 &&  & \vdash 0 \in \mathbb{C}&(\mbox{\ref{eq:0cn}})\notag
%  \\ 16 &&  & \vdash B \in \mathbb{C}&\text{Hyp~npucani.2}\notag%bl.npucani.2
%  \\ 17 &&  & \vdash ( ( A \minus A ) \in \mathbb{C} \wedge 0 \in \mathbb{C}
%       \wedge B \in \mathbb{C}
%       )&(\mbox{\ref{eq:3pm3.2i}},14,15,16)\notag
%  \\ 18 &&  & \vdash ( ( ( A \minus A ) \in \mathbb{C} \wedge 0 \in \mathbb{C}
%       \wedge B \in \mathbb{C} ) \rightarrow ( ( ( A \minus
%       A ) \minus B ) = ( 0 \minus B ) \leftrightarrow ( A
%       \minus A ) = 0 ) )&(\mbox{\ref{eq:subcan2}})\notag
%  \\ 19 &&  & \vdash ( ( ( A \minus A ) \minus B ) = ( 0 \minus B )
%       \leftrightarrow ( A \minus A ) = 0
%       )&(\mbox{\ref{eq:ax-mp}},17,18)\notag
%  \\ 20 &&  & \vdash ( ( A \minus A ) \minus B ) = ( 0 \minus B
%       )&(\mbox{\ref{eq:mpbir}},9,19)\notag
%  \\ 21 &&  & \vdash ( A \minus ( A + B ) ) = ( 0 \minus B
%       )&(\mbox{\ref{eq:eqtr3i}},6,20)\notag
%  \\ 22 &&  & \vdash \textrm{-} B = ( 0 \minus B
%       )&(\mbox{\ref{eq:df-neg}})\notag
%  \\ 23 &&  & \vdash ( A \minus ( A + B ) ) = \textrm{-}
%       B&(\mbox{\ref{eq:eqtr4i}},21,22)\notag
%\end{align}
%\end{proof}

\begin{lemma}\blTitle{bl.npucan} Add/Subtract Same: $a \minus ( a + b ) =
\textrm{-} b$
\begin{align}
\vdash ( ( A \in \mathbb{C} \wedge B \in \mathbb{C} ) \rightarrow ( A \minus (
    A + B ) ) = \textrm{-} B )\label{eq:bl.npucan}\tag{bl.npucan}
\end{align}
\end{lemma}

\begin{proof}
\begin{align}
  1 && @1:\  & \vdash ( ( A \in \mathbb{C} \wedge B \in \mathbb{C} ) \rightarrow A \in \mathbb{C} )&(\mbox{\ref{eq:simpl}})\notag
  \\ 2 &&  & \vdash ( ( A \in \mathbb{C} \wedge B \in \mathbb{C} ) \rightarrow A \in \mathbb{C} )&(@1)\notag
  \\ 3 && @3:\  & \vdash ( ( A \in \mathbb{C} \wedge B \in \mathbb{C} ) \rightarrow B \in \mathbb{C} )&(\mbox{\ref{eq:simpr}})\notag
  \\ 4 &&  & \vdash ( ( A \in \mathbb{C} \wedge B \in \mathbb{C} ) \rightarrow ( A \in \mathbb{C} \wedge A \in \mathbb{C} \wedge B \in \mathbb{C} ) )&(\mbox{\ref{eq:3jca}},1,2,3)\notag
  \\ 5 &&  & \vdash ( ( A \in \mathbb{C} \wedge A \in \mathbb{C} \wedge B \in \mathbb{C} ) \rightarrow ( ( A \minus A ) \minus B ) = ( A \minus ( A + B ) ) )&(\mbox{\ref{eq:subsub4}})\notag
  \\ 6 &&  & \vdash ( ( A \in \mathbb{C} \wedge B \in \mathbb{C} ) \rightarrow ( ( A \minus A ) \minus B ) = ( A \minus ( A + B ) ) )&(\mbox{\ref{eq:syl}},4,5)\notag
  \\ 7 &&  & \vdash ( A \in \mathbb{C} \rightarrow ( A \minus A ) = 0 )&(\mbox{\ref{eq:subid}})\notag
  \\ 8 &&  & \vdash ( ( A \in \mathbb{C} \wedge B \in \mathbb{C} ) \rightarrow ( A \minus A ) = 0 )&(\mbox{\ref{eq:adantr}},7)\notag
  \\ 9 &&  & \vdash ( ( A \in \mathbb{C} \wedge B \in \mathbb{C} ) \rightarrow A \in \mathbb{C} )&(@1)\notag
  \\ 10 &&  & \vdash ( ( A \in \mathbb{C} \wedge A \in \mathbb{C} ) \rightarrow ( A \minus A ) \in \mathbb{C} )&(\mbox{\ref{eq:subcl}})\notag
  \\ 11 &&  & \vdash ( ( A \in \mathbb{C} \wedge B \in \mathbb{C} ) \rightarrow ( A \minus A ) \in \mathbb{C} )&(\mbox{\ref{eq:syldan}},9,10)\notag
  \\ 12 &&  & \vdash 0 \in \mathbb{C}&(\mbox{\ref{eq:0cn}})\notag
  \\ 13 &&  & \vdash ( ( A \in \mathbb{C} \wedge B \in \mathbb{C} ) \rightarrow 0 \in \mathbb{C} )&(\mbox{\ref{eq:a1i}},12)\notag
  \\ 14 &&  & \vdash ( ( A \in \mathbb{C} \wedge B \in \mathbb{C} ) \rightarrow B \in \mathbb{C} )&(@3)\notag
  \\ 15 &&  & \vdash ( ( A \in \mathbb{C} \wedge B \in \mathbb{C} ) \rightarrow ( ( A \minus A ) \in \mathbb{C} \wedge 0 \in \mathbb{C} \wedge B \in \mathbb{C} ) )&(\mbox{\ref{eq:3jca}},11,13,14)\notag
  \\ 16 &&  & \vdash ( ( ( A \minus A ) \in \mathbb{C} \wedge 0 \in \mathbb{C} \wedge B \in \mathbb{C} ) \rightarrow ( ( ( A \minus A ) \minus B ) = ( 0 \minus B ) \leftrightarrow ( A \minus A ) = 0 ) )&(\mbox{\ref{eq:subcan2}})\notag
  \\ 17 &&  & \vdash ( ( A \in \mathbb{C} \wedge B \in \mathbb{C} ) \rightarrow ( ( ( A \minus A ) \minus B ) = ( 0 \minus B ) \leftrightarrow ( A \minus A ) = 0 ) )&(\mbox{\ref{eq:syl}},15,16)\notag
  \\ 18 &&  & \vdash ( ( A \in \mathbb{C} \wedge B \in \mathbb{C} ) \rightarrow ( ( A \minus A ) \minus B ) = ( 0 \minus B ) )&(\mbox{\ref{eq:mpbird}},8,17)\notag
  \\ 19 &&  & \vdash ( ( A \in \mathbb{C} \wedge B \in \mathbb{C} ) \rightarrow ( A \minus ( A + B ) ) = ( 0 \minus B ) )&(\mbox{\ref{eq:eqtr3d}},6,18)\notag
  \\ 20 &&  & \vdash \textrm{-} B = ( 0 \minus B )&(\mbox{\ref{eq:df-neg}})\notag
  \\ 21 &&  & \vdash ( ( A \in \mathbb{C} \wedge B \in \mathbb{C} ) \rightarrow ( A \minus ( A + B ) ) = \textrm{-} B )&(\mbox{\ref{eq:syl6eqr}},19,20)\notag
\end{align}
\end{proof}

\begin{lemma}\blTitle{bl.npucan2i} Add/Subtract Same: $b \minus ( a + b )
= \textrm{-} a$
\begin{align}
\vdash A \in \mathbb{C}\quad\&\quad \vdash B \in
    \mathbb{C}\quad\Rightarrow\quad \vdash ( B \minus ( A + B ) ) = \textrm{-}
    A\label{eq:bl.npucan2i}\tag{bl.npucan2i}
\end{align}
\end{lemma}

\begin{proof}
\begin{align}
  1 &&  & \vdash B \in \mathbb{C}&\text{Hyp~npucan2i.2}\notag%bl.npucan2i.2
  \\ 2 &&  & \vdash A \in \mathbb{C}&\text{Hyp~npucan2i.1}\notag%bl.npucan2i.1
  \\ 3 &&  & \vdash B \in \mathbb{C}&\text{Hyp~npucan2i.2}\notag%bl.npucan2i.2
  \\ 4 &&  & \vdash ( A + B ) \in \mathbb{C}&(\mbox{\ref{eq:addcli}},2,3)\notag
  \\ 5 &&  & \vdash ( B \in \mathbb{C} \wedge ( A + B ) \in \mathbb{C} )&(\mbox{\ref{eq:pm3.2i}},1,4)\notag
  \\ 6 &&  & \vdash ( ( B \in \mathbb{C} \wedge ( A + B ) \in \mathbb{C} ) \rightarrow ( B + \textrm{-} ( A + B ) ) = ( B \minus ( A + B ) ) )&(\mbox{\ref{eq:negsub}})\notag
  \\ 7 &&  & \vdash ( B + \textrm{-} ( A + B ) ) = ( B \minus ( A + B ) )&(\mbox{\ref{eq:ax-mp}},5,6)\notag
  \\ 8 &&  & \vdash A \in \mathbb{C}&\text{Hyp~npucan2i.1}\notag%bl.npucan2i.1
  \\ 9 &&  & \vdash B \in \mathbb{C}&\text{Hyp~npucan2i.2}\notag%bl.npucan2i.2
  \\ 10 &&  & \vdash ( A \in \mathbb{C} \wedge B \in \mathbb{C} )&(\mbox{\ref{eq:pm3.2i}},8,9)\notag
  \\ 11 &&  & \vdash ( ( A \in \mathbb{C} \wedge B \in \mathbb{C} ) \rightarrow \textrm{-} ( A + B ) = ( \textrm{-} A + \textrm{-} B ) )&(\mbox{\ref{eq:negdi}})\notag
  \\ 12 &&  & \vdash \textrm{-} ( A + B ) = ( \textrm{-} A + \textrm{-} B )&(\mbox{\ref{eq:ax-mp}},10,11)\notag
  \\ 13 &&  & \vdash ( B + \textrm{-} ( A + B ) ) = ( B + ( \textrm{-} A + \textrm{-} B ) )&(\mbox{\ref{eq:oveq2i}},12)\notag
  \\ 14 &&  & \vdash ( B \minus ( A + B ) ) = ( B + ( \textrm{-} A + \textrm{-} B ) )&(\mbox{\ref{eq:eqtr3i}},7,13)\notag
  \\ 15 &&  & \vdash B \in \mathbb{C}&\text{Hyp~npucan2i.2}\notag%bl.npucan2i.2
  \\ 16 &&  & \vdash A \in \mathbb{C}&\text{Hyp~npucan2i.1}\notag%bl.npucan2i.1
  \\ 17 && @17:\  & \vdash \textrm{-} A \in \mathbb{C}&(\mbox{\ref{eq:negcli}},16)\notag
  \\ 18 &&  & \vdash B \in \mathbb{C}&\text{Hyp~npucan2i.2}\notag%bl.npucan2i.2
  \\ 19 &&  & \vdash \textrm{-} B \in \mathbb{C}&(\mbox{\ref{eq:negcli}},18)\notag
  \\ 20 &&  & \vdash ( B \in \mathbb{C} \wedge \textrm{-} A \in \mathbb{C} \wedge \textrm{-} B \in \mathbb{C} )&(\mbox{\ref{eq:3pm3.2i}},15,17,19)\notag
  \\ 21 &&  & \vdash ( ( B \in \mathbb{C} \wedge \textrm{-} A \in \mathbb{C} \wedge \textrm{-} B \in \mathbb{C} ) \rightarrow ( B + ( \textrm{-} A + \textrm{-} B ) ) = ( \textrm{-} A + ( B + \textrm{-} B ) ) )&(\mbox{\ref{eq:add12}})\notag
  \\ 22 &&  & \vdash ( B + ( \textrm{-} A + \textrm{-} B ) ) = ( \textrm{-} A + ( B + \textrm{-} B ) )&(\mbox{\ref{eq:ax-mp}},20,21)\notag
  \\ 23 &&  & \vdash ( B \minus ( A + B ) ) = ( \textrm{-} A + ( B + \textrm{-} B ) )&(\mbox{\ref{eq:eqtri}},14,22)\notag
  \\ 24 &&  & \vdash B \in \mathbb{C}&\text{Hyp~npucan2i.2}\notag%bl.npucan2i.2
  \\ 25 &&  & \vdash ( B \in \mathbb{C} \rightarrow ( B + \textrm{-} B ) = 0 )&(\mbox{\ref{eq:negid}})\notag
  \\ 26 &&  & \vdash ( B + \textrm{-} B ) = 0&(\mbox{\ref{eq:ax-mp}},24,25)\notag
  \\ 27 &&  & \vdash ( \textrm{-} A + ( B + \textrm{-} B ) ) = ( \textrm{-} A + 0 )&(\mbox{\ref{eq:oveq2i}},26)\notag
  \\ 28 &&  & \vdash \textrm{-} A \in \mathbb{C}&(@17)\notag
  \\ 29 &&  & \vdash ( \textrm{-} A \in \mathbb{C} \rightarrow ( \textrm{-} A + 0 ) = \textrm{-} A )&(\mbox{\ref{eq:addid1}})\notag
  \\ 30 &&  & \vdash ( \textrm{-} A + 0 ) = \textrm{-} A&(\mbox{\ref{eq:ax-mp}},28,29)\notag
  \\ 31 &&  & \vdash ( B \minus ( A + B ) ) = \textrm{-} A&(\mbox{\ref{eq:3eqtri}},23,27,30)\notag
\end{align}
\end{proof}

\begin{lemma}\blTitle{bl.npucan2d} Add/Subtract Same: $b \minus ( a + b )
= \textrm{-} a$ deduction form
\begin{align}
\vdash ( \mw{\varphi} \rightarrow A \in \mathbb{C} )\quad\&\quad \vdash (
    \mw{\varphi} \rightarrow B \in \mathbb{C} )\quad\Rightarrow\quad \vdash (
    \mw{\varphi} \rightarrow ( B \minus ( A + B ) ) = \textrm{-} A
    )\label{eq:bl.npucan2d}\tag{bl.npucan2d}
\end{align}
\end{lemma}

\begin{proof}
\begin{align}
  1 &&  & &(\mbox{\ref{eq:?}})\notag
\end{align}
\end{proof}

\begin{lemma}\blTitle{bl.npucan2} Add/Subtract Same: $b \minus ( a + b ) =
\textrm{-} a$
\begin{align}
\vdash ( ( A \in \mathbb{C} \wedge B \in \mathbb{C} ) \rightarrow ( B \minus (
    A + B ) ) = \textrm{-} A )\label{eq:bl.npucan2}\tag{bl.npucan2}
\end{align}
\end{lemma}

\begin{proof}
\begin{align}
  1 &&  & &(\mbox{\ref{eq:?}})\notag
\end{align}
\end{proof}

\begin{lemma}\blTitle{bl.subadd4} Associativity: $( a \minus b ) + ( c
\minus d ) = ( a + c ) \minus ( b + d )$
\begin{align}
\vdash ( ( ( A \in \mathbb{C} \wedge B \in \mathbb{C} ) \wedge ( C \in
    \mathbb{C} \wedge D \in \mathbb{C} ) ) \rightarrow ( ( A \minus B ) + ( C
    \minus D ) ) = ( ( A + C ) \minus ( B + D ) )
    )\label{eq:bl.subadd4}\tag{bl.subadd4}
\end{align}
\end{lemma}

\begin{proof}
\begin{align}
  1 &&  & \vdash ( ( ( A \in \mathbb{C} \wedge C \in \mathbb{C} ) \wedge ( B \in \mathbb{C} \wedge D \in \mathbb{C} ) ) \rightarrow ( ( A + C ) \minus ( B + D ) ) = ( ( A \minus B ) + ( C \minus D ) ) )&(\mbox{\ref{eq:addsub4}})\notag
  \\ 2 &&  & \vdash ( ( ( A \in \mathbb{C} \wedge C \in \mathbb{C} ) \wedge ( B \in \mathbb{C} \wedge D \in \mathbb{C} ) ) \leftrightarrow ( A \in \mathbb{C} \wedge ( C \in \mathbb{C} \wedge ( B \in \mathbb{C} \wedge D \in \mathbb{C} ) ) ) )&(\mbox{\ref{eq:anass}})\notag
  \\ 3 &&  & \vdash ( ( B \in \mathbb{C} \wedge D \in \mathbb{C} ) \leftrightarrow ( D \in \mathbb{C} \wedge B \in \mathbb{C} ) )&(\mbox{\ref{eq:ancom}})\notag
  \\ 4 &&  & \vdash ( ( C \in \mathbb{C} \wedge ( B \in \mathbb{C} \wedge D \in \mathbb{C} ) ) \leftrightarrow ( C \in \mathbb{C} \wedge ( D \in \mathbb{C} \wedge B \in \mathbb{C} ) ) )&(\mbox{\ref{eq:anbi2i}},3)\notag
  \\ 5 &&  & \vdash ( ( ( C \in \mathbb{C} \wedge D \in \mathbb{C} ) \wedge B \in \mathbb{C} ) \leftrightarrow ( C \in \mathbb{C} \wedge ( D \in \mathbb{C} \wedge B \in \mathbb{C} ) ) )&(\mbox{\ref{eq:anass}})\notag
  \\ 6 &&  & \vdash ( ( C \in \mathbb{C} \wedge ( D \in \mathbb{C} \wedge B \in \mathbb{C} ) ) \leftrightarrow ( ( C \in \mathbb{C} \wedge D \in \mathbb{C} ) \wedge B \in \mathbb{C} ) )&(\mbox{\ref{eq:bicomi}},5)\notag
  \\ 7 &&  & \vdash ( ( C \in \mathbb{C} \wedge ( B \in \mathbb{C} \wedge D \in \mathbb{C} ) ) \leftrightarrow ( ( C \in \mathbb{C} \wedge D \in \mathbb{C} ) \wedge B \in \mathbb{C} ) )&(\mbox{\ref{eq:bitri}},4,6)\notag
  \\ 8 &&  & \vdash ( ( A \in \mathbb{C} \wedge ( C \in \mathbb{C} \wedge ( B \in \mathbb{C} \wedge D \in \mathbb{C} ) ) ) \leftrightarrow ( A \in \mathbb{C} \wedge ( ( C \in \mathbb{C} \wedge D \in \mathbb{C} ) \wedge B \in \mathbb{C} ) ) )&(\mbox{\ref{eq:anbi2i}},7)\notag
  \\ 9 &&  & \vdash ( ( ( C \in \mathbb{C} \wedge D \in \mathbb{C} ) \wedge B \in \mathbb{C} ) \leftrightarrow ( B \in \mathbb{C} \wedge ( C \in \mathbb{C} \wedge D \in \mathbb{C} ) ) )&(\mbox{\ref{eq:ancom}})\notag
  \\ 10 &&  & \vdash ( ( A \in \mathbb{C} \wedge ( ( C \in \mathbb{C} \wedge D \in \mathbb{C} ) \wedge B \in \mathbb{C} ) ) \leftrightarrow ( A \in \mathbb{C} \wedge ( B \in \mathbb{C} \wedge ( C \in \mathbb{C} \wedge D \in \mathbb{C} ) ) ) )&(\mbox{\ref{eq:anbi2i}},9)\notag
  \\ 11 &&  & \vdash ( ( A \in \mathbb{C} \wedge ( C \in \mathbb{C} \wedge ( B \in \mathbb{C} \wedge D \in \mathbb{C} ) ) ) \leftrightarrow ( A \in \mathbb{C} \wedge ( B \in \mathbb{C} \wedge ( C \in \mathbb{C} \wedge D \in \mathbb{C} ) ) ) )&(\mbox{\ref{eq:bitri}},8,10)\notag
  \\ 12 &&  & \vdash ( ( ( A \in \mathbb{C} \wedge B \in \mathbb{C} ) \wedge ( C \in \mathbb{C} \wedge D \in \mathbb{C} ) ) \leftrightarrow ( A \in \mathbb{C} \wedge ( B \in \mathbb{C} \wedge ( C \in \mathbb{C} \wedge D \in \mathbb{C} ) ) ) )&(\mbox{\ref{eq:anass}})\notag
  \\ 13 &&  & \vdash ( ( A \in \mathbb{C} \wedge ( B \in \mathbb{C} \wedge ( C \in \mathbb{C} \wedge D \in \mathbb{C} ) ) ) \leftrightarrow ( ( A \in \mathbb{C} \wedge B \in \mathbb{C} ) \wedge ( C \in \mathbb{C} \wedge D \in \mathbb{C} ) ) )&(\mbox{\ref{eq:bicomi}},12)\notag
  \\ 14 &&  & \vdash ( ( A \in \mathbb{C} \wedge ( C \in \mathbb{C} \wedge ( B \in \mathbb{C} \wedge D \in \mathbb{C} ) ) ) \leftrightarrow ( ( A \in \mathbb{C} \wedge B \in \mathbb{C} ) \wedge ( C \in \mathbb{C} \wedge D \in \mathbb{C} ) ) )&(\mbox{\ref{eq:bitri}},11,13)\notag
  \\ 15 &&  & \vdash ( ( ( A \in \mathbb{C} \wedge C \in \mathbb{C} ) \wedge ( B \in \mathbb{C} \wedge D \in \mathbb{C} ) ) \leftrightarrow ( ( A \in \mathbb{C} \wedge B \in \mathbb{C} ) \wedge ( C \in \mathbb{C} \wedge D \in \mathbb{C} ) ) )&(\mbox{\ref{eq:bitri}},2,14)\notag
  \\ 16 &&  & \vdash ( ( ( ( A \in \mathbb{C} \wedge C \in \mathbb{C} ) \wedge ( B \in \mathbb{C} \wedge D \in \mathbb{C} ) ) \rightarrow ( ( A + C ) \minus ( B + D ) ) = ( ( A \minus B ) + ( C \minus D ) ) ) \leftrightarrow ( ( ( A \in \mathbb{C} \wedge B \in \mathbb{C} ) \wedge ( C \in \mathbb{C} \wedge D \in \mathbb{C} ) ) \rightarrow ( ( A + C ) \minus ( B + D ) ) = ( ( A \minus B ) + ( C \minus D ) ) ) )&(\mbox{\ref{eq:imbi1i}},15)\notag
  \\ 17 &&  & \vdash ( ( ( A + C ) \minus ( B + D ) ) = ( ( A \minus B ) + ( C \minus D ) ) \leftrightarrow ( ( A \minus B ) + ( C \minus D ) ) = ( ( A + C ) \minus ( B + D ) ) )&(\mbox{\ref{eq:eqcom}})\notag
  \\ 18 &&  & \vdash ( ( ( ( A \in \mathbb{C} \wedge B \in \mathbb{C} ) \wedge ( C \in \mathbb{C} \wedge D \in \mathbb{C} ) ) \rightarrow ( ( A + C ) \minus ( B + D ) ) = ( ( A \minus B ) + ( C \minus D ) ) ) \leftrightarrow ( ( ( A \in \mathbb{C} \wedge B \in \mathbb{C} ) \wedge ( C \in \mathbb{C} \wedge D \in \mathbb{C} ) ) \rightarrow ( ( A \minus B ) + ( C \minus D ) ) = ( ( A + C ) \minus ( B + D ) ) ) )&(\mbox{\ref{eq:imbi2i}},17)\notag
  \\ 19 &&  & \vdash ( ( ( ( A \in \mathbb{C} \wedge C \in \mathbb{C} ) \wedge ( B \in \mathbb{C} \wedge D \in \mathbb{C} ) ) \rightarrow ( ( A + C ) \minus ( B + D ) ) = ( ( A \minus B ) + ( C \minus D ) ) ) \leftrightarrow ( ( ( A \in \mathbb{C} \wedge B \in \mathbb{C} ) \wedge ( C \in \mathbb{C} \wedge D \in \mathbb{C} ) ) \rightarrow ( ( A \minus B ) + ( C \minus D ) ) = ( ( A + C ) \minus ( B + D ) ) ) )&(\mbox{\ref{eq:bitri}},16,18)\notag
  \\ 20 &&  & \vdash ( ( ( A \in \mathbb{C} \wedge B \in \mathbb{C} ) \wedge ( C \in \mathbb{C} \wedge D \in \mathbb{C} ) ) \rightarrow ( ( A \minus B ) + ( C \minus D ) ) = ( ( A + C ) \minus ( B + D ) ) )&(\mbox{\ref{eq:mpbi}},1,19)\notag
\end{align}
\end{proof}

\begin{lemma}\blTitle{bl.rnegsub} Unary Minus: $( \textrm{-} a + b ) = ( b
\minus a )$
\begin{align}
\vdash ( ( A \in \mathbb{C} \wedge B \in \mathbb{C} ) \rightarrow ( \textrm{-}
    A + B ) = ( B \minus A ) )\label{eq:bl.rnegsub}\tag{bl.rnegsub}
\end{align}
\end{lemma}

\begin{proof}
\begin{align}
  1 &&  & \vdash ( ( B \in \mathbb{C} \wedge A \in \mathbb{C} ) \rightarrow ( B + \textrm{-} A ) = ( B \minus A ) )&(\mbox{\ref{eq:negsub}})\notag
  \\ 2 &&  & \vdash ( ( B \in \mathbb{C} \wedge A \in \mathbb{C} ) \leftrightarrow ( A \in \mathbb{C} \wedge B \in \mathbb{C} ) )&(\mbox{\ref{eq:ancom}})\notag
  \\ 3 &&  & \vdash ( ( ( B \in \mathbb{C} \wedge A \in \mathbb{C} ) \rightarrow ( B + \textrm{-} A ) = ( B \minus A ) ) \leftrightarrow ( ( A \in \mathbb{C} \wedge B \in \mathbb{C} ) \rightarrow ( B + \textrm{-} A ) = ( B \minus A ) ) )&(\mbox{\ref{eq:imbi1i}},2)\notag
  \\ 4 &&  & \vdash ( B + \textrm{-} A ) = ( \textrm{-} A + B )&(\mbox{\ref{eq:addcomgi}})\notag
  \\ 5 &&  & \vdash ( ( B + \textrm{-} A ) = ( B \minus A ) \leftrightarrow ( \textrm{-} A + B ) = ( B \minus A ) )&(\mbox{\ref{eq:eqeq1i}},4)\notag
  \\ 6 &&  & \vdash ( ( ( A \in \mathbb{C} \wedge B \in \mathbb{C} ) \rightarrow ( B + \textrm{-} A ) = ( B \minus A ) ) \leftrightarrow ( ( A \in \mathbb{C} \wedge B \in \mathbb{C} ) \rightarrow ( \textrm{-} A + B ) = ( B \minus A ) ) )&(\mbox{\ref{eq:imbi2i}},5)\notag
  \\ 7 &&  & \vdash ( ( ( B \in \mathbb{C} \wedge A \in \mathbb{C} ) \rightarrow ( B + \textrm{-} A ) = ( B \minus A ) ) \leftrightarrow ( ( A \in \mathbb{C} \wedge B \in \mathbb{C} ) \rightarrow ( \textrm{-} A + B ) = ( B \minus A ) ) )&(\mbox{\ref{eq:bitri}},3,6)\notag
  \\ 8 &&  & \vdash ( ( A \in \mathbb{C} \wedge B \in \mathbb{C} ) \rightarrow ( \textrm{-} A + B ) = ( B \minus A ) )&(\mbox{\ref{eq:mpbi}},1,7)\notag
\end{align}
\end{proof}

\begin{lemma}\blTitle{bl.uaddsub} Unary Minus: $( \textrm{-} a + b )
\minus c = b \minus ( a + c )$
\begin{align}
\vdash ( ( A \in \mathbb{C} \wedge B \in \mathbb{C} \wedge C \in \mathbb{C} )
    \rightarrow ( ( \textrm{-} A + B ) \minus C ) = ( B \minus ( A + C ) )
    )\label{eq:bl.uaddsub}\tag{bl.uaddsub}
\end{align}
\end{lemma}

\begin{proof}
\begin{align}
  1 &&  & \vdash ( ( A \in \mathbb{C} \wedge B \in \mathbb{C} ) \rightarrow ( \textrm{-} A + B ) = ( B \minus A ) )&(\mbox{\ref{eq:bl.rnegsub}})\notag
  \\ 2 &&  & \vdash ( ( A \in \mathbb{C} \wedge B \in \mathbb{C} \wedge C \in \mathbb{C} ) \rightarrow ( \textrm{-} A + B ) = ( B \minus A ) )&(\mbox{\ref{eq:3adant3}},1)\notag
  \\ 3 && @3:\  & \vdash ( ( A \in \mathbb{C} \wedge B \in \mathbb{C} \wedge C \in \mathbb{C} ) \rightarrow A \in \mathbb{C} )&(\mbox{\ref{eq:simp1}})\notag
  \\ 4 &&  & \vdash ( A \in \mathbb{C} \rightarrow \textrm{-} A \in \mathbb{C} )&(\mbox{\ref{eq:negcl}})\notag
  \\ 5 &&  & \vdash ( ( A \in \mathbb{C} \wedge B \in \mathbb{C} \wedge C \in \mathbb{C} ) \rightarrow \textrm{-} A \in \mathbb{C} )&(\mbox{\ref{eq:syl}},3,4)\notag
  \\ 6 && @6:\  & \vdash ( ( A \in \mathbb{C} \wedge B \in \mathbb{C} \wedge C \in \mathbb{C} ) \rightarrow B \in \mathbb{C} )&(\mbox{\ref{eq:simp2}})\notag
  \\ 7 &&  & \vdash ( ( \textrm{-} A \in \mathbb{C} \wedge B \in \mathbb{C} ) \rightarrow ( \textrm{-} A + B ) \in \mathbb{C} )&(\mbox{\ref{eq:addcl}})\notag
  \\ 8 &&  & \vdash ( ( A \in \mathbb{C} \wedge B \in \mathbb{C} \wedge C \in \mathbb{C} ) \rightarrow ( \textrm{-} A + B ) \in \mathbb{C} )&(\mbox{\ref{eq:syl2anc}},5,6,7)\notag
  \\ 9 &&  & \vdash ( ( A \in \mathbb{C} \wedge B \in \mathbb{C} \wedge C \in \mathbb{C} ) \rightarrow B \in \mathbb{C} )&(@6)\notag
  \\ 10 &&  & \vdash ( ( A \in \mathbb{C} \wedge B \in \mathbb{C} \wedge C \in \mathbb{C} ) \rightarrow A \in \mathbb{C} )&(@3)\notag
  \\ 11 &&  & \vdash ( ( A \in \mathbb{C} \wedge B \in \mathbb{C} \wedge C \in \mathbb{C} ) \rightarrow ( B \in \mathbb{C} \wedge A \in \mathbb{C} ) )&(\mbox{\ref{eq:jca}},9,10)\notag
  \\ 12 &&  & \vdash ( ( B \in \mathbb{C} \wedge A \in \mathbb{C} ) \rightarrow ( B \minus A ) \in \mathbb{C} )&(\mbox{\ref{eq:subcl}})\notag
  \\ 13 &&  & \vdash ( ( A \in \mathbb{C} \wedge B \in \mathbb{C} \wedge C \in \mathbb{C} ) \rightarrow ( B \minus A ) \in \mathbb{C} )&(\mbox{\ref{eq:syl}},11,12)\notag
  \\ 14 && @14:\  & \vdash ( ( A \in \mathbb{C} \wedge B \in \mathbb{C} \wedge C \in \mathbb{C} ) \rightarrow C \in \mathbb{C} )&(\mbox{\ref{eq:simp3}})\notag
  \\ 15 &&  & \vdash ( ( A \in \mathbb{C} \wedge B \in \mathbb{C} \wedge C \in \mathbb{C} ) \rightarrow ( ( \textrm{-} A + B ) \in \mathbb{C} \wedge ( B \minus A ) \in \mathbb{C} \wedge C \in \mathbb{C} ) )&(\mbox{\ref{eq:3jca}},8,13,14)\notag
  \\ 16 &&  & \vdash ( ( ( \textrm{-} A + B ) \in \mathbb{C} \wedge ( B \minus A ) \in \mathbb{C} \wedge C \in \mathbb{C} ) \rightarrow ( ( ( \textrm{-} A + B ) \minus C ) = ( ( B \minus A ) \minus C ) \leftrightarrow ( \textrm{-} A + B ) = ( B \minus A ) ) )&(\mbox{\ref{eq:subcan2}})\notag
  \\ 17 &&  & \vdash ( ( A \in \mathbb{C} \wedge B \in \mathbb{C} \wedge C \in \mathbb{C} ) \rightarrow ( ( ( \textrm{-} A + B ) \minus C ) = ( ( B \minus A ) \minus C ) \leftrightarrow ( \textrm{-} A + B ) = ( B \minus A ) ) )&(\mbox{\ref{eq:syl}},15,16)\notag
  \\ 18 &&  & \vdash ( ( A \in \mathbb{C} \wedge B \in \mathbb{C} \wedge C \in \mathbb{C} ) \rightarrow ( ( \textrm{-} A + B ) \minus C ) = ( ( B \minus A ) \minus C ) )&(\mbox{\ref{eq:mpbird}},2,17)\notag
  \\ 19 &&  & \vdash ( ( A \in \mathbb{C} \wedge B \in \mathbb{C} \wedge C \in \mathbb{C} ) \rightarrow B \in \mathbb{C} )&(@6)\notag
  \\ 20 &&  & \vdash ( ( A \in \mathbb{C} \wedge B \in \mathbb{C} \wedge C \in \mathbb{C} ) \rightarrow A \in \mathbb{C} )&(@3)\notag
  \\ 21 &&  & \vdash ( ( A \in \mathbb{C} \wedge B \in \mathbb{C} \wedge C \in \mathbb{C} ) \rightarrow C \in \mathbb{C} )&(@14)\notag
  \\ 22 &&  & \vdash ( ( A \in \mathbb{C} \wedge B \in \mathbb{C} \wedge C \in \mathbb{C} ) \rightarrow ( B \in \mathbb{C} \wedge A \in \mathbb{C} \wedge C \in \mathbb{C} ) )&(\mbox{\ref{eq:3jca}},19,20,21)\notag
  \\ 23 &&  & \vdash ( ( B \in \mathbb{C} \wedge A \in \mathbb{C} \wedge C \in \mathbb{C} ) \rightarrow ( ( B \minus A ) \minus C ) = ( B \minus ( A + C ) ) )&(\mbox{\ref{eq:subsub4}})\notag
  \\ 24 &&  & \vdash ( ( A \in \mathbb{C} \wedge B \in \mathbb{C} \wedge C \in \mathbb{C} ) \rightarrow ( ( B \minus A ) \minus C ) = ( B \minus ( A + C ) ) )&(\mbox{\ref{eq:syl}},22,23)\notag
  \\ 25 &&  & \vdash ( ( A \in \mathbb{C} \wedge B \in \mathbb{C} \wedge C \in \mathbb{C} ) \rightarrow ( ( \textrm{-} A + B ) \minus C ) = ( B \minus ( A + C ) ) )&(\mbox{\ref{eq:eqtrd}},18,24)\notag
\end{align}
\end{proof}

\begin{lemma}\blTitle{bl.addusub} Unary Minus: $( a + \textrm{-} b )
\minus c = a \minus ( b + c )$
\begin{align}
\vdash ( ( A \in \mathbb{C} \wedge B \in \mathbb{C} \wedge C \in \mathbb{C} )
    \rightarrow ( ( A + \textrm{-} B ) \minus C ) = ( A \minus ( B + C ) )
    )\label{eq:bl.addusub}\tag{bl.addusub}
\end{align}
\end{lemma}

\begin{proof}
\begin{align}
  1 &&  & &(\mbox{\ref{eq:?}})\notag
\end{align}
\end{proof}

\begin{lemma}\blTitle{bl.erla} Empty Range Law: all a:t in false are P is
tautology
\begin{align}
\vdash \forall x \in \varnothing \mw{\varphi}\label{eq:bl.erla}\tag{bl.erla}
\end{align}
\end{lemma}

\begin{proof}
\begin{align}
  1 &&  & &(\mbox{\ref{eq:?}})\notag
\end{align}
\end{proof}

\begin{lemma}\blTitle{bl.erle} Empty Range Law: exists a:t in false that P
is false
\begin{align}
\vdash \lnot \exists x \in \varnothing
    \mw{\varphi}\label{eq:bl.erle}\tag{bl.erle}
\end{align}
\end{lemma}

\begin{proof}
\begin{align}
  1 &&  & &(\mbox{\ref{eq:?}})\notag
\end{align}
\end{proof}

\begin{lemma}\blTitle{bl.erls} Empty Range Law: (sum a:t in false of P) =
0
\begin{align}
\vdash ( ( \Sigma k \in \varnothing A ) = 0 )\label{eq:bl.erls}\tag{bl.erls}
\end{align}
\end{lemma}

\begin{proof}
\begin{align}
  1 &&  & &(\mbox{\ref{eq:?}})\notag
\end{align}
\end{proof}

\begin{lemma}\blTitle{bl.srla} Solitary Range Law: all a:t in j..j are P
is P[j/a]
\begin{align}
\vdash ( \forall x \in A [,] A \mw{\varphi} \leftrightarrow [ A / x ]
    \mw{\varphi} )\label{eq:bl.srla}\tag{bl.srla}
\end{align}
\end{lemma}

\begin{proof}
\begin{align}
  1 &&  & &(\mbox{\ref{eq:?}})\notag
\end{align}
\end{proof}

\begin{lemma}\blTitle{bl.srle} Solitary Range Law: exists a:t in j..j that
P is P[j/a]
\begin{align}
\vdash ( \exists x \in A [,] A \mw{\varphi} \leftrightarrow [ A / x ]
    \mw{\varphi} )\label{eq:bl.srle}\tag{bl.srle}
\end{align}
\end{lemma}

\begin{proof}
\begin{align}
  1 &&  & &(\mbox{\ref{eq:?}})\notag
\end{align}
\end{proof}

\begin{lemma}\blTitle{bl.srls} Solitary Range Law: (sum a:t in j..j of P)
= P[j/a]
\begin{align}
\vdash ( \Sigma x \in A [,] A \mw{\varphi} \leftrightarrow [ A / x ]
    \mw{\varphi} )\label{eq:bl.srls}\tag{bl.srls}
\end{align}
\end{lemma}

\begin{proof}
\begin{align}
  1 &&  & &(\mbox{\ref{eq:?}})\notag
\end{align}
\end{proof}

\begin{lemma}\blTitle{bl.arla} Assumed Range Law: (all a:t which Q are P)
iff (all a:t are (Q and P))
\begin{align}
\vdash ( \forall x \in \{ \mw{\psi} \} \mw{\varphi} \leftrightarrow \forall x
    \in ( \mw{\psi} \wedge \mw{\varphi} ) )\label{eq:bl.arla}\tag{bl.arla}
\end{align}
\end{lemma}

\begin{proof}
\begin{align}
  1 &&  & &(\mbox{\ref{eq:?}})\notag
\end{align}
\end{proof}

\begin{lemma}\blTitle{bl.arle} Assumed Range Law: (exists a:t which Q that
P) iff (exists a:t that (Q and P))
\begin{align}
\vdash ( \exists x \in \{ \mw{\psi} \} \mw{\varphi} \leftrightarrow \exists x
    \in ( \mw{\psi} \wedge \mw{\varphi} ) )\label{eq:bl.arle}\tag{bl.arle}
\end{align}
\end{lemma}

\begin{proof}
\begin{align}
  1 &&  & &(\mbox{\ref{eq:?}})\notag
\end{align}
\end{proof}

\begin{lemma}\blTitle{bl.tbla} True Body Law: all a:t in R are true is
tautology
\begin{align}
\vdash \forall x \in A \top\label{eq:bl.tbla}\tag{bl.tbla}
\end{align}
\end{lemma}

\begin{proof}
\begin{align}
  1 &&  & &(\mbox{\ref{eq:?}})\notag
\end{align}
\end{proof}

\begin{lemma}\blTitle{bl.fbla} False Body Law:  all a:t in R are false is
false
\begin{align}
\vdash \lnot \forall x \in A \bot\label{eq:bl.fbla}\tag{bl.fbla}
\end{align}
\end{lemma}

\begin{proof}
\begin{align}
  1 &&  & &(\mbox{\ref{eq:?}})\notag
\end{align}
\end{proof}

\begin{lemma}\blTitle{bl.fble} False Body Law:  exists a:t in R that false
is false
\begin{align}
\vdash \lnot \exists x \in A \bot\label{eq:bl.fble}\tag{bl.fble}
\end{align}
\end{lemma}

\begin{proof}
\begin{align}
  1 &&  & &(\mbox{\ref{eq:?}})\notag
\end{align}
\end{proof}

\begin{lemma}\blTitle{bl.solrla} Solitary Open Left Range Law: exists a:t
in j,.j that P is false
\begin{align}
\vdash \lnot \exists x \in A (,] A
    \mw{\varphi}\label{eq:bl.solrla}\tag{bl.solrla}
\end{align}
\end{lemma}

\begin{proof}
\begin{align}
  1 &&  & &(\mbox{\ref{eq:?}})\notag
\end{align}
\end{proof}

\begin{lemma}\blTitle{bl.sorrla} Solitary Open Right Range Law: exists a:t
in j.,j that P is false
\begin{align}
\vdash \lnot \exists x \in A [,) A
    \mw{\varphi}\label{eq:bl.sorrla}\tag{bl.sorrla}
\end{align}
\end{lemma}

\begin{proof}
\begin{align}
  1 &&  & &(\mbox{\ref{eq:?}})\notag
\end{align}
\end{proof}

\begin{lemma}\blTitle{bl.sorla} Solitary Open Range Law: exists a:t in
j,,j that P is false
\begin{align}
\vdash \lnot \exists x \in A (,) A
    \mw{\varphi}\label{eq:bl.sorla}\tag{bl.sorla}
\end{align}
\end{lemma}

\begin{proof}
\begin{align}
  1 &&  & &(\mbox{\ref{eq:?}})\notag
\end{align}
\end{proof}

\begin{lemma}\blTitle{bl.ancomphfirst} make any wff first in conjunction
wff-list
\begin{align}
\vdash ( \bigwedge( \mwl{l_1} \mw{\varphi} \mwl{l_2} ) \leftrightarrow
    \bigwedge( \mw{\varphi} \mwl{l_2} \mwl{l_1} )
    )\label{eq:bl.ancomphfirst}\tag{bl.ancomphfirst}
\end{align}
\end{lemma}

\begin{proof}
\begin{align}
  1 &&  & \vdash ( \bigwedge( \mwl{l_1} \mw{\varphi} \mwl{l_2} ) \leftrightarrow \bigwedge( \mw{\varphi} \mwl{l_2} \mwl{l_1} ) )&(\mbox{\ref{eq:bl.ancomwlwl}})\notag
\end{align}
\end{proof}

\begin{lemma}\blTitle{bl.ancomphlast} make any wff last in conjunction
wff-list
\begin{align}
\vdash ( \bigwedge( \mwl{l_1} \mw{\varphi} \mwl{l_2} ) \leftrightarrow
    \bigwedge( \mwl{l_2} \mwl{l_1} \mw{\varphi} )
    )\label{eq:bl.ancomphlast}\tag{bl.ancomphlast}
\end{align}
\end{lemma}

\begin{proof}
\begin{align}
  1 &&  & \vdash ( \bigwedge( \mwl{l_1} \mw{\varphi} \mwl{l_2} ) \leftrightarrow \bigwedge( \mwl{l_2} \mwl{l_1} \mw{\varphi} ) )&(\mbox{\ref{eq:bl.ancomwlwl}})\notag
\end{align}
\end{proof}

\begin{lemma}\blTitle{bl.anpmw} pull middle wff from conjunction wff-list
\begin{align}
\vdash ( \bigwedge( \mwl{l_1} \mw{\varphi} \mwl{l_2} ) \leftrightarrow (
    \mw{\varphi} \wedge \bigwedge( \mwl{l_2} \mwl{l_1} ) )
    )\label{eq:bl.anpmw}\tag{bl.anpmw}
\end{align}
\end{lemma}

\begin{proof}
\begin{align}
  1 &&  & \vdash ( \bigwedge( \mwl{l_1} \mw{\varphi} \mwl{l_2} ) \leftrightarrow \bigwedge( \mw{\varphi} \mwl{l_2} \mwl{l_1} ) )&(\mbox{\ref{eq:bl.ancomphfirst}})\notag
  \\ 2 &&  & \vdash ( \bigwedge( \mw{\varphi} \mwl{l_2} \mwl{l_1} ) \leftrightarrow ( \mw{\varphi} \wedge \bigwedge( \mwl{l_2} \mwl{l_1} ) ) )&(\mbox{\ref{eq:bl.anpfw}})\notag
  \\ 3 &&  & \vdash ( \bigwedge( \mwl{l_1} \mw{\varphi} \mwl{l_2} ) \leftrightarrow ( \mw{\varphi} \wedge \bigwedge( \mwl{l_2} \mwl{l_1} ) ) )&(\mbox{\ref{eq:bitri}},1,2)\notag
\end{align}
\end{proof}

\begin{lemma}\blTitle{bl.anabpm} absorb parentheses, middle term
\begin{align}
\vdash ( \bigwedge( \mwl{l_1} \bigwedge( \mwl{l_2} ) \mwl{l_3} )
    \leftrightarrow \bigwedge( \mwl{l_1} \mwl{l_2} \mwl{l_3} )
    )\label{eq:bl.anabpm}\tag{bl.anabpm}
\end{align}
\end{lemma}

\begin{proof}
\begin{align}
  1 &&  & \vdash ( \bigwedge( \mwl{l_1} \bigwedge( \mwl{l_2} ) \mwl{l_3} ) \leftrightarrow \bigwedge( \bigwedge( \mwl{l_2} ) \mwl{l_3} \mwl{l_1} ) )&(\mbox{\ref{eq:bl.ancomwlwl}})\notag
  \\ 2 &&  & \vdash ( \bigwedge( \bigwedge( \mwl{l_2} ) \mwl{l_3} \mwl{l_1} ) \leftrightarrow \bigwedge( \mwl{l_2} \mwl{l_3} \mwl{l_1} ) )&(\mbox{\ref{eq:bl.anabpf}})\notag
  \\ 3 &&  & \vdash ( \bigwedge( \mwl{l_2} \mwl{l_3} \mwl{l_1} ) \leftrightarrow \bigwedge( \mwl{l_1} \mwl{l_2} \mwl{l_3} ) )&(\mbox{\ref{eq:bl.ancomwlwl}})\notag
  \\ 4 &&  & \vdash ( \bigwedge( \mwl{l_1} \bigwedge( \mwl{l_2} ) \mwl{l_3} ) \leftrightarrow \bigwedge( \mwl{l_1} \mwl{l_2} \mwl{l_3} ) )&(\mbox{\ref{eq:3bitri}},1,2,3)\notag
\end{align}
\end{proof}

\begin{lemma}\blTitle{bl.orcomphfirst} make any wff first in disjunction
wff-list
\begin{align}
\vdash ( \bigvee( \mwl{l_1} \mw{\varphi} \mwl{l_2} ) \leftrightarrow \bigvee(
    \mw{\varphi} \mwl{l_2} \mwl{l_1} )
    )\label{eq:bl.orcomphfirst}\tag{bl.orcomphfirst}
\end{align}
\end{lemma}

\begin{proof}
\begin{align}
  1 &&  & \vdash ( \bigvee( \mwl{l_1} \mw{\varphi} \mwl{l_2} ) \leftrightarrow \bigvee( \mw{\varphi} \mwl{l_2} \mwl{l_1} ) )&(\mbox{\ref{eq:bl.orcomwlwl}})\notag
\end{align}
\end{proof}

\begin{lemma}\blTitle{bl.orcomphlast} make any wff last in disjunction
wff-list
\begin{align}
\vdash ( \bigvee( \mwl{l_1} \mw{\varphi} \mwl{l_2} ) \leftrightarrow \bigvee(
    \mwl{l_2} \mwl{l_1} \mw{\varphi} )
    )\label{eq:bl.orcomphlast}\tag{bl.orcomphlast}
\end{align}
\end{lemma}

\begin{proof}
\begin{align}
  1 &&  & \vdash ( \bigvee( \mwl{l_1} \mw{\varphi} \mwl{l_2} ) \leftrightarrow \bigvee( \mwl{l_2} \mwl{l_1} \mw{\varphi} ) )&(\mbox{\ref{eq:bl.orcomwlwl}})\notag
\end{align}
\end{proof}

\begin{lemma}\blTitle{bl.orpmw} pull middle wff from disjunction wff-list
\begin{align}
\vdash ( \bigvee( \mwl{l_1} \mw{\varphi} \mwl{l_2} ) \leftrightarrow (
    \mw{\varphi} \vee \bigvee( \mwl{l_2} \mwl{l_1} ) )
    )\label{eq:bl.orpmw}\tag{bl.orpmw}
\end{align}
\end{lemma}

\begin{proof}
\begin{align}
  1 &&  & \vdash ( \bigvee( \mwl{l_1} \mw{\varphi} \mwl{l_2} ) \leftrightarrow \bigvee( \mw{\varphi} \mwl{l_2} \mwl{l_1} ) )&(\mbox{\ref{eq:bl.orcomphfirst}})\notag
  \\ 2 &&  & \vdash ( \bigvee( \mw{\varphi} \mwl{l_2} \mwl{l_1} ) \leftrightarrow ( \mw{\varphi} \vee \bigvee( \mwl{l_2} \mwl{l_1} ) ) )&(\mbox{\ref{eq:bl.orpfw}})\notag
  \\ 3 &&  & \vdash ( \bigvee( \mwl{l_1} \mw{\varphi} \mwl{l_2} ) \leftrightarrow ( \mw{\varphi} \vee \bigvee( \mwl{l_2} \mwl{l_1} ) ) )&(\mbox{\ref{eq:bitri}},1,2)\notag
\end{align}
\end{proof}

\begin{lemma}\blTitle{bl.orabpm} absorb parentheses, middle term in
disjunction wff-list
\begin{align}
\vdash ( \bigvee( \mwl{l_1} \bigvee( \mwl{l_2} ) \mwl{l_3} ) \leftrightarrow
    \bigvee( \mwl{l_1} \mwl{l_2} \mwl{l_3} )
    )\label{eq:bl.orabpm}\tag{bl.orabpm}
\end{align}
\end{lemma}

\begin{proof}
\begin{align}
  1 &&  & \vdash ( \bigvee( \mwl{l_1} \bigvee( \mwl{l_2} ) \mwl{l_3} ) \leftrightarrow \bigvee( \bigvee( \mwl{l_2} ) \mwl{l_3} \mwl{l_1} ) )&(\mbox{\ref{eq:bl.orcomwlwl}})\notag
  \\ 2 &&  & \vdash ( \bigvee( \bigvee( \mwl{l_2} ) \mwl{l_3} \mwl{l_1} ) \leftrightarrow \bigvee( \mwl{l_2} \mwl{l_3} \mwl{l_1} ) )&(\mbox{\ref{eq:bl.orabpf}})\notag
  \\ 3 &&  & \vdash ( \bigvee( \mwl{l_2} \mwl{l_3} \mwl{l_1} ) \leftrightarrow \bigvee( \mwl{l_1} \mwl{l_2} \mwl{l_3} ) )&(\mbox{\ref{eq:bl.orcomwlwl}})\notag
  \\ 4 &&  & \vdash ( \bigvee( \mwl{l_1} \bigvee( \mwl{l_2} ) \mwl{l_3} ) \leftrightarrow \bigvee( \mwl{l_1} \mwl{l_2} \mwl{l_3} ) )&(\mbox{\ref{eq:3bitri}},1,2,3)\notag
\end{align}
\end{proof}

\begin{lemma}\blTitle{bl.ais} And-Introduction Schema of wff
\begin{align}
\vdash ( \bigwedge( \mwl{l_1} \mw{\varphi} \mwl{l_2} ) \rightarrow \mw{\varphi}
    )\label{eq:bl.ais}\tag{bl.ais}
\end{align}
\end{lemma}

\begin{proof}
\begin{align}
  1 &&  & \vdash ( \bigwedge( \mwl{l_1} \mw{\varphi} \mwl{l_2} ) \leftrightarrow ( \mw{\varphi} \wedge \bigwedge( \mwl{l_2} \mwl{l_1} ) ) )&(\mbox{\ref{eq:bl.anpmw}})\notag
  \\ 2 &&  & \vdash ( ( \mw{\varphi} \wedge \bigwedge( \mwl{l_2} \mwl{l_1} ) ) \rightarrow \mw{\varphi} )&(\mbox{\ref{eq:simpl}})\notag
  \\ 3 &&  & \vdash ( \bigwedge( \mwl{l_1} \mw{\varphi} \mwl{l_2} ) \rightarrow \mw{\varphi} )&(\mbox{\ref{eq:sylbi}},1,2)\notag
\end{align}
\end{proof}

\begin{lemma}\blTitle{bl.aiswl} And-Introduction Schema of wff-List
\begin{align}
\vdash ( \bigwedge( \mwl{l_1} \mwl{l_2} \mwl{l_3} ) \rightarrow \bigwedge(
    \mwl{l_2} ) )\label{eq:bl.aiswl}\tag{bl.aiswl}
\end{align}
\end{lemma}

\begin{proof}
\begin{align}
  1 &&  & \vdash ( \bigwedge( \mwl{l_1} \bigwedge( \mwl{l_2} ) \mwl{l_3} ) \leftrightarrow \bigwedge( \mwl{l_1} \mwl{l_2} \mwl{l_3} ) )&(\mbox{\ref{eq:bl.anabpm}})\notag
  \\ 2 &&  & \vdash ( \bigwedge( \mwl{l_1} \bigwedge( \mwl{l_2} ) \mwl{l_3} ) \leftrightarrow ( \bigwedge( \mwl{l_2} ) \wedge \bigwedge( \mwl{l_3} \mwl{l_1} ) ) )&(\mbox{\ref{eq:bl.anpmw}})\notag
  \\ 3 &&  & \vdash ( \bigwedge( \mwl{l_1} \mwl{l_2} \mwl{l_3} ) \leftrightarrow ( \bigwedge( \mwl{l_2} ) \wedge \bigwedge( \mwl{l_3} \mwl{l_1} ) ) )&(\mbox{\ref{eq:bitr3i}},1,2)\notag
  \\ 4 &&  & \vdash ( ( \bigwedge( \mwl{l_2} ) \wedge \bigwedge( \mwl{l_3} \mwl{l_1} ) ) \rightarrow \bigwedge( \mwl{l_2} ) )&(\mbox{\ref{eq:simpl}})\notag
  \\ 5 &&  & \vdash ( \bigwedge( \mwl{l_1} \mwl{l_2} \mwl{l_3} ) \rightarrow \bigwedge( \mwl{l_2} ) )&(\mbox{\ref{eq:sylbi}},3,4)\notag
\end{align}
\end{proof}

